% Las tecnologías digitales y la comunicación — Espagnol 1ère A — Cours signé OnBuch+
% Programme officiel MINESEC. Compiler avec : tectonic EA18-tecnologias-digitales-comunicacion.tex
\newcommand{\DOCMATIERE}{Espagnol}
\newcommand{\DOCNIVEAU}{1ère A}
\newcommand{\DOCMODULE}{Módulo 5 : Les mass médias et la communication}
\newcommand{\DOCLECON}{EA18}
\newcommand{\DOCTITRE}{Las tecnologías digitales y la comunicación}
% OnBuch+ — préambule « cours signés » ENRICHI (compilé avec tectonic / XeLaTeX)
% Système visuel « Pop » d'OnBuch+ : fond crème (jamais blanc), bordures encre
% épaisses, ombres « dures » décalées, titres Archivo Black en capitales.
% Version MATHÉMATIQUES : graphes (pgfplots/tikz), boîtes pédagogiques
% (définition, propriété, méthode, attention, le savais-tu, exercice résolu,
% exercice, TP, bilan), page de garde et signature OnBuch+ ancrée en bas.
%
% Macros attendues dans main.tex AVANT \input{preamble} :
%   \DOCMATIERE  \DOCNIVEAU  \DOCMODULE  \DOCLECON (numéro)  \DOCTITRE
\documentclass[11pt]{article}

\usepackage{fontspec}
\usepackage[spanish,french]{babel} % français = langue principale ; l'espagnol via \esp{..} / espagnol
\usepackage[a4paper,margin=2cm,top=3.4cm,bottom=2.8cm,headheight=1.4cm,footskip=1.3cm]{geometry}
\usepackage{amsmath,amssymb,amsfonts}
\usepackage{mathtools} % \xmapsto, \coloneqq... souvent utilisés par les modèles
\usepackage{cancel} % simplifications barrées (bilans de forces)
% \abs{z} (module, valeur absolue) et \norm{u} (norme), utilisés par les modèles
\providecommand{\abs}{}\let\abs\relax\DeclarePairedDelimiter{\abs}{\lvert}{\rvert}
\providecommand{\norm}{}\let\norm\relax\DeclarePairedDelimiter{\norm}{\lVert}{\rVert}
\usepackage[table]{xcolor} % [table] : fournit aussi \rowcolors (alternance de lignes)
\usepackage{pagecolor}
\usepackage{tikz}
\usetikzlibrary{arrows.meta,positioning,calc,shapes.geometric,shapes.misc,shapes.arrows,decorations.pathmorphing,decorations.pathreplacing,decorations.markings,patterns,shadows,mindmap,trees,fit,backgrounds,matrix,intersections,angles,quotes}
\usepackage{pgfplots}
\usepackage[version=4]{mhchem} % équations nucléaires \ce{^{235}_{92}U}
\usepackage[european,siunitx]{circuitikz} % schémas électriques
\pgfplotsset{compat=1.18}
\usepgfplotslibrary{fillbetween,groupplots} % aires entre courbes (calcul intégral)
\usepackage{enumitem}
\usepackage{fancyhdr}
% [most] charge toutes les bibliothèques tcolorbox (listings, minted, raster,
% documentation...) et gonfle inutilement la mémoire TeX sur un document long
% (~300 boîtes) : on ne charge que ce qui est réellement utilisé.
\usepackage[skins,breakable]{tcolorbox}
\usepackage{graphicx}
\usepackage{colortbl}
\usepackage{booktabs}
\usepackage{tabularx}
\usepackage{diagbox} % case d'en-tête diagonale des tableaux à double entrée
% Colonnes extensibles centrée / gauche / droite (C, L, R), souvent utilisées par les modèles
\newcolumntype{C}{>{\centering\arraybackslash}X}
\newcolumntype{L}{>{\raggedright\arraybackslash}X}
\newcolumntype{R}{>{\raggedleft\arraybackslash}X}
\usepackage{array}
\usepackage{multicol}
\usepackage{siunitx}
% parse-numbers=false : accepte aussi \SI{2,5 \times 10^{-3}}{...}
\sisetup{locale=FR,output-decimal-marker={,},group-minimum-digits=4,parse-numbers=false}
\DeclareSIUnit\ohm{\ensuremath{\symup{\Omega}}} % U+2126 absent de la police
\DeclareSIUnit\division{div} % réglages d'oscilloscope (ms/div, V/div)
\DeclareSIUnit\tr{tr} % tours (vitesse de rotation, ex. \SI{1500}{\tr\per\minute})
\DeclareSIUnit\molar{M} % concentration molaire (ex. \SI{3}{\molar}, \SI{1}{\micro\molar})
\DeclareSIUnit\an{an} % année (le modèle confond parfois avec \year, un compteur TeX) — ex. \SI{5}{\kilo\meter\per\an}
\DeclareSIUnit\FCFA{FCFA} % franc CFA (ex. \SI{500}{\FCFA\per\kilo\gram})
\DeclareSIUnit\degreeBrix{\textdegree Brix} % teneur en sucre d'un sirop/jus (réfractométrie)
\DeclareSIUnit\month{mois} % le modèle utilise parfois \month, un compteur TeX, comme unité de durée
\usepackage{qrcode}
% \degree : commande fréquemment utilisée par les modèles pour un angle ou une
% température en dehors de \SI{}{} (non fournie par les paquets chargés).
\providecommand{\degree}{^{\circ}}
% \qed (fin de démonstration), utilisé par les modèles sans amsthm
\providecommand{\qed}{\hfill$\square$}
% \barre : double barre des valeurs interdites dans un tableau de variations (tkz-tab)
\providecommand{\barre}{\ensuremath{\|}}
% environnement proof (amsthm non chargé) utilisé par les modèles
\ifdefined\proof\else\newenvironment{proof}[1][Démonstration]{\par\noindent\textit{#1.}\ }{\qed\par}\fi
\usepackage{lastpage}
\usepackage{hyperref}
\hypersetup{hidelinks}

% ============================================================
% Palette « Pop » OnBuch+ — mêmes hex que la classe P (plus_ui.dart)
% ============================================================
\definecolor{poporange}{HTML}{F59321}
\definecolor{popdark}{HTML}{E07A0C}
\definecolor{popcream}{HTML}{FFF6E8}
\definecolor{popink}{HTML}{1C1714}
\definecolor{poppurple}{HTML}{7A5AE0}
\definecolor{popgreen}{HTML}{2E9E6A}
\definecolor{poppink}{HTML}{E0457B}
\definecolor{popblue}{HTML}{2F7DE1}
\definecolor{popgold}{HTML}{C98A12}
\definecolor{popmuted}{HTML}{8A7F76}
\definecolor{popline}{HTML}{EADFD2}
\definecolor{popgray}{HTML}{8A7F76}
\colorlet{popgrey}{popgray}
% Alias tolérants (le modèle nomme parfois les couleurs différemment)
\colorlet{popwhite}{white}
\colorlet{popblack}{popink}
\colorlet{popred}{poppink}
\colorlet{popyellow}{popgold}
% Teintes claires (fonds de tableaux/encadrés) — le modèle nomme parfois
% les couleurs avec un suffixe L (ex. popblueL) sans les avoir définies.
\colorlet{poporangeL}{poporange!20}
\colorlet{popdarkL}{popdark!20}
\colorlet{poppurpleL}{poppurple!20}
\colorlet{popgreenL}{popgreen!20}
\colorlet{poppinkL}{poppink!20}
\colorlet{popblueL}{popblue!20}
\colorlet{popgoldL}{popgold!20}
\colorlet{popredL}{popred!20}
\colorlet{popgrayL}{popgray!20}
% Teintes claires pour fonds de tableaux / zones
\colorlet{popblueL}{popblue!10!white}
\colorlet{poporangeL}{poporange!14!white}
\colorlet{popgreenL}{popgreen!12!white}
\colorlet{poppurpleL}{poppurple!12!white}
\colorlet{poppinkL}{poppink!12!white}
\colorlet{popgoldL}{popgold!14!white}

\pagecolor{popcream}

% ============================================================
% Typographie
% ============================================================
\defaultfontfeatures{Ligatures=TeX}
\setmainfont{PlusJakartaSans-Medium.ttf}[Path=../../../fonts/,BoldFont=PlusJakartaSans-Bold.ttf]
% Maths en sans-serif (Fira Math) avec chiffres et lettres droites de Plus
% Jakarta, pour que les formules se fondent dans le texte.
\usepackage[math-style=ISO,bold-style=ISO]{unicode-math}
\setmathfont{FiraMath-Regular.otf}[Path=../../../fonts/]
\setmathfont{PlusJakartaSans-Medium.ttf}[Path=../../../fonts/,range={up/num,up/{Latin,latin},\mathup}]
\usepackage{icomma} % virgule décimale sans espace en mode math
% Caractères mathématiques Unicode absents des polices de texte (Plus Jakarta,
% Archivo Black) : rendus par leur équivalent mathématique, en texte comme en
% formule — sinon ils s'affichent en carré vide (ex. « Arithmétique dans ℤ »).
\usepackage{newunicodechar}
\newunicodechar{ℝ}{\ensuremath{\mathbb{R}}}
\newunicodechar{ℤ}{\ensuremath{\mathbb{Z}}}
\newunicodechar{ℂ}{\ensuremath{\mathbb{C}}}
\newunicodechar{ℕ}{\ensuremath{\mathbb{N}}}
\newunicodechar{ℚ}{\ensuremath{\mathbb{Q}}}
\newunicodechar{𝔻}{\ensuremath{\mathbb{D}}}
\newunicodechar{α}{\ensuremath{\alpha}}
\newunicodechar{β}{\ensuremath{\beta}}
\newunicodechar{γ}{\ensuremath{\gamma}}
\newunicodechar{δ}{\ensuremath{\delta}}
\newunicodechar{ε}{\ensuremath{\varepsilon}}
\newunicodechar{θ}{\ensuremath{\theta}}
\newunicodechar{λ}{\ensuremath{\lambda}}
\newunicodechar{μ}{\ensuremath{\mu}}
\newunicodechar{σ}{\ensuremath{\sigma}}
\newunicodechar{τ}{\ensuremath{\tau}}
\newunicodechar{φ}{\ensuremath{\varphi}}
\newunicodechar{ψ}{\ensuremath{\psi}}
\newunicodechar{ω}{\ensuremath{\omega}}
\newunicodechar{Σ}{\ensuremath{\Sigma}}
% majuscules produites par \MakeUppercase dans les titres (α-aminés -> Α)
\newunicodechar{Α}{\ensuremath{\alpha}}
\newunicodechar{Β}{\ensuremath{\beta}}
\newunicodechar{Γ}{\ensuremath{\Gamma}}
\newunicodechar{Δ}{\ensuremath{\Delta}}
\newunicodechar{Ε}{\ensuremath{\varepsilon}}
\newunicodechar{Θ}{\ensuremath{\Theta}}
\newunicodechar{Λ}{\ensuremath{\Lambda}}
\newunicodechar{Μ}{\ensuremath{\mu}}
\newunicodechar{Τ}{\ensuremath{\tau}}
\newunicodechar{Φ}{\ensuremath{\Phi}}
\newunicodechar{Ψ}{\ensuremath{\Psi}}
\newunicodechar{Ω}{\ensuremath{\Omega}}
\newunicodechar{↦}{\ensuremath{\mapsto}}
\newunicodechar{∈}{\ensuremath{\in}}
\newunicodechar{∉}{\ensuremath{\notin}}
\newunicodechar{∧}{\ensuremath{\wedge}}
\newunicodechar{⇔}{\ensuremath{\Leftrightarrow}}
\newunicodechar{⟺}{\ensuremath{\Longleftrightarrow}}
\newunicodechar{⇒}{\ensuremath{\Rightarrow}}
\newunicodechar{⟹}{\ensuremath{\Longrightarrow}}
\newunicodechar{∘}{\ensuremath{\circ}}
\newunicodechar{∪}{\ensuremath{\cup}}
\newunicodechar{∩}{\ensuremath{\cap}}
\newunicodechar{≡}{\ensuremath{\equiv}}
\newunicodechar{⊂}{\ensuremath{\subset}}
\newunicodechar{⊥}{\ensuremath{\perp}}
\newunicodechar{∃}{\ensuremath{\exists}}
\newunicodechar{∀}{\ensuremath{\forall}}
\newunicodechar{∛}{\ensuremath{\sqrt[3]{\,}}}
\newunicodechar{∜}{\ensuremath{\sqrt[4]{\,}}}
\newunicodechar{⃗}{\ensuremath{{}^{\to}}}
\newunicodechar{ⁿ}{\ifmmode^{n}\else\textsuperscript{\ensuremath{n}}\fi}
\newunicodechar{ˣ}{\ifmmode^{x}\else\textsuperscript{\ensuremath{x}}\fi}
\newunicodechar{ᵇ}{\ifmmode^{b}\else\textsuperscript{\ensuremath{b}}\fi}
\newunicodechar{ʳ}{\ifmmode^{r}\else\textsuperscript{\ensuremath{r}}\fi}
\newunicodechar{ᵘ}{\ifmmode^{u}\else\textsuperscript{\ensuremath{u}}\fi}
\newunicodechar{ʸ}{\ifmmode^{y}\else\textsuperscript{\ensuremath{y}}\fi}
\newunicodechar{ᵃ}{\ifmmode^{a}\else\textsuperscript{\ensuremath{a}}\fi}
\newunicodechar{ᵉ}{\ifmmode^{e}\else\textsuperscript{\ensuremath{e}}\fi}
\newunicodechar{ᵏ}{\ifmmode^{k}\else\textsuperscript{\ensuremath{k}}\fi}
\newunicodechar{ᵖ}{\ifmmode^{p}\else\textsuperscript{\ensuremath{p}}\fi}
\newunicodechar{ᴺ}{\ifmmode^{N}\else\textsuperscript{\ensuremath{N}}\fi}
\newunicodechar{ᶜ}{\ifmmode^{c}\else\textsuperscript{\ensuremath{c}}\fi}
\newunicodechar{ʰ}{\ifmmode^{h}\else\textsuperscript{\ensuremath{h}}\fi}
\newunicodechar{ᵠ}{\ifmmode^{q}\else\textsuperscript{\ensuremath{q}}\fi}
\newunicodechar{ₙ}{\ifmmode_{n}\else\textsubscript{\ensuremath{n}}\fi}
\newunicodechar{ₐ}{\ifmmode_{a}\else\textsubscript{\ensuremath{a}}\fi}
\newunicodechar{ᵢ}{\ifmmode_{i}\else\textsubscript{\ensuremath{i}}\fi}
\newunicodechar{ᵣ}{\ifmmode_{r}\else\textsubscript{\ensuremath{r}}\fi}
\newunicodechar{ₖ}{\ifmmode_{k}\else\textsubscript{\ensuremath{k}}\fi}
\newunicodechar{ᵦ}{\ifmmode_{\beta}\else\textsubscript{\ensuremath{\beta}}\fi}
\newunicodechar{ₓ}{\ifmmode_{x}\else\textsubscript{\ensuremath{x}}\fi}
\newfontfamily\popdisplay{ArchivoBlack-Regular.ttf}[Path=../../../fonts/]
\newfontfamily\poplogo{SpaceGrotesk-Bold.ttf}[Path=../../../fonts/]
\newfontfamily\popsemi{PlusJakartaSans-SemiBold.ttf}[Path=../../../fonts/]
\newfontfamily\popxbold{PlusJakartaSans-ExtraBold.ttf}[Path=../../../fonts/]
\newfontfamily\popxboldls{PlusJakartaSans-ExtraBold.ttf}[Path=../../../fonts/,LetterSpace=3.0]

\setlist[enumerate]{leftmargin=1.7em,itemsep=5pt,topsep=4pt}
\setlist[itemize]{leftmargin=1.7em,itemsep=4pt,topsep=4pt,label={\color{poporange}\textbullet}}
\setlength{\parindent}{0pt}
\setlength{\parskip}{5pt}
\renewcommand{\arraystretch}{1.25}

% pgfplots : style Pop par défaut
\pgfplotsset{
  every axis/.append style={axis line style={popink,line width=1pt},tick style={popink},
    grid style={popline,line width=0.5pt},label style={font=\small},tick label style={font=\footnotesize}},
  popcurve/.style={poporange,line width=1.8pt,no markers,smooth},
}
% Même style disponible dans un \draw TikZ simple (hors pgfplots)
\tikzset{popcurve/.style={poporange,line width=1.8pt}}
\tikzset{popgrid/.style={popline,very thin}}
\colorlet{popcurve}{poporange} % le nom du style est parfois utilisé comme couleur

% ============================================================
% Pill / squiggle
% ============================================================
\newcommand{\poppill}[3]{%
  \tikz[baseline=-0.55ex]\node[draw=popink,line width=1.2pt,rounded corners=2.6pt,%
    fill=#2,inner xsep=6.5pt,inner ysep=3pt]{%
    \popxboldls\fontsize{7}{7}\selectfont\color{#3}\MakeUppercase{#1}};%
}
\newcommand{\popsquiggle}[1]{%
  \tikz[baseline=-0.4ex]{
    \draw[#1,line width=2.6pt,line cap=round]
      plot [smooth] coordinates {(0,0.09) (0.34,0.01) (0.68,0.21) (1.02,0.01) (1.36,0.10)};
  }%
}
% Mot-clé surligné (marqueur orange sous le texte)
\newcommand{\cle}[1]{{\popxbold\color{popdark}#1}}

% ============================================================
% En-tête / pied
% ============================================================
\pagestyle{fancy}
\fancyhf{}
\renewcommand{\headrulewidth}{0pt}
\renewcommand{\footrulewidth}{0pt}
\lhead{\raisebox{-0.15ex}{\poplogo\fontsize{13}{13}\selectfont\color{popink}OnBuch\color{poporange}+}}
\rhead{\poppill{\DOCMATIERE\ \textbullet\ \DOCNIVEAU\ \textbullet\ Leçon \DOCLECON}{white}{popink}}
\lfoot{\popxbold\fontsize{7.6}{7.6}\selectfont\color{popmuted}\MakeUppercase{Cours signé OnBuch+}\ \textbullet\ {\popsemi\DOCTITRE}}
\rfoot{\tikz[baseline=-0.6ex]\node[rounded corners=8.5pt,fill=popink,minimum height=17pt,minimum width=17pt,inner xsep=5pt,inner sep=0]{\popxbold\fontsize{8}{8}\selectfont\color{popcream}\thepage\,/\,\pageref*{LastPage}};}

% ============================================================
% Titres
% ============================================================
\newcommand{\chaptitre}[2]{%
  \begin{tcolorbox}[enhanced,colback=poporange,colframe=popink,boxrule=2.6pt,arc=11pt,
      drop shadow=popink,left=15pt,right=15pt,top=12pt,bottom=11pt,before skip=3pt,after skip=18pt]
    \poppill{#2}{popink}{popcream}\par\vspace{9pt}
    {\raggedright\hyphenpenalty=10000\exhyphenpenalty=10000\popdisplay\fontsize{22}{24}\selectfont\color{popcream}\MakeUppercase{#1}\par}
  \end{tcolorbox}
}
% Section numérotée : pastille orange avec le numéro + titre Archivo
\newcounter{popsec}
\newcommand{\coursec}[1]{%
  \stepcounter{popsec}\setcounter{popsub}{0}%
  \par\vspace{16pt}\noindent
  \begin{minipage}{\linewidth}
  \tikz[baseline=-0.7ex]\node[draw=popink,line width=1.6pt,fill=poporange,rounded corners=4pt,minimum size=22pt,inner sep=0]{\popdisplay\fontsize{12}{12}\selectfont\color{popcream}\thepopsec};%
  \hspace{8pt}{\popdisplay\fontsize{15}{17}\selectfont\color{popink}\MakeUppercase{#1}}\par
  \vspace{4pt}\noindent{\color{poporange}\rule{36pt}{3.6pt}}
  \end{minipage}\par\nopagebreak\vspace{8pt}
}
\newcounter{popsub}
\newcommand{\courssub}[1]{%
  \stepcounter{popsub}%
  \par\vspace{9pt}\noindent{\popxbold\fontsize{11.5}{13}\selectfont\color{popink}{\color{poporange}\thepopsec.\thepopsub}\hspace{6pt}#1}\par\nopagebreak\vspace{4pt}
}

% ============================================================
% Boîtes pédagogiques — même recette : fond blanc, bordure épaisse colorée,
% ombre dure encre, badge plein. Titre optionnel [#1].
% ============================================================
\tcbset{popbase/.style={breakable,enhanced,colback=white,boxrule=2.3pt,arc=9pt,
  shadow={3pt}{-3pt}{0pt}{popink},left=14pt,right=14pt,top=11pt,bottom=11pt,before skip=11pt,after skip=15pt,
  fontupper=\popsemi\fontsize{10.6}{14.5}\selectfont\color{popink},
  fontlower=\popsemi\fontsize{10.6}{14.5}\selectfont\color{popink}}}
% Coche dessinée (le glyphe ✓ n'existe pas dans les polices du document)
\newcommand{\popcheck}[1]{\tikz[baseline=-0.5ex]\draw[#1,line width=1.6pt,line cap=round,line join=round](-0.13,0.0)--(-0.04,-0.09)--(0.13,0.1);}
\providecommand{\coche}{\popcheck{popgreen}}
\providecommand{\resultat}[1]{\par\smallskip\noindent\fcolorbox{popgreen}{popgreenL}{\parbox{\dimexpr\linewidth-2\fboxsep-2\fboxrule\relax}{\textbf{Résultat.} #1}}\par\smallskip}
\newcommand{\popboxhead}[3]{\poppill{#1}{#2}{white}\ifx\relax#3\relax\else\hspace{6pt}{\popxbold\fontsize{10.6}{12}\selectfont\color{popink}#3}\fi\par\vspace{7pt}}

\newtcolorbox{aretenir}[1][]{popbase,colframe=popgreen,before upper={\popboxhead{À retenir}{popgreen}{#1}}}
% Texte en espagnol (document de lecture, dialogue, poème, extrait) : cadre
% d'étude. Un titre optionnel (source, auteur) s'affiche dans l'en-tête.
\newtcolorbox{textebox}[1][]{popbase,colframe=popblue,before upper={\popboxhead{Texto}{popblue}{#1}\begin{otherlanguage}{spanish}},after upper={\end{otherlanguage}}}
% Marques ✓ / ✗ (forme correcte / fautive) souvent utilisées par les modèles
\providecommand{\cross}{\textcolor{poppink}{\ensuremath{\times}}}
\providecommand{\xmark}{\cross}\providecommand{\crossmark}{\cross}\providecommand{\cmark}{\ensuremath{\checkmark}}
% Ligne de réponse / justification à compléter (inventée par les modèles)
\providecommand{\justification}[1]{\par\noindent\textit{Justification~:} \dotfill\par}
% \textipa (package tipa non chargé) : la police ne couvre pas l'API -> simple passage
\providecommand{\textipa}[1]{#1}
% Soulignement (ulem non chargé) : \uline, \ul
\providecommand{\uline}[1]{\underline{#1}}\providecommand{\ul}[1]{\underline{#1}}
% Ordinaux français écrits en macro par les modèles : 1\ière, 3\ième
\newcommand{\ière}{\textsuperscript{re}}\newcommand{\ième}{\textsuperscript{e}}
% \exemple (inventée par les modèles) : étiquette « Exemple : »
\providecommand{\exemple}{\textbf{Exemple~:}\ }
% Mot ou phrase en espagnol à l'intérieur d'un paragraphe français
\newcommand{\esp}[1]{\ifmmode\text{\textit{#1}}\else\begingroup\selectlanguage{spanish}\itshape #1\endgroup\fi}
\newtcolorbox{exemplebox}[1][]{popbase,colframe=poporange,before upper={\popboxhead{Exemple}{poporange}{#1}}}
\newtcolorbox{definition}[1][]{popbase,colframe=popblue,before upper={\popboxhead{Définition}{popblue}{#1}}}
\newtcolorbox{propriete}[1][]{popbase,colframe=poppurple,before upper={\popboxhead{Propriété}{poppurple}{#1}}}
\newtcolorbox{methode}[1][]{popbase,colframe=popgold,before upper={\popboxhead{Méthode}{popgold}{#1}}}
\newtcolorbox{attention}[1][]{popbase,colframe=poppink,before upper={\popboxhead{Attention}{poppink}{#1}}}
\newtcolorbox{experience}[1][]{popbase,colframe=popblue,colback=popblueL,before upper={\popboxhead{Expérience}{popblue}{#1}}}
% « Le savais-tu ? » : carte violette pleine (contexte, culture, Cameroun)
\newtcolorbox{savaistu}[1][]{popbase,colback=poppurple,colframe=popink,
  fontupper=\popsemi\fontsize{10.4}{14.2}\selectfont\color{white},
  before upper={\poppill{Le savais-tu\,?}{white}{poppurple}\ifx\relax#1\relax\else\hspace{6pt}{\popxbold\color{white}#1}\fi\par\vspace{7pt}}}
% Exercice résolu : énoncé en haut, \tcblower puis la solution sur fond crème
\newcounter{popexo}\newcounter{popexr}
\newtcolorbox{exoresolu}[1][]{popbase,colframe=poporange,colbacklower=popcream,
  segmentation style={popink,line width=1.2pt,dashed},
  before upper={\stepcounter{popexr}\popboxhead{Exercice résolu \thepopexr}{poporange}{#1}},
  before lower={\poppill{Solution}{popink}{popcream}\par\vspace{6pt}}}
% Exercice (non résolu en place) — la correction est dans la partie Corrigés
\newtcolorbox{exercice}[1][]{popbase,colframe=popink,
  before upper={\stepcounter{popexo}\popboxhead{Exercice \thepopexo}{popink}{#1}}}
% Corrigé
\newtcolorbox{corrige}[1][]{popbase,colframe=popgreen,colback=popgreenL,
  before upper={\popboxhead{Corrigé}{popgreen}{#1}}}
% Objectifs / prérequis / bilan
\newtcolorbox{objectifs}{popbase,colframe=popink,colback=poporangeL,
  before upper={\popboxhead{Objectifs de la leçon}{popink}{}}}
\newtcolorbox{prerequis}{popbase,colframe=popink,colback=popblueL,
  before upper={\popboxhead{Prérequis}{popblue}{}}}
\newtcolorbox{bilan}{popbase,colframe=popink,colback=white,boxrule=2.8pt,
  before upper={\popboxhead{Fiche bilan}{poporange}{L'essentiel à connaître pour l'examen}}}
% Figure encadrée avec légende
\newtcolorbox{popfigure}[1][]{enhanced,breakable=false,colback=white,colframe=popline,boxrule=1.4pt,arc=8pt,
  left=10pt,right=10pt,top=10pt,bottom=8pt,before skip=10pt,after skip=12pt,halign=center,
  lower separated=false,halign lower=center,
  fontlower=\popsemi\fontsize{9}{11.5}\selectfont\color{popmuted},
  #1}
\newcommand{\legende}[1]{\par\vspace{6pt}{\popsemi\fontsize{9}{11.5}\selectfont\color{popmuted}#1}}

% ============================================================
% Signature OnBuch+
% ============================================================
\newcommand{\poplogoinline}[1]{{\poplogo\fontsize{#1}{#1}\selectfont\color{popink}OnBuch\color{poporange}+}}
\newcommand{\signatureonbuch}{%
  \begin{tcolorbox}[enhanced,colback=popink,colframe=popink,boxrule=2.2pt,arc=10pt,
      drop shadow=poporange,left=14pt,right=14pt,top=10pt,bottom=10pt,before skip=0pt,after skip=0pt]
    \tikz[baseline=-0.7ex]\node[circle,fill=popgreen,draw=popcream,line width=1.3pt,minimum size=22pt,inner sep=0]{\popcheck{white}};%
    \hspace{9pt}\begin{minipage}[c]{0.62\linewidth}
      {\popxbold\fontsize{12}{13}\selectfont\color{popcream}Signé\ \textbullet\ {\poplogo OnBuch\color{poporange}+}}\par\vspace{2pt}
      {\raggedright\popsemi\fontsize{8}{10.5}\selectfont\color{popline}\MakeUppercase{Équipe pédagogique OnBuch+}\\\MakeUppercase{Conforme au programme officiel MINESEC}\par}
    \end{minipage}\hfill
    {\poplogo\fontsize{22}{22}\selectfont\color{popcream}OnBuch\color{poporange}+}
  \end{tcolorbox}%
}

% ============================================================
% Page de garde
% #1 titre · #2 matière · #3 niveau · #4 module · #5 n° leçon · #6 durée ·
% #7 description / objectifs (texte ou itemize)
% ============================================================
\newcommand{\pagegarde}[7]{%
  \thispagestyle{empty}
  \newgeometry{a4paper,margin=1.7cm,top=1.6cm,bottom=1.4cm}%
  \begin{tikzpicture}[remember picture,overlay]
    \fill[poporange!18] ([xshift=-3.2cm,yshift=-3.6cm]current page.north east) circle (4.6cm);
    \fill[poppurple!14] ([xshift=2.4cm,yshift=7.6cm]current page.south west) circle (3.8cm);
  \end{tikzpicture}%
  {\poplogo\fontsize{30}{30}\selectfont\color{popink}OnBuch\color{poporange}+}\hfill
  \raisebox{0.6ex}{\poppill{Cours signé \textbullet\ Premium}{popink}{popcream}}\par
  \vspace{1pt}\popsquiggle{poppurple}\par
  \vspace{8pt}
  \begin{tcolorbox}[enhanced,colback=poporange,colframe=popink,boxrule=2.8pt,arc=13pt,
      drop shadow=popink,left=18pt,right=18pt,top=12pt,bottom=13pt,after skip=10pt]
    \poppill{#2 \textbullet\ #3}{popink}{popcream}\hspace{5pt}\poppill{#4}{white}{popink}\par\vspace{8pt}
    {\popdisplay\fontsize{14}{14}\selectfont\color{popink}LEÇON #5}\par\vspace{4pt}
    {\raggedright\hyphenpenalty=10000\exhyphenpenalty=10000\popdisplay\fontsize{25}{27}\selectfont\color{popcream}\MakeUppercase{#1}\par}
    \vspace{8pt}
    {\popxbold\fontsize{9.5}{9.5}\selectfont\color{popink}\MakeUppercase{Durée conseillée : #6}}
  \end{tcolorbox}
  \begin{tcolorbox}[enhanced,colback=white,colframe=popink,boxrule=2.2pt,arc=10pt,
      drop shadow=popink,left=16pt,right=16pt,top=10pt,bottom=10pt,after skip=8pt]
    \poppill{Dans cette leçon}{poppurple}{white}\par\vspace{5pt}
    \popsemi\fontsize{9.3}{12.6}\selectfont\color{popink}#7
  \end{tcolorbox}
  \noindent\begin{minipage}[t]{0.487\linewidth}
    \begin{tcolorbox}[enhanced,colback=popgreen,colframe=popink,boxrule=2.2pt,arc=10pt,
        drop shadow=popink,left=11pt,right=11pt,top=10pt,bottom=10pt,height=3.1cm,valign=center]
      \tcbox[colback=white,colframe=popink,boxrule=1.3pt,arc=5pt,boxsep=0pt,nobeforeafter,
        left=3pt,right=3pt,top=3pt,bottom=3pt,valign=center]{\qrcode[height=1.9cm]{https://play.google.com/store/apps/details?id=com.luvvix.onbuch&hl=fr}}%
      \hspace{8pt}\begin{minipage}[c]{0.5\linewidth}
        {\popdisplay\fontsize{11}{12.5}\selectfont\color{white}\MakeUppercase{Télécharge OnBuch}}\par\vspace{4pt}
        {\raggedright\popsemi\fontsize{8.4}{11}\selectfont\color{white}Gratuit \textbullet\ Google Play\\Tuteur IA, annales, concours, résultats\par}
      \end{minipage}
    \end{tcolorbox}
  \end{minipage}\hfill
  \begin{minipage}[t]{0.487\linewidth}
    \begin{tcolorbox}[enhanced,colback=poppurple,colframe=popink,boxrule=2.2pt,arc=10pt,
        drop shadow=popink,left=11pt,right=11pt,top=10pt,bottom=10pt,height=3.1cm,valign=center]
      \tikz[baseline=-0.7ex]\node[circle,fill=white,draw=popink,line width=1.3pt,minimum size=1.6cm,inner sep=0]{\popdisplay\fontsize{24}{24}\selectfont\color{poppurple}+};%
      \hspace{8pt}\begin{minipage}[c]{0.6\linewidth}
        {\popdisplay\fontsize{11}{12.5}\selectfont\color{white}\MakeUppercase{Passe à OnBuch+}}\par\vspace{4pt}
        {\raggedright\popsemi\fontsize{8.4}{11}\selectfont\color{white}Tous les cours signés, Tuteur IA illimité, fiches PDF — onglet OnBuch+ de l'app\par}
      \end{minipage}
    \end{tcolorbox}
  \end{minipage}
  \vspace{8pt}
  \popcontenu
  \vfill
  \signatureonbuch
  \restoregeometry
  \newpage
}

% Rangée « Ce cours contient » (page de garde)
\newcommand{\poptile}[3]{% #1 couleur #2 gros texte #3 légende
  \begin{tcolorbox}[enhanced,colback=#1,colframe=popink,boxrule=2pt,arc=9pt,shadow={2.5pt}{-2.5pt}{0pt}{popink},
      left=6pt,right=6pt,top=9pt,bottom=9pt,halign=center,valign=center,height=1.75cm,width=\linewidth,nobeforeafter]
    {\popdisplay\fontsize{12.5}{13}\selectfont\color{white}\MakeUppercase{#2}}\par\vspace{3pt}
    {\centering\popsemi\fontsize{7.8}{9.5}\selectfont\color{white}#3\par}
  \end{tcolorbox}}
\newcommand{\popcontenu}{%
  \par\noindent{\popxboldls\fontsize{7.5}{7.5}\selectfont\color{popmuted}CE COURS SIGNÉ CONTIENT}\par\vspace{7pt}
  \noindent\begin{minipage}[t]{0.235\linewidth}\poptile{popblue}{Cours}{complet, illustré, pas à pas}\end{minipage}\hfill
  \begin{minipage}[t]{0.235\linewidth}\poptile{poporange}{Méthodes}{et exercices résolus}\end{minipage}\hfill
  \begin{minipage}[t]{0.235\linewidth}\poptile{poppink}{Exercices}{type examen + corrigés}\end{minipage}\hfill
  \begin{minipage}[t]{0.235\linewidth}\poptile{popgreen}{Fiche bilan}{l'essentiel à retenir}\end{minipage}\par}

% Sommaire
\newtcolorbox{sommaire}{enhanced,colback=white,colframe=popink,boxrule=2.2pt,arc=10pt,
  drop shadow=popink,left=18pt,right=18pt,top=13pt,bottom=13pt,before skip=6pt,after skip=20pt,
  before upper={\poppill{Sommaire}{poppurple}{white}\par\vspace{9pt}\popsemi\fontsize{10.6}{14.5}\selectfont\color{popink}}}

% Signature de fin de cours — poussée tout en bas de la dernière page
\newcommand{\findecours}{%
  \par\vfill\nopagebreak
  \begin{center}\popsquiggle{poporange}\end{center}\vspace{4pt}
  \signatureonbuch
}
\AtBeginDocument{\def\llbracket{\mathopen{[\mkern-3.5mu[}}\def\rrbracket{\mathclose{]\mkern-3.5mu]}}} % intervalles d'entiers (glyphe ⟦ absent de la police)
% Glyphes absents de Fira Math : redéfinis à partir de glyphes disponibles
\AtBeginDocument{%
  \def\setminus{\mathbin{\backslash}}\let\smallsetminus\setminus
  \def\vdots{\mathord{\vbox{\baselineskip3.2pt\lineskiplimit0pt\kern4pt\hbox{.}\hbox{.}\hbox{.}}}}%
  \def\ddots{\mathinner{\mkern1mu\raise6.5pt\hbox{.}\mkern2mu\raise3.6pt\hbox{.}\mkern2mu\raise0.7pt\hbox{.}\mkern1mu}}%
  \def\triangle{\mathord{\scalebox{0.9}{$\symup{\Delta}$}}}%
  \def\nabla{\mathord{\raisebox{\depth}{\scalebox{1}[-1]{$\symup{\Delta}$}}}}%
  \def\checkmark{\popcheck{popdark}}%
  \def\star{\mathbin{\ast}}\def\bigstar{\mathord{\ast}}%
  \def\diamond{\mathbin{\scalebox{0.6}{$\Diamond$}}}%
  \def\bigcup{\mathop{\mathchoice{\raisebox{-0.3em}{\scalebox{1.7}{$\cup$}}}{\scalebox{1.25}{$\cup$}}{\cup}{\cup}}\displaylimits}%
  \def\bigcap{\mathop{\mathchoice{\raisebox{-0.3em}{\scalebox{1.7}{$\cap$}}}{\scalebox{1.25}{$\cap$}}{\cap}{\cap}}\displaylimits}%
  \def\lesseqqgtr{\mathrel{\lessgtr}}\def\bigoplus{\mathop{\oplus}\displaylimits}%
}
\providecommand{\ssi}{si et seulement si} % abréviation fréquente
\AtBeginDocument{\providecommand{\wideparen}{\overparen}} % arc de cercle

%% FIGFIX-BEGIN v5 — correctifs globaux des figures (docs/onbuch-generation/tools/figfix)
\usepackage{adjustbox}\usepackage{etoolbox}\tcbuselibrary{raster}
% toute figure TikZ/pgfplots plus large que la ligne est réduite à la largeur disponible
\BeforeBeginEnvironment{tikzpicture}{\begin{adjustbox}{max width=\linewidth}}
\AfterEndEnvironment{tikzpicture}{\end{adjustbox}}
% légendes pgfplots lisibles par-dessus les courbes
\pgfplotsset{every axis/.append style={legend style={fill=white,fill opacity=0.92,text opacity=1,draw=black!15,inner sep=3pt}}}
% étiquettes d'arêtes (syntaxe edge["..."]) avec fond blanc
\tikzset{every edge quotes/.append style={fill=white,inner sep=1.2pt}}
%% FIGFIX-END
\begin{document}

\pagegarde{Las tecnologías digitales y la comunicación}{Espagnol}{1ère A}{Módulo 5 : Les mass médias et la communication}{EA18}{2 h}{Cette leçon mixte du module « Les mass médias et la communication » propose un texte support sur les technologies numériques, un lexique thématique riche et l'étude de l'impératif négatif et du subjonctif pour conseiller. Elle conduit l'élève à évaluer les risques du numérique (cyberharcèlement, addiction, vie privée) et à formuler des comportements responsables en espagnol.
\vspace{4pt}
\begin{multicols}{2}\small
\begin{itemize}
\item À la fin de cette leçon, je sais lire et comprendre un texte espagnol sur les technologies numériques
\item À la fin de cette leçon, je sais utiliser le lexique du numérique : internet, móvil, aplicación, contraseña, ciberacoso, adicción, privacidad, descargar
\item À la fin de cette leçon, je sais identifier et nommer les risques du numérique : ciberacoso, adicción, pérdida de privacidad
\item À la fin de cette leçon, je sais former l'impératif affirmatif et négatif pour donner des conseils : no compartas, protege tus datos
\item À la fin de cette leçon, je sais employer le subjonctif après « es importante que » pour recommander : es importante que protejas tu privacidad
\item À la fin de cette leçon, je sais comparer les structures de conseil en français et en espagnol et éviter les pièges de traduction
\end{itemize}\end{multicols}}

\begin{sommaire}
\begin{multicols}{2}
\begin{enumerate}
\item Texte support : « Conectados, ¿pero protegidos? »
\item Lexique thématique : el mundo digital
\item Grammaire 1 : l'impératif pour conseiller
\item Grammaire 2 : le subjonctif après « es importante que »
\item Méthode du commentaire et commentaire modèle
\item Production guidée : proposer des comportements responsables
\item Activité d'intégration
\item Méthodes et exercices résolus
\item Exercices
\item Corrigés des exercices
\item Fiche bilan
\end{enumerate}
\end{multicols}
\end{sommaire}

\begin{objectifs}
\begin{itemize}
\item À la fin de cette leçon, je sais lire et comprendre un texte espagnol sur les technologies numériques
\item À la fin de cette leçon, je sais utiliser le lexique du numérique : internet, móvil, aplicación, contraseña, ciberacoso, adicción, privacidad, descargar
\item À la fin de cette leçon, je sais identifier et nommer les risques du numérique : ciberacoso, adicción, pérdida de privacidad
\item À la fin de cette leçon, je sais former l'impératif affirmatif et négatif pour donner des conseils : no compartas, protege tus datos
\item À la fin de cette leçon, je sais employer le subjonctif après « es importante que » pour recommander : es importante que protejas tu privacidad
\item À la fin de cette leçon, je sais comparer les structures de conseil en français et en espagnol et éviter les pièges de traduction
\item À la fin de cette leçon, je sais commenter un texte sur le numérique selon la méthode du commentaire composé
\item À la fin de cette leçon, je sais rédiger un texte argumentatif proposant des comportements numériques responsables
\end{itemize}
\end{objectifs}

\begin{prerequis}
\begin{itemize}
\item Le présent de l'indicatif des verbes réguliers et irréguliers courants (Seconde)
\item Les pronoms personnels sujets et compléments (Seconde)
\item Le vocabulaire des médias et de la communication vu dans les leçons précédentes du module 5
\item Les structures de base de l'impératif affirmatif (début d'année de Première)
\item L'expression de l'opinion : pienso que, me parece que (Seconde)
\end{itemize}
\end{prerequis}

\newpage

\coursec{Texte support : « Conectados, ¿pero protegidos? »}

\courssub{Situation d'accroche et problématique}

À Yaoundé comme à Madrid, presque tous les lycéens possèdent un \cle{smartphone} et passent plusieurs heures par jour sur \esp{WhatsApp}, \esp{TikTok} ou \esp{Instagram}. Au lycée, dans les bendskins, au marché Mokolo : partout, des yeux rivés sur les écrans. Mais que se passe-t-il quand un camarade diffuse une photo humiliante dans un groupe de classe ? Quand on ne peut plus se déconnecter ? Quand une application utilise nos données personnelles sans notre accord ? Selon les études de l'UNESCO, environ un adolescent sur dix dans le monde déclare avoir déjà été victime de \esp{ciberacoso} (cyberharcèlement).

\textbf{Problématique :} comment parler de ces risques en espagnol et conseiller des comportements responsables ? C'est tout l'enjeu de cette leçon. Pour y répondre, nous allons d'abord lire l'histoire d'Andrés, un adolescent camerounais comme vous, victime de cyberharcèlement dans le groupe WhatsApp de sa classe.

\courssub{Lecture globale du texte}

\textbf{Déroulement de la lecture.} D'abord, une \cle{lecture silencieuse} individuelle (3 minutes) : soulignez les mots inconnus. Ensuite, l'enseignant lit le texte \cle{à voix haute} ; écoutez et suivez avec les yeux. Enfin, travail sur la prononciation de deux mots difficiles :

\begin{itemize}
\item \esp{ciberacoso} se prononce « si-bé-ra-ko-so » : le \esp{c} devant \esp{i} se prononce « s » en Amérique latine et « th » (comme dans l'anglais \emph{think}) dans la plupart des régions d'Espagne.
\item \esp{aplicación} se prononce « a-pli-ka-SION » : l'accent écrit sur le \esp{ó} indique que la dernière syllabe est accentuée.
\end{itemize}

\begin{textebox}[Conectados, ¿pero protegidos?]
Andrés tiene dieciséis años y vive en Yaundé. Como muchos jóvenes de su edad, pasa cinco horas al día delante de la pantalla de su móvil: descarga aplicaciones, comparte fotos con sus amigos y chatea en las redes sociales. Para él, el móvil es una ventana al mundo.

Pero un día, todo cambia. Un compañero de clase publica una foto humillante de Andrés en el grupo de WhatsApp de la clase. En pocas horas, la foto circula por todas partes y algunos alumnos escriben insultos y bromas crueles. Andrés sufre ciberacoso: no quiere ir al colegio, no duerme bien y se siente solo y triste. Además, se da cuenta de que pasa demasiado tiempo con el móvil: empieza a tener una verdadera adicción a las pantallas. También descubre que una aplicación que descargó la semana pasada utiliza sus datos personales sin su permiso. Su privacidad está en peligro.

Afortunadamente, Andrés decide hablar con su profesora de español. Ella le escucha con atención y le explica: «No compartas tus datos personales en internet y protege tu contraseña. No descargues aplicaciones desconocidas. Es importante que hables con un adulto y que denuncies el ciberacoso.»

Gracias a estos consejos, Andrés cambia sus hábitos: crea una contraseña segura, limita su tiempo delante de la pantalla a dos horas al día y habla con sus padres. El compañero que publicó la foto es castigado y pide perdón. Al final, Andrés comprende una idea esencial: las tecnologías digitales son muy útiles, pero hay que usarlas con inteligencia y con responsabilidad. Estar conectado está bien; estar protegido es todavía mejor.
\end{textebox}

\textbf{Traduction du texte (pour vérifier votre compréhension) :}

« Andrés a seize ans et vit à Yaoundé. Comme beaucoup de jeunes de son âge, il passe cinq heures par jour devant l'écran de son portable : il télécharge des applications, partage des photos avec ses amis et tchatche sur les réseaux sociaux. Pour lui, le portable est une fenêtre sur le monde.

Mais un jour, tout change. Un camarade de classe publie une photo humiliante d'Andrés dans le groupe WhatsApp de la classe. En quelques heures, la photo circule partout et certains élèves écrivent des insultes et des plaisanteries cruelles. Andrés subit du cyberharcèlement : il ne veut plus aller au collège, il ne dort pas bien et il se sent seul et triste. De plus, il se rend compte qu'il passe trop de temps avec son portable : il commence à avoir une véritable addiction aux écrans. Il découvre aussi qu'une application téléchargée la semaine dernière utilise ses données personnelles sans sa permission. Sa vie privée est en danger.

Heureusement, Andrés décide de parler avec sa professeure d'espagnol. Elle l'écoute attentivement et lui explique : "Ne partage pas tes données personnelles sur internet et protège ton mot de passe. Ne télécharge pas d'applications inconnues. Il est important que tu parles à un adulte et que tu dénonces le cyberharcèlement."

Grâce à ces conseils, Andrés change ses habitudes : il crée un mot de passe sûr, limite son temps d'écran à deux heures par jour et parle avec ses parents. Le camarade qui a publié la photo est puni et demande pardon. Au final, Andrés comprend une idée essentielle : les technologies numériques sont très utiles, mais il faut les utiliser avec intelligence et responsabilité. Être connecté, c'est bien ; être protégé, c'est encore mieux. »

\textbf{Glossaire (vocabulaire du texte) :}

\begin{center}
\begin{tabularx}{\linewidth}{|X|X|X|}
\hline
\rowcolor{popblueL} \textbf{Español} & \textbf{Français} & \textbf{Exemple du texte} \\
\hline
\esp{el ciberacoso} & le cyberharcèlement & \esp{Andrés sufre ciberacoso.} \\
\hline
\esp{la contraseña} & le mot de passe & \esp{protege tu contraseña} \\
\hline
\esp{descargar} & télécharger & \esp{descarga aplicaciones} \\
\hline
\esp{la privacidad} & la vie privée & \esp{Su privacidad está en peligro.} \\
\hline
\esp{la adicción} & l'addiction & \esp{una adicción a las pantallas} \\
\hline
\esp{compartir} & partager & \esp{comparte fotos con sus amigos} \\
\hline
\esp{la pantalla} & l'écran & \esp{delante de la pantalla} \\
\hline
\esp{la red social} & le réseau social & \esp{chatea en las redes sociales} \\
\hline
\esp{humillante} & humiliant & \esp{una foto humillante} \\
\hline
\esp{denunciar} & dénoncer, signaler & \esp{que denuncies el ciberacoso} \\
\hline
\esp{pedir perdón} & demander pardon & \esp{pide perdón} \\
\hline
\esp{castigar} & punir & \esp{es castigado} \\
\hline
\end{tabularx}
\end{center}

\begin{definition}[El ciberacoso]
Le \esp{ciberacoso} (cyberharcèlement) est l'ensemble des agressions — insultes, humiliations, diffusion d'images ou de rumeurs — commises contre une personne à travers internet ou le téléphone portable. À la différence du harcèlement classique (\esp{el acoso escolar}), il peut se produire à toute heure du jour et de la nuit, et la victime ne trouve aucun refuge, même chez elle.
\end{definition}

\begin{attention}[Faux amis et pièges de vocabulaire]
\begin{itemize}
\item \esp{la contraseña} signifie « le mot de passe », et non « la contrainte ». On dit aussi \esp{la clave} ou \esp{el código}.
\item \esp{descargar} signifie « télécharger » (d'internet vers ton appareil) ; « mettre en ligne » se dit \esp{subir} ou \esp{cargar}. Ne confonds pas les deux directions !
\item \esp{compartir} signifie « partager » ; ne dis jamais \esp{partir} (qui veut dire « partir, s'en aller »).
\item \esp{la aplicación} (souvent abrégé en \esp{la app}) est féminin : \esp{una aplicación útil}.
\item \esp{estar conectado} utilise \esp{estar} (état temporaire) ; mais dans la phrase finale du texte, \esp{estar protegido es todavía mejor}, l'infinitif \esp{estar} est le sujet de \esp{es} : « être protégé, c'est encore mieux ».
\end{itemize}
\end{attention}

\courssub{Repérage des informations}

Après la lecture, répondons aux quatre questions fondamentales du lecteur : \textbf{qui ? où ? quel problème ? quelles solutions ?}

\begin{center}
\begin{tabularx}{\linewidth}{|X|X|}
\hline
\rowcolor{popgreenL} \textbf{Pregunta} & \textbf{Respuesta} \\
\hline
\esp{¿Quién?} (qui ?) & \esp{Andrés, un adolescente de dieciséis años, y su profesora de español.} \\
\hline
\esp{¿Dónde?} (où ?) & \esp{En Yaundé, en el colegio y en el grupo de WhatsApp de la clase.} \\
\hline
\esp{¿Qué problema?} (quel problème ?) & \esp{Un compañero publica una foto humillante: Andrés sufre ciberacoso; además, tiene adicción a las pantallas y su privacidad está en peligro.} \\
\hline
\esp{¿Qué soluciones?} (quelles solutions ?) & \esp{Hablar con un adulto, proteger la contraseña, no compartir datos personales, limitar el tiempo de pantalla y denunciar.} \\
\hline
\end{tabularx}
\end{center}

\begin{popfigure}
\begin{center}\tcbox[colback=popblue,colframe=popblue,coltext=white,arc=9pt,boxsep=4pt,left=6pt,right=6pt]{\bfseries\small Conectados, ¿pero protegidos?}\end{center}
\vspace{-2pt}
\begin{tcbraster}[raster columns=2,raster equal height=rows,raster column skip=6pt,raster row skip=6pt,enhanced,blankest]
\begin{tcolorbox}[enhanced,colback=poporange,colframe=poporange,coltext=white,boxrule=0.9pt,arc=3pt,left=4pt,right=4pt,top=3pt,bottom=3pt,boxsep=1pt,halign=left]\bfseries\scriptsize el problema \esp{ciberacoso}\end{tcolorbox}
\begin{tcolorbox}[enhanced,colback=white,colframe=poppink,boxrule=0.9pt,arc=3pt,title={los riesgos},fonttitle=\bfseries\scriptsize,coltitle=white,colbacktitle=poppink,left=4pt,right=4pt,top=2pt,bottom=2pt,boxsep=1pt,toptitle=1.5pt,bottomtitle=1.5pt,halign=left]
\begin{itemize}[nosep,leftmargin=*,itemsep=1pt,font=\scriptsize]
\item \esp{adicción}
\item \esp{privacidad}
\end{itemize}
\end{tcolorbox}
\begin{tcolorbox}[enhanced,colback=white,colframe=popgreen,boxrule=0.9pt,arc=3pt,title={los consejos},fonttitle=\bfseries\scriptsize,coltitle=white,colbacktitle=popgreen,left=4pt,right=4pt,top=2pt,bottom=2pt,boxsep=1pt,toptitle=1.5pt,bottomtitle=1.5pt,halign=left]
\begin{itemize}[nosep,leftmargin=*,itemsep=1pt,font=\scriptsize]
\item \esp{contraseña} \esp{segura}
\item \esp{no compartir} \esp{datos}
\end{itemize}
\end{tcolorbox}
\begin{tcolorbox}[enhanced,colback=popgold,colframe=popgold,coltext=white,boxrule=0.9pt,arc=3pt,left=4pt,right=4pt,top=3pt,bottom=3pt,boxsep=1pt,halign=left]\bfseries\scriptsize la solución \esp{hablar con} \esp{un adulto}\end{tcolorbox}
\end{tcbraster}

\legende{Carte mentale du texte : problème, risques, conseils et solution.}
\end{popfigure}

\courssub{Questions de compréhension}

Répondez par des phrases complètes, en espagnol, en reprenant les mots de la question.

\begin{enumerate}
\item \esp{¿Qué le pasó a Andrés en el grupo de WhatsApp?}
\item \esp{¿Cuántas horas pasa Andrés delante de la pantalla?}
\item \esp{¿Qué consejos le da su profesora?}
\item \esp{¿Por qué es importante proteger la contraseña?}
\item \esp{¿Qué decide hacer Andrés al final?}
\end{enumerate}

\begin{exoresolu}[Questions de compréhension : corrigé modèle]
Réponds aux cinq questions de compréhension ci-dessus par des phrases complètes.
\tcblower
\textbf{1.} \esp{Un compañero publica una foto humillante de él en el grupo de WhatsApp y algunos alumnos escriben insultos: Andrés sufre ciberacoso.} « Un camarade publie une photo humiliante de lui et certains élèves écrivent des insultes. »

\textbf{2.} \esp{Andrés pasa cinco horas al día delante de la pantalla de su móvil.} « Andrés passe cinq heures par jour devant l'écran. » Remarque : avec \esp{pasar} + durée, on utilise la structure \esp{pasa cinco horas al día} (cinq heures par jour).

\textbf{3.} \esp{Su profesora le dice: no compartas tus datos personales, protege tu contraseña, no descargues aplicaciones desconocidas y habla con un adulto.} « Sa professeure lui dit : ne partage pas tes données, protège ton mot de passe, ne télécharge pas d'applications inconnues et parle à un adulte. »

\textbf{4.} \esp{Es importante proteger la contraseña porque, si no, otras personas pueden entrar en nuestras cuentas y robar nuestros datos personales.} « Sinon, d'autres personnes peuvent entrer dans nos comptes et voler nos données. »

\textbf{5.} \esp{Al final, Andrés decide cambiar sus hábitos: crea una contraseña segura, limita su tiempo de pantalla y habla con sus padres.} « À la fin, Andrés décide de changer ses habitudes. »
\end{exoresolu}

\courssub{Méthode : lire un texte narratif à visée argumentative}

Notre texte raconte une histoire, mais il défend aussi une idée : il faut utiliser le numérique avec responsabilité. Voici la démarche à suivre pour lire ce type de texte, très fréquent à l'examen.

\begin{methode}[Lire et comprendre un texte en quatre étapes]
\begin{enumerate}
\item \textbf{Identifier le thème.} Pose la question : « De quoi parle le texte ? » Ici : les technologies numériques et leurs risques. Le titre \esp{Conectados, ¿pero protegidos?} annonce déjà le débat.
\item \textbf{Repérer la thèse et les arguments.} La thèse est l'idée défendue : \esp{las tecnologías son útiles, pero hay que usarlas con responsabilidad}. Les arguments sont les exemples qui la prouvent : le \esp{ciberacoso}, la \esp{adicción}, le vol de données.
\item \textbf{Relever le lexique thématique.} Classe les mots par familles : les outils (\esp{móvil, aplicación, red social}), les risques (\esp{ciberacoso, adicción, privacidad}), les solutions (\esp{contraseña, denunciar, proteger}).
\item \textbf{Dégager la visée de l'auteur.} Demande-toi : « Pourquoi l'auteur écrit-il ce texte ? » Ici : informer les jeunes, les mettre en garde et proposer des comportements responsables.
\end{enumerate}
\end{methode}

\courssub{Reformulation : résumer le texte en trois phrases}

Savoir résumer est une compétence clé de l'expression écrite. Consigne : résumez le texte en \textbf{trois phrases} en utilisant les connecteurs chronologiques \esp{primero} (d'abord), \esp{luego} (ensuite), \esp{finalmente} (finalement).

\begin{exemplebox}[Modèle de résumé en trois phrases]
\esp{Primero, Andrés, un joven de Yaundé, pasa cinco horas al día con su móvil y sufre ciberacoso cuando un compañero publica una foto humillante en WhatsApp.}

« D'abord, Andrés, un jeune de Yaoundé, passe cinq heures par jour sur son portable et subit du cyberharcèlement quand un camarade publie une photo humiliante sur WhatsApp. »

\esp{Luego, habla con su profesora, que le da consejos: proteger su contraseña, no compartir sus datos personales y denunciar el acoso.}

« Ensuite, il parle avec sa professeure, qui lui donne des conseils : protéger son mot de passe, ne pas partager ses données personnelles et dénoncer le harcèlement. »

\esp{Finalmente, Andrés cambia sus hábitos y comprende que hay que usar las tecnologías con responsabilidad.}

« Finalement, Andrés change ses habitudes et comprend qu'il faut utiliser les technologies avec responsabilité. »
\end{exemplebox}

\begin{attention}[Erreurs fréquentes des francophones dans le résumé]
\begin{itemize}
\item Ne traduis pas mot à mot : « passer du temps » se dit \esp{pasar tiempo}, avec un seul \esp{s} ; n'écris jamais une forme avec \esp{-ss-} sur le modèle français.
\item « Il y a cinq heures » se dit \esp{hace cinco horas} ; mais « cinq heures par jour » se dit \esp{cinco horas al día} (\esp{al} = \esp{a} + \esp{el}).
\item Les connecteurs \esp{primero, luego, finalmente} se placent en début de phrase, suivis d'une virgule, comme en français.
\item Attention au genre : \esp{el problema}, \esp{el móvil}, \esp{el colegio} sont masculins ; \esp{la pantalla}, \esp{la contraseña}, \esp{la privacidad} sont féminins.
\end{itemize}
\end{attention}

\begin{savaistu}[Le cyberharcèlement et la loi dans le monde hispanique]
Selon les études de l'UNESCO sur la violence à l'école, environ un adolescent sur dix dans le monde déclare avoir déjà été victime de cyberharcèlement. L'Espagne a adopté en 2018 une loi organique de protection des données personnelles (la \esp{LOPDGDD}, \esp{Ley Orgánica de Protección de Datos}) qui fixe à quatorze ans l'âge à partir duquel un jeune peut consentir seul au traitement de ses données sur les réseaux sociaux. En Amérique latine, plusieurs pays comme le Mexique et la Colombie ont aussi créé des lois contre le \esp{ciberacoso}. Et au Cameroun, la loi de 2010 sur la cybersécurité et la cybercriminalité punit déjà les atteintes aux personnes commises par voie numérique : le harcèlement en ligne n'est donc pas un jeu, c'est un délit !
\end{savaistu}

\begin{exercice}[À toi de jouer : vrai ou faux]
Réponds par \esp{verdadero} (vrai) ou \esp{falso} (faux) et justifie avec une phrase du texte.

\begin{enumerate}
\item \esp{Andrés vive en Duala.}
\item \esp{Andrés pasa dos horas al día con su móvil al principio del texto.}
\item \esp{Un compañero publica una foto humillante en el grupo de WhatsApp.}
\item \esp{La profesora aconseja a Andrés que comparta sus datos personales.}
\item \esp{Al final, Andrés limita su tiempo delante de la pantalla.}
\end{enumerate}
\end{exercice}

\begin{corrige}[Exercice : vrai ou faux]
\textbf{1.} \esp{Falso}: \esp{Andrés vive en Yaundé}, pas à Douala.

\textbf{2.} \esp{Falso}: au début, \esp{pasa cinco horas al día delante de la pantalla} ; les deux heures, c'est à la fin, après les conseils.

\textbf{3.} \esp{Verdadero}: \esp{Un compañero de clase publica una foto humillante de Andrés en el grupo de WhatsApp de la clase.}

\textbf{4.} \esp{Falso}: la professeure dit le contraire : \esp{No compartas tus datos personales}.

\textbf{5.} \esp{Verdadero}: \esp{limita su tiempo delante de la pantalla a dos horas al día}.
\end{corrige}

\begin{aretenir}[Ce qu'il faut retenir de cette première étape]
\begin{itemize}
\item Le texte raconte l'histoire d'Andrés, victime de \esp{ciberacoso}, et défend une thèse : les technologies sont utiles mais il faut les utiliser avec \esp{responsabilidad}.
\item Le lexique clé du monde numérique : \esp{móvil, aplicación, contraseña, red social, pantalla, descargar, compartir, privacidad, adicción, ciberacoso}.
\item Pour lire un texte : identifier le thème, repérer la thèse et les arguments, relever le lexique, dégager la visée de l'auteur.
\item Pour résumer : trois phrases maximum, avec \esp{primero, luego, finalmente}.
\item Dans le texte, la professeure conseille avec l'impératif (\esp{no compartas, protege}) et le subjonctif (\esp{es importante que hables}) : ce sera l'objet des deux prochaines sections de grammaire.
\end{itemize}
\end{aretenir}

\coursec{Lexique thématique : el mundo digital}

Dans la section précédente, nous avons lu et commenté le texte \esp{Conectados, ¿pero protegidos?}. Ce texte regorgeait de mots nouveaux : \esp{móvil}, \esp{contraseña}, \esp{ciberacoso}, \esp{descargar}... Pour bien comprendre un article de presse sur le numérique, pour en parler à l'oral et pour le traduire, il faut d'abord maîtriser ce vocabulaire. C'est l'objectif de cette section : organiser, comprendre et mémoriser le lexique du monde digital, classé en cinq champs sémantiques.

\courssub{Observer avant d'apprendre : le lexique en contexte}

Avant toute liste de mots, lisons ce court dialogue entre deux élèves d'un lycée de Yaoundé. Soulignez mentalement tous les mots liés au numérique.

\begin{textebox}[Diálogo : « El móvil de Carlos »]
— Carlos, ¿por qué estás tan cansado?\\
— Porque pasé cuatro horas delante de la pantalla anoche. Descargué una aplicación nueva y estuve chateando hasta las dos de la mañana.\\
— ¡Cuidado! Estás enganchado al móvil. ¿Y tu contraseña? ¿La compartes con tus amigos?\\
— Sí, con mi mejor amigo. ¿Por qué?\\
— ¡Qué peligroso! Debes proteger tus datos personales y cambiar la configuración de privacidad. Y si alguien te molesta en las redes sociales, bloquéalo y denúncialo.\\
— Tienes razón. Voy a limitar mi tiempo de pantalla y a pensar antes de publicar.
\end{textebox}

\textbf{Glossaire} : \esp{estar enganchado a} = être accro à ; \esp{peligroso} = dangereux ; \esp{molestar} = importuner, harceler ; \esp{bloquéalo} = bloque-le (impératif + pronom) ; \esp{denúncialo} = dénonce-le ; \esp{tener razón} = avoir raison.

Comptez : ce dialogue de huit répliques contient plus de quinze mots de notre lexique. Observons maintenant ces mots champ par champ.

\courssub{Champ 1 — Les outils (las herramientas)}

Ce sont les objets et les supports du monde numérique.

\begin{center}
\begin{tabularx}{\linewidth}{|l|X|X|}
\hline
\rowcolor{popblueL} Español & Français & Remarque \\
\hline
\esp{internet} & internet & souvent sans article : \esp{navegar por internet} \\
\hline
\esp{el móvil / el celular} & le portable, le téléphone mobile & \esp{celular} en Amérique latine \\
\hline
\esp{el ordenador} & l'ordinateur & \esp{la computadora} en Amérique latine \\
\hline
\esp{la tableta} & la tablette & féminin \\
\hline
\esp{la aplicación (la app)} & l'application & féminin : tous les noms en \esp{-ción} le sont \\
\hline
\esp{la red social} & le réseau social & pluriel : \esp{las redes sociales} \\
\hline
\esp{la pantalla} & l'écran & \esp{tiempo de pantalla} = temps d'écran \\
\hline
\esp{el correo electrónico} & le courriel, l'e-mail & on dit aussi \esp{el email} \\
\hline
\end{tabularx}
\end{center}

\begin{exemplebox}[Les outils en phrases]
\esp{Uso el ordenador para hacer mis tareas.} « J'utilise l'ordinateur pour faire mes devoirs. »

\esp{Descargo muchas aplicaciones en mi móvil.} « Je télécharge beaucoup d'applications sur mon portable. »

\esp{Las redes sociales son muy populares entre los jóvenes cameruneses.} « Les réseaux sociaux sont très populaires chez les jeunes Camerounais. »
\end{exemplebox}

\courssub{Champ 2 — Les actions (las acciones)}

Ce sont les verbes de la vie numérique. Presque tous sont réguliers en \esp{-ar} ou \esp{-ir}, ce qui est une bonne nouvelle pour la conjugaison.

\begin{center}
\begin{tabularx}{\linewidth}{|l|X|X|}
\hline
\rowcolor{popgreenL} Español & Français & Exemple (1\textsuperscript{re} pers. sg.) \\
\hline
\esp{descargar} & télécharger (depuis internet) & \esp{Descargo música.} \\
\hline
\esp{compartir} & partager & \esp{Comparto fotos.} \\
\hline
\esp{publicar} & publier, poster & \esp{Publico un mensaje.} \\
\hline
\esp{chatear} & tchatter, discuter en ligne & \esp{Chateamos por la noche.} \\
\hline
\esp{conectarse} & se connecter & \esp{Me conecto a internet.} \\
\hline
\esp{navegar por internet} & naviguer sur internet & \esp{Navego por internet.} \\
\hline
\esp{enviar mensajes} & envoyer des messages & \esp{Envío mensajes a mis amigos.} \\
\hline
\esp{bloquear} & bloquer (un contact) & \esp{Bloqueo al desconocido.} \\
\hline
\esp{denunciar} & signaler, dénoncer & \esp{Denuncio el ciberacoso.} \\
\hline
\end{tabularx}
\end{center}

\begin{attention}[Descargar ou cargar ?]
Ne confondez pas \esp{descargar} (télécharger un fichier, une application) et \esp{cargar} (charger la batterie du téléphone). \esp{Descargo una aplicación} = je télécharge une application ; \esp{Cargo el móvil} = je charge le portable (branchement). Le préfixe \esp{des-} indique ici le transfert vers soi (téléchargement). Autre piège : \esp{enviar} se conjugue avec un accent sur le \esp{í} aux personnes où la tonique tombe sur le radical : \esp{envío, envías, envía} (même mécanisme que \esp{confiar} : \esp{confío}).
\end{attention}

\courssub{Champ 3 — La sécurité (la seguridad)}

\begin{definition}[La privacidad]
\esp{La privacidad} désigne le droit de contrôler ses informations personnelles (\esp{fotos, datos, contraseñas}) sur internet : qui peut les voir, les utiliser, les partager. En français : la vie privée, la confidentialité.
\end{definition}

\begin{center}
\begin{tabularx}{\linewidth}{|l|X|X|}
\hline
\rowcolor{popgoldL} Español & Français & Exemple \\
\hline
\esp{la contraseña} & le mot de passe & \esp{Cambia tu contraseña.} \\
\hline
\esp{la privacidad} & la vie privée, la confidentialité & \esp{Respeta mi privacidad.} \\
\hline
\esp{los datos personales} & les données personnelles & \esp{No publiques tus datos personales.} \\
\hline
\esp{proteger} & protéger & \esp{Protege tu contraseña.} \\
\hline
\esp{la configuración de privacidad} & les paramètres de confidentialité & \esp{Revisa la configuración de privacidad.} \\
\hline
\esp{el usuario / la usuaria} & l'utilisateur / l'utilisatrice & \esp{El usuario debe registrarse.} \\
\hline
\end{tabularx}
\end{center}

\begin{exemplebox}[La sécurité en phrases]
\esp{Protege tu contraseña y no la compartas con nadie.} « Protège ton mot de passe et ne le partage avec personne. »

\esp{Los usuarios jóvenes deben revisar la configuración de privacidad de sus cuentas.} « Les jeunes utilisateurs doivent vérifier les paramètres de confidentialité de leurs comptes. »
\end{exemplebox}

\courssub{Champ 4 — Les risques (los riesgos)}

\begin{center}
\begin{tabularx}{\linewidth}{|l|X|X|}
\hline
\rowcolor{poppinkL} Español & Français & Exemple \\
\hline
\esp{el ciberacoso} & le cyberharcèlement & \esp{El ciberacoso es un problema grave.} \\
\hline
\esp{la adicción} & l'addiction, la dépendance & \esp{La adicción al móvil afecta el sueño.} \\
\hline
\esp{las noticias falsas} & les fausses informations (fake news) & \esp{No creas las noticias falsas.} \\
\hline
\esp{el robo de datos} & le vol de données & \esp{El robo de datos es frecuente.} \\
\hline
\esp{el contenido inapropiado} & le contenu inapproprié & \esp{Hay contenido inapropiado en algunas páginas.} \\
\hline
\esp{el aislamiento} & l'isolement & \esp{El exceso de pantalla causa aislamiento.} \\
\hline
\end{tabularx}
\end{center}

Notez la formation des mots composés avec le préfixe \esp{ciber-} : \esp{ciberacoso} (cyberharcèlement), \esp{ciberseguridad} (cybersécurité), \esp{ciberdelincuente} (cyberdélinquant). Ce préfixe, très productif, s'écrit en un seul mot, sans trait d'union.

\courssub{Champ 5 — Les bons comportements (los buenos comportamientos)}

\begin{itemize}
\item \esp{hablar con un adulto} = parler à un adulte (en cas de problème) ;
\item \esp{limitar el tiempo de pantalla} = limiter le temps d'écran ;
\item \esp{respetar a los demás} = respecter les autres ;
\item \esp{verificar la información} = vérifier l'information avant de la croire ou de la partager ;
\item \esp{pensar antes de publicar} = réfléchir avant de publier.
\end{itemize}

\begin{exemplebox}[Expressions utiles]
\esp{Paso muchas horas delante de la pantalla.} « Je passe beaucoup d'heures devant l'écran. »

\esp{Está enganchado al móvil.} « Il est accro à son portable. » (attention : \esp{enganchado a} + article contracté \esp{al} devant un nom masculin)

\esp{Voy a desconectar esta noche.} « Je vais me déconnecter ce soir. »

\esp{Debemos usar el móvil con responsabilidad.} « Nous devons utiliser le portable de manière responsable. »
\end{exemplebox}

\courssub{Carte mentale et tableau récapitulatif}

Voici le tableau à double entrée du lexique, à recopier et à compléter dans votre cahier au fur et à mesure de l'année.

\begin{popfigure}
\begin{tikzpicture}[scale=0.92, every node/.style={transform shape}]
\draw[fill=popblueL] (0,0) rectangle (2.6,1);
\node[align=center] at (1.3,0.5) {\textbf{Herramientas}};
\draw[fill=popcream] (0,1) rectangle (2.6,4.6);
\node[align=left, anchor=north west, text width=2.4cm] at (0.05,4.55) {el móvil\\ el ordenador\\ la tableta\\ la aplicación\\ la red social\\ la pantalla\\ el correo};

\draw[fill=popgreenL] (2.7,0) rectangle (5.3,1);
\node[align=center] at (4,0.5) {\textbf{Acciones}};
\draw[fill=popcream] (2.7,1) rectangle (5.3,4.6);
\node[align=left, anchor=north west, text width=2.4cm] at (2.75,4.55) {descargar\\ compartir\\ publicar\\ chatear\\ conectarse\\ navegar\\ bloquear\\ denunciar};

\draw[fill=popgoldL] (5.4,0) rectangle (8,1);
\node[align=center] at (6.7,0.5) {\textbf{Seguridad}};
\draw[fill=popcream] (5.4,1) rectangle (8,4.6);
\node[align=left, anchor=north west, text width=2.4cm] at (5.45,4.55) {la contraseña\\ la privacidad\\ los datos\\ proteger\\ el usuario};

\draw[fill=poppinkL] (8.1,0) rectangle (10.7,1);
\node[align=center] at (9.4,0.5) {\textbf{Riesgos}};
\draw[fill=popcream] (8.1,1) rectangle (10.7,4.6);
\node[align=left, anchor=north west, text width=2.4cm] at (8.15,4.55) {el ciberacoso\\ la adicción\\ noticias falsas\\ robo de datos\\ aislamiento};

\draw[fill=poppurpleL] (10.8,0) rectangle (13.4,1);
\node[align=center, text width=2.5cm] at (12.1,0.5) {\textbf{Buenos comporta- mientos}};
\draw[fill=popcream] (10.8,1) rectangle (13.4,4.6);
\node[align=left, anchor=north west, text width=2.4cm] at (10.85,4.55) {hablar con un adulto\\ limitar la pantalla\\ respetar\\ verificar\\ pensar antes};
\end{tikzpicture}
\legende{Tableau à double entrée : les cinq champs du monde digital}
\end{popfigure}

\begin{aretenir}[Les points clés du lexique]
\begin{itemize}
\item \esp{La aplicación} est féminin : tous les noms en \esp{-ción} sont féminins (\esp{la configuración}, \esp{la adicción}, \esp{la información}).
\item \esp{Internet} s'emploie souvent sans article : \esp{navegar por internet}, \esp{estar en internet}.
\item \esp{El móvil} (Espagne) = \esp{el celular} (Amérique latine) ; les deux se comprennent partout.
\item \esp{Descargar} = télécharger ; \esp{cargar} = charger la batterie.
\item Le préfixe \esp{ciber-} s'écrit collé : \esp{ciberacoso}, \esp{ciberseguridad}.
\end{itemize}
\end{aretenir}

\begin{savaistu}[El móvil o el celular ?]
En Espagne, on dit \esp{el móvil} ; dans la plupart des pays d'Amérique latine (Mexique, Colombie, Argentine...), on préfère \esp{el celular}. En Guinée équatoriale, le seul pays hispanophone d'Afrique subsaharienne, voisin du Cameroun, les deux formes coexistent dans la presse et la conversation. De même, \esp{el ordenador} (Espagne) devient souvent \esp{la computadora} ou \esp{la computador} outre-Atlantique. Un bon hispanophone comprend toutes ces variantes !
\end{savaistu}

\courssub{Exercice résolu et activités de fixation}

\begin{exoresolu}[Compléter avec le mot juste]
Énoncé — Complétez les phrases avec : \esp{contraseña}, \esp{descargar}, \esp{ciberacoso}, \esp{pantalla}, \esp{verificar}.

1. No debes compartir tu \ldots{} con tus amigos.
2. Quiero \ldots{} una aplicación de música.
3. El \ldots{} es un problema grave en las redes sociales.
4. Pasas demasiadas horas delante de la \ldots{}.
5. Antes de compartir una noticia, hay que \ldots{} la información.

\tcblower
Solution détaillée :
1. \esp{contraseña} : on ne partage jamais son mot de passe (sécurité).
2. \esp{descargar} : on télécharge une application (infinitif après \esp{querer}).
3. \esp{ciberacoso} : le cyberharcèlement, problème grave des réseaux sociaux ; le mot est masculin, d'où \esp{el}.
4. \esp{pantalla} : expression figée \esp{delante de la pantalla} (devant l'écran).
5. \esp{verificar} : bon comportement numéro 4 : vérifier l'information avant de la partager.
\end{exoresolu}

\begin{exercice}[Activité 1 : classement dans les cinq champs]
Classez ces 15 mots dans les cinq colonnes du tableau (\emph{Herramientas / Acciones / Seguridad / Riesgos / Buenos comportamientos}) :

\esp{la tableta} — \esp{bloquear} — \esp{la contraseña} — \esp{la adicción} — \esp{respetar a los demás} — \esp{el correo electrónico} — \esp{chatear} — \esp{la privacidad} — \esp{las noticias falsas} — \esp{limitar el tiempo de pantalla} — \esp{la red social} — \esp{denunciar} — \esp{el usuario} — \esp{el aislamiento} — \esp{verificar la información}.
\end{exercice}

\begin{corrige}[Activité 1]
\emph{Herramientas} : la tableta, el correo electrónico, la red social. \emph{Acciones} : bloquear, chatear, denunciar. \emph{Seguridad} : la contraseña, la privacidad, el usuario. \emph{Riesgos} : la adicción, las noticias falsas, el aislamiento. \emph{Buenos comportamientos} : respetar a los demás, limitar el tiempo de pantalla, verificar la información.
\end{corrige}

\begin{experience}[Activité 2 : jeu de devinettes « definición → palabra »]
Par groupes de deux. L'élève A lit une définition en espagnol, l'élève B doit trouver le mot du lexique. Exemples de définitions :

1. \esp{Es una palabra secreta para entrar en tu cuenta.} (→ \esp{la contraseña})
2. \esp{Es cuando una persona molesta a otra por internet, muchas veces.} (→ \esp{el ciberacoso})
3. \esp{Es lo que haces cuando quieres tener una aplicación nueva en tu móvil.} (→ \esp{descargar})
4. \esp{Es una persona que utiliza internet o una aplicación.} (→ \esp{el usuario})
5. \esp{Es lo contrario de compartir tus fotos con todo el mundo.} (→ \esp{la privacidad})

Ensuite, inversez les rôles et inventez vos propres définitions avec les autres mots du tableau.
\end{experience}

\begin{attention}[Erreurs fréquentes des francophones]
1. Ne dites pas \esp{la internet} : en espagnol courant, \esp{navegar por internet}, \esp{buscar en internet}, sans article.
2. \esp{La aplicación} : féminin, comme tous les noms en \esp{-ción}. On dit \esp{una aplicación útil}.
3. « Être accro à » se traduit par \esp{estar enganchado a} (avec \esp{estar}, car c'est un état) : \esp{Está enganchado al móvil}, jamais \esp{es enganchado}.
4. « Télécharger » = \esp{descargar} ou \esp{bajar} ; ne jamais calquer le français avec \esp{telecargar}, qui n'existe pas.
5. \esp{Los datos} est masculin pluriel : \esp{mis datos personales}, pas \esp{las datas}.
\end{attention}

Ce lexique est désormais votre boîte à outils. Dans les sections suivantes, nous l'utiliserons pour construire des conseils à l'impératif (\esp{no compartas tu contraseña}) et au subjonctif (\esp{es importante que protejas tus datos}), afin de proposer des comportements numériques responsables.

\coursec{Grammaire 1 : l'impératif pour conseiller}

Après avoir exploré le lexique du monde numérique (\emph{el mundo digital}) dans la section précédente, nous disposons maintenant du vocabulaire nécessaire pour formuler des conseils concrets : \esp{proteger la contraseña}, \esp{evitar el ciberacoso}, \esp{controlar la adicción}. En espagnol, pour exprimer une recommandation, un ordre ou une consigne adressée à quelqu'un, on utilise le \cle{mode impératif}. C'est la forme verbale privilégiée dans les campagnes de prévention, les tutoriels et les dialogues éducatifs. Observons d'abord comment il fonctionne dans notre texte support \emph{« Conectados, ¿pero protegidos? »}.

\courssub{Observation dans le texte support}

Reprenons trois phrases extraites du texte lu en Section 1. Elles illustrent les deux visages de l'impératif : l'affirmatif (faire quelque chose) et le négatif (ne pas faire quelque chose).

\begin{textebox}[Extraits du texte : Conseils de sécurité]
\textbf{Consejo 1:} \esp{Protege tu contraseña con caracteres especiales.}\par
\textbf{Consejo 2:} \esp{No compartas tus datos personales con extraños.}\par
\textbf{Consejo 3:} \esp{Habla con un adulto de confianza si sufres ciberacoso.}
\end{textebox}

\textbf{Glossaire :} \esp{caracteres especiales} (caractères spéciaux) ; \esp{extraños} (inconnus / étrangers) ; \esp{sufres} (tu subis, présent indicatif) ; \esp{ciberacoso} (cyberharcèlement).

\emph{Analyse :}
\begin{itemize}
    \item \esp{Protege} et \esp{Habla} terminent par \cle{-e}. Ce sont des formes affirmatives adressées à \cle{tú} (toi).
    \item \esp{No compartas} commence par \cle{no} et se termine par \cle{-as}. C'est une forme négative adressée à \cle{tú}.
    \item La fonction commune est le \cle{conseil} / la \cle{recommandation} de comportement responsable.
\end{itemize}

\courssub{Formation de l'impératif affirmatif (tú)}

\begin{propriete}[Règle de base : Impératif affirmatif \textit{tú}]
Pour les verbes réguliers, l'impératif affirmatif de la 2\textsuperscript{e} personne du singulier (\textit{tú}) est \textbf{identique à la 3\textsuperscript{e} personne du singulier du présent de l'indicatif} (\textit{él/ella/usted}).
\end{propriete}

\begin{center}
\begin{tabularx}{\linewidth}{|X|X|X|X|}
\hline
\rowcolor{popblueL}
\textbf{Infinitif} & \textbf{Présent indicatif (él/ella)} & $\rightarrow$ & \textbf{Impératif affirmatif (tú)} \\
\hline
\esp{compartir} & \esp{comparte} & $\rightarrow$ & \esp{¡Comparte!} (Partage !) \\
\hline
\esp{proteger} & \esp{protege} & $\rightarrow$ & \esp{¡Protege!} (Protège !) \\
\hline
\esp{hablar} & \esp{habla} & $\rightarrow$ & \esp{¡Habla!} (Parle !) \\
\hline
\esp{escribir} & \esp{escribe} & $\rightarrow$ & \esp{¡Escribe!} (Écris !) \\
\hline
\esp{descargar} & \esp{descarga} & $\rightarrow$ & \esp{¡Descarga!} (Télécharge !) \\
\hline
\end{tabularx}
\end{center}

\emph{Remarque :} L'accentuation suit la règle générale (paroxytones terminées par voyelle, \emph{n}, \emph{s} $\rightarrow$ pas d'accent écrit : \esp{habla}, \esp{protege}, \esp{comparte}). Seuls les verbes en \emph{-ir} comme \esp{escribir} $\rightarrow$ \esp{escribe} ou \esp{vivir} $\rightarrow$ \esp{vive} gardent l'accent sur la radicale si le radical le porte (ex. \esp{huir} $\rightarrow$ \esp{huye}).

\begin{attention}[Les 8 impératifs affirmatifs irréguliers (tú) — \textbf{À apprendre par cœur}]
Ces verbes très fréquents ne suivent pas la règle de la 3\textsuperscript{e} personne. Ils sont courts et portent souvent un accent diacritique pour les distinguer d'autres mots (ex. \esp{sé} vs \esp{se}).
\begin{center}
\begin{tabular}{|l|l|l|l|l|l|l|l|l|}
\hline
\rowcolor{poporangeL}
\textbf{decir} & \textbf{hacer} & \textbf{ir} & \textbf{poner} & \textbf{salir} & \textbf{ser} & \textbf{tener} & \textbf{venir} \\
\hline
\esp{di} & \esp{haz} & \esp{ve} & \esp{pon} & \esp{sal} & \esp{sé} & \esp{ten} & \esp{ven} \\
\hline
\end{tabular}
\end{center}
\emph{Astuce mnémotechnique :} \textbf{« Di haz ve pon sal sé ten ven »} (Dis fais va mets sors sois tiens viens). Répétez-les en rythme.
\end{attention}

\begin{exemplebox}[Exemples affirmatifs avec verbes irréguliers]
\begin{itemize}
    \item \esp{Di la verdad a tus padres.} « Dis la vérité à tes parents. » (\esp{decir} $\rightarrow$ \esp{di})
    \item \esp{Haz la tarea antes de jugar.} « Fais tes devoirs avant de jouer. » (\esp{hacer} $\rightarrow$ \esp{haz})
    \item \esp{Ve al colegio temprano.} « Va à l'école tôt. » (\esp{ir} $\rightarrow$ \esp{ve})
    \item \esp{Pon el móvil en modo avión.} « Mets le portable en mode avion. » (\esp{poner} $\rightarrow$ \esp{pon})
    \item \esp{Sé responsable en redes.} « Sois responsable sur les réseaux. » (\esp{ser} $\rightarrow$ \esp{sé})
    \item \esp{Ten cuidado con los enlaces.} « Fais attention aux liens. » (\esp{tener cuidado} $\rightarrow$ \esp{ten})
    \item \esp{Ven a la charla sobre ciberseguridad.} « Viens à la conférence sur la cybersécurité. » (\esp{venir} $\rightarrow$ \esp{ven})
\end{itemize}
\end{exemplebox}

\courssub{Formation de l'impératif négatif (tú)}

C'est ici que se situe la difficulté majeure pour les francophones. En français, l'interdiction ou le conseil négatif s'exprime souvent à l'infinitif (\emph{Ne pas partager}, \emph{Ne pas cliquer}). En espagnol, \textbf{l'impératif négatif n'existe pas sous forme simple} : il emprunte obligatoirement les formes du \cle{subjonctif présent}.

\begin{propriete}[Règle d'or : Impératif négatif = \textbf{no} + Subjonctif présent (tú)]
Pour interdire ou déconseiller (\textit{ne fais pas...}), on place \textbf{no} devant le verbe conjugué au \textbf{subjonctif présent} à la 2\textsuperscript{e} personne du singulier (\textit{tú}).
\end{propriete}

\begin{center}
\begin{tabularx}{\linewidth}{|X|X|X|X|}
\hline
\rowcolor{popgreenL}
\textbf{Infinitif} & \textbf{Subjonctif présent (tú)} & $\rightarrow$ & \textbf{Impératif négatif (tú)} \\
\hline
\esp{compartir} & \esp{compartas} & $\rightarrow$ & \esp{¡No compartas!} (Ne partage pas !) \\
\hline
\esp{proteger} & \esp{protejas} & $\rightarrow$ & \esp{¡No protejas!} (Ne protège pas !) \\
\hline
\esp{hablar} & \esp{hables} & $\rightarrow$ & \esp{¡No hables!} (Ne parle pas !) \\
\hline
\esp{escribir} & \esp{escribas} & $\rightarrow$ & \esp{¡No escribas!} (N'écris pas !) \\
\hline
\esp{publicar} & \esp{publiques} & $\rightarrow$ & \esp{¡No publiques!} (Ne publie pas !) \\
\hline
\esp{abrir} & \esp{abras} & $\rightarrow$ & \esp{¡No abras!} (N'ouvre pas !) \\
\hline
\esp{descargar} & \esp{descargues} & $\rightarrow$ & \esp{¡No descargues!} (Ne télécharge pas !) \\
\hline
\end{tabularx}
\end{center}

\emph{Observation des voyelles thématiques (subjonctif = voyelle inverse) :}
\begin{itemize}
    \item Verbes en \textbf{-ar} (\esp{hablar}, \esp{publicar}) $\rightarrow$ subjonctif en \textbf{-e} (\esp{hables}, \esp{publiques}).
    \item Verbes en \textbf{-er/-ir} (\esp{compartir}, \esp{proteger}, \esp{escribir}, \esp{abrir}, \esp{descargar}) $\rightarrow$ subjonctif en \textbf{-a} (\esp{compartas}, \esp{protejas}, \esp{escribas}, \esp{abras}, \esp{descargues}).
\end{itemize}

\begin{exemplebox}[Exemples négatifs traduits (contexte numérique)]
\begin{itemize}
    \item \esp{No publiques fotos sin permiso.} « Ne publie pas de photos sans permission. »
    \item \esp{No abras correos de remitentes desconocidos.} « N'ouvre pas les mails d'expéditeurs inconnus. »
    \item \esp{No descargues aplicaciones de fuentes no oficiales.} « Ne télécharge pas d'applications de sources non officielles. »
    \item \esp{No pases tantas horas delante de la pantalla.} « Ne passe pas tant d'heures devant l'écran. »
    \item \esp{Verifica la información antes de reenviar.} (Affirmatif) « Vérifie l'info avant de la transférer. »
    \item \esp{No te fíes de perfiles sin foto.} « Ne te fie pas aux profils sans photo. » (Réfléchi \esp{fiarses} $\rightarrow$ \esp{te fíes})
\end{itemize}
\end{exemplebox}

\courssub{Place des pronoms : la grande différence Affirmatif / Négatif}

C'est une règle mécanique stricte. La position du pronom (COD, COI, réfléchi) change radicalement selon la polarité.

\begin{propriete}[Position des pronoms avec l'impératif]
\begin{itemize}
    \item \textbf{Affirmatif :} Le pronom se \textbf{soude APRÈS} le verbe (enclitique). Un accent écrit est souvent nécessaire pour conserver la tonicité originelle. \esp{Protege} $\rightarrow$ \esp{Protégelo} ; \esp{Di} $\rightarrow$ \esp{Dímelo}.
    \item \textbf{Négatif :} Le pronom se place \textbf{AVANT} le verbe, après \esp{no} (proclitique). Pas de soudure, pas d'accent ajouté. \esp{No lo protejas} ; \esp{No me lo digas}.
\end{itemize}
\end{propriete}

\begin{center}
\begin{tabularx}{\linewidth}{|X|X|X|}
\hline
\rowcolor{poppurpleL}
\textbf{Infinitif + Pronom} & \textbf{Impératif Affirmatif (tú)} & \textbf{Impératif Négatif (tú)} \\
\hline
\esp{Protegerla} (la = la contraseña) & \esp{Protégela.} & \esp{No la protejas.} \\
\hline
\esp{Decírselo} (se + lo) & \esp{Díselo.} & \esp{No se lo digas.} \\
\hline
\esp{Enviármelo} (me + lo) & \esp{Envíamelo.} & \esp{No me lo envíes.} \\
\hline
\esp{Conectarse} (se) & \esp{Conéctate.} & \esp{No te conectes.} \\
\hline
\end{tabularx}
\end{center}

\emph{Attention aux accents :} \esp{protege} (pas d'accent) $\rightarrow$ \esp{protégela} (accent sur \emph{é} pour garder la tonicité sur l'avant-dernière syllabe). \esp{di} (accent diacritique) $\rightarrow$ \esp{dímelo} (accent sur \emph{í}). \esp{haz} $\rightarrow$ \esp{hazlo} (monosyllabe, pas d'accent sauf si deux pronoms : \esp{házmelo}).

\courssub{Les autres personnes : \textit{usted} et \textit{vosotros}}

Pour la version (thème) et la compréhension de textes officiels ou de politesse, il faut connaître les formes de \cle{usted} (politesse) et \cle{vosotros} (collectif Espagne / familiarité groupe). Elles suivent la même logique : Affirmatif = forme spécifique ; Négatif = Subjonctif.

\begin{center}
\begin{tabularx}{\linewidth}{|X|X|X|X|X|X|}
\hline
\rowcolor{popblueL}
\textbf{Verbe} & \textbf{Tú (+)} & \textbf{Tú (-)} & \textbf{Usted (+/-)} & \textbf{Vosotros (+)} & \textbf{Vosotros (-)} \\
\hline
\esp{Compartir} & \esp{Comparte} & \esp{No compartas} & \esp{Comparta / No comparta} & \esp{Compartid} & \esp{No compartáis} \\
\hline
\esp{Proteger} & \esp{Protege} & \esp{No protejas} & \esp{Proteja / No proteja} & \esp{Proteged} & \esp{No protejáis} \\
\hline
\esp{Escribir} & \esp{Escribe} & \esp{No escribas} & \esp{Escriba / No escriba} & \esp{Escribid} & \esp{No escribáis} \\
\hline
\end{tabularx}
\end{center}

\emph{Règles de formation rapides :}
\begin{itemize}
    \item \textbf{Usted} : Toujours le subjonctif présent (3\textsuperscript{e} pers. sing.) pour \textbf{les deux polarités}. \esp{Comparta} (aff.) / \esp{No comparta} (nég.).
    \item \textbf{Vosotros Affirmatif} : Remplacer \emph{-r} de l'infinitif par \emph{-d}. \esp{Compartir} $\rightarrow$ \esp{Compartid}. \esp{Proteger} $\rightarrow$ \esp{Proteged}. \esp{Ir} $\rightarrow$ \esp{Id} (irrégulier).
    \item \textbf{Vosotros Négatif} : Subjonctif présent 2\textsuperscript{e} pers. plur. \esp{No compartáis}, \esp{No protejáis}, \esp{No escribáis}.
\end{itemize}

\courssub{Comparaison Français / Espagnol : le piège de l'infinitif}

C'est l'erreur n°1 des apprenants francophones. En français, l'impersonnel ou la consigne générale utilise l'infinitif (\emph{Ne pas fumer}, \emph{Cliquer ici}, \emph{Ne pas partager}). En espagnol, \textbf{une consigne adressée à quelqu'un (même "vous" général) exige une personne verbale}.

\begin{center}
\begin{tabularx}{\linewidth}{|X|X|X|}
\hline
\rowcolor{poppinkL}
\textbf{Contexte français (Infinitif / Impersonnel)} & \textbf{Espagnol correct (Impératif / Subjonctif)} & \textbf{Erreur fréquente (Calque)} \\
\hline
Ne pas partager vos données. & \esp{No compartas tus datos.} (tú) / \esp{No compartan sus datos.} (ustedes) & \esp{*No compartir tus datos.} (\textbf{Incorrect}) \\
\hline
Cliquer sur le lien. & \esp{Haz clic en el enlace.} (tú) / \esp{Hagan clic...} (ustedes) & \esp{*Hacer clic...} (\textbf{Incorrect}) \\
\hline
Vérifier l'expéditeur. & \esp{Verifica el remitente.} / \esp{Verifiquen...} & \esp{*Verificar...} (\textbf{Incorrect}) \\
\hline
Ne pas ouvrir la pièce jointe. & \esp{No abras el archivo adjunto.} & \esp{*No abrir el archivo...} (\textbf{Incorrect}) \\
\hline
\end{tabularx}
\end{center}

\begin{methode}[Traduire une consigne française (infinitif) en espagnol (impératif)]
\begin{enumerate}
    \item \textbf{Identifier le destinataire :} \textit{tú} (familier, pair), \textit{usted} (formel, adulte), \textit{vosotros} (groupe en Espagne), \textit{ustedes} (groupe formel / AmLat).
    \item \textbf{Si la phrase est NÉGATIVE (Ne pas...)} $\rightarrow$ \textbf{No + Subjonctif présent} de la personne correspondante. \emph{Jamais d'infinitif.}
    \item \textbf{Si la phrase est AFFIRMATIVE} $\rightarrow$ \textbf{Impératif affirmatif} de la personne correspondante (3\textsuperscript{e} pers. sing. pour \textit{tú} ; subjonctif pour \textit{usted} ; infinitif -r + -d pour \textit{vosotros}).
    \item \textbf{Gérer les pronoms :} Les souder à la fin (affirmatif) ou les placer avant le verbe (négatif). Vérifier l'accent écrit.
\end{enumerate}
\end{methode}

\courssub{Exercice résolu : Thème (Français $\rightarrow$ Espagnol)}

\begin{exoresolu}[Traduire des consignes de sécurité numérique]
\textbf{Énoncé :} Traduisez en espagnol (registre \textit{tú}, conseil entre jeunes) les phrases suivantes :
\begin{enumerate}
    \item Ne partage pas ton mot de passe.
    \item Protège ton profil avec la double authentification.
    \item N'ouvre pas ce lien suspect.
    \item Dis-le à tes parents immédiatement.
    \item Ne te connecte pas au Wi-Fi public sans VPN.
\end{enumerate}

\tcblower

\textbf{Solution détaillée :}

\begin{enumerate}
    \item \esp{No compartas tu contraseña.}
    \emph{Analyse :} Négatif \textit{tú}. \esp{Compartir} $\rightarrow$ Subj. \textit{tú} \esp{compartas}. \esp{No} + subjonctif. \esp{Tu contraseña} (ton mot de passe).
    \item \esp{Protege tu perfil con la doble autenticación.}
    \emph{Analyse :} Affirmatif \textit{tú}. \esp{Proteger} $\rightarrow$ 3\textsuperscript{e} pers. sing. \esp{protege}. \esp{Doble autenticación} = double authentification (2FA).
    \item \esp{No abras ese enlace sospechoso.}
    \emph{Analyse :} Négatif \textit{tú}. \esp{Abrir} $\rightarrow$ Subj. \textit{tú} \esp{abras} (verbe en \emph{-ir} $\rightarrow$ voyelle \emph{a}). \esp{Ese enlace sospechoso} = ce lien suspect.
    \item \esp{Díselo a tus padres inmediatamente.}
    \emph{Analyse :} Affirmatif \textit{tú} + pronoms. \esp{Decir} $\rightarrow$ irrégulier \esp{di}. Pronoms \esp{se} (à eux) + \esp{lo} (le/la chose) $\rightarrow$ \esp{selo} $\rightarrow$ soudés : \esp{díselo}. Accent sur \emph{í} (monosyllabe + 2 pronoms). \emph{Immédiatement} = \esp{inmediatamente}.
    \item \esp{No te conectes a la Wi-Fi pública sin VPN.}
    \emph{Analyse :} Négatif \textit{tú} + pronom réfléchi. \esp{Conectarse} $\rightarrow$ Subj. \textit{tú} \esp{conectes}. Pronom \esp{te} AVANT le verbe : \esp{No te conectes}. \esp{Wi-Fi pública} (féminin car \esp{red} sous-entendu).
\end{enumerate}
\end{exoresolu}

\courssub{Synthèse visuelle : Organigramme de décision}

\begin{popfigure}
\begin{tikzpicture}[
    node distance=1.2cm and 1.5cm,
    decision/.style={diamond, draw=popdark, thick, fill=popgoldL, text width=2.8cm, align=center, inner sep=2pt},
    process/.style={rectangle, draw=popblue, thick, fill=popblueL, text width=3.2cm, align=center, rounded corners, inner sep=4pt},
    start/.style={ellipse, draw=popgreen, thick, fill=popgreenL, text width=2cm, align=center, inner sep=4pt},
    arrow/.style={-Stealth, thick, draw=popdark},
    yesno/.style={midway, fill=popcream, rounded corners, inner sep=2pt, font=\small}
]
\node[start] (start) {Conseil / Consigne};
\node[decision, below of=start] (pol) {Polarité ?};
\node[process, left of=pol, xshift=-2.5cm, align=center] (aff){AFFIRMATIF\\ (Fais-le !)};
\node[process, right of=pol, xshift=2.5cm, align=center] (neg){NÉGATIF\\ (Ne le fais pas !)};
\node[process, below of=aff, yshift=-0.8cm, align=center] (affrule){Règle \textit{tú}:\\ 3\textsuperscript{e} pers. sing. Présent Ind.\\(\esp{habla, protege, di, haz})\\+ Pronoms \textbf{APRÈS} (soudés)};
\node[process, below of=neg, yshift=-0.8cm, align=center] (negrule){Règle \textit{tú}:\\ \textbf{no} + Subjonctif Présent\\(\esp{no hables, no protejas, no digas})\\+ Pronoms \textbf{AVANT} (séparés)};
\node[process, below of=affrule, yshift=-0.5cm, align=center] (ust){\textit{Usted} : Subjonctif (les 2 cas)\\(\esp{Comparta / No comparta})};
\node[process, below of=negrule, yshift=-0.5cm, align=center] (vos){\textit{Vosotros} : Aff=\esp{-d} / Nég=Subj\\(\esp{Compartid / No compartáis})};

\draw[arrow] (start) -- (pol);
\draw[arrow] (pol) -- node[yesno, above left] {Oui} (aff);
\draw[arrow] (pol) -- node[yesno, above right] {Non} (neg);
\draw[arrow] (aff) -- (affrule);
\draw[arrow] (neg) -- (negrule);
\draw[arrow] (affrule) -- (ust);
\draw[arrow] (negrule) -- (vos);
\end{tikzpicture}
\legende{Organigramme de formation de l'impératif pour conseiller (niveau 1ère A).}
\end{popfigure}

\begin{savaistu}[Campagnes institutionnelles : l'impératif au service de la citoyenneté]
En Espagne et en Amérique latine, les gouvernements lancent des campagnes massives contre le cyberharcèlement (\emph{ciberacoso}) et la désinformation. Leurs slogans utilisent presque exclusivement l'impératif pour interpeller le citoyen :
\begin{itemize}
    \item \esp{« No lo ignores, denúncialo. »} (Ne l'ignore pas, dénonce-le.) — Campagne \emph{Ministerio de Educación} (Espagne).
    \item \esp{« Piensa antes de compartir. »} (Réfléchis avant de partager.) — \emph{Safer Internet Day} (International).
    \item \esp{« Bloquea, denuncia, habla. »} (Bloque, signale, parle.) — Protocole standard dans les écoles latino-américaines.
    \item \esp{« Sé legal en la red. »} (Sois légal sur le net.) — Jeu de mots \esp{sé} (impératif \esp{ser}) / \esp{sé} (je sais).
\end{itemize}
Au Cameroun, l'ANTIC (Agence Nationale des Technologies de l'Information et de la Communication) diffuse aussi des messages bilingues : \esp{« Protege tus datos, no los compartas. »} / « Protégez vos données, ne les partagez pas. »
\end{savaistu}

\courssub{Exercice d'application}

\begin{exercice}[Mettez au bon mode : Impératif Affirmatif ou Négatif \textit{tú}]
Complétez avec la forme correcte du verbe entre parenthèses. Attention aux irréguliers et aux pronoms !
\begin{enumerate}
    \item \esp{\_\_\_\_\_\_\_\_ (No / publicar) fotos de otras personas sin su consentimiento.}
    \item \esp{\_\_\_\_\_\_\_\_ (Configurar) la privacidad de tu perfil en "Solo amigos".}
    \item \esp{\_\_\_\_\_\_\_\_ (No / creer) todo lo que ves en TikTok o Instagram.}
    \item \esp{\_\_\_\_\_\_\_\_ (Decir) la verdad si ves ciberacoso.}
    \item \esp{\_\_\_\_\_\_\_\_ (No / dar) tus datos bancarios por correo electrónico.}
    \item \esp{\_\_\_\_\_\_\_\_ (Poner) una contraseña fuerte y \_\_\_\_\_\_\_\_ (cambiarla) cada tres meses.}
    \item \esp{\_\_\_\_\_\_\_\_ (No / fiarse) de mensajes que piden dinero urgente.}
    \item \esp{\_\_\_\_\_\_\_\_ (Hablar) con un profesor o tu tutor si hay problemas.}
\end{enumerate}
\end{exercice}

\begin{corrige}[Exercice d'application]
\begin{enumerate}
    \item \esp{No publiques} (Négatif \textit{tú} : \esp{no} + subj. \esp{publiques} — verbe en \emph{-ar} $\rightarrow$ \emph{-e}).
    \item \esp{Configura} (Affirmatif \textit{tú} : 3\textsuperscript{e} pers. sing. prés. ind. \esp{configura}).
    \item \esp{No creas} (Négatif \textit{tú} : \esp{no} + subj. \esp{creas}. \esp{Creer} (verbe en \emph{-er}) $\rightarrow$ subj. \textit{tú} \esp{creas} (voyelle thématique \emph{a}). L'indicatif \textit{tú} est \esp{crees} ; le subjonctif change le \emph{e} en \emph{a}).
    \item \esp{Di} (Affirmatif \textit{tú} : Irrégulier \esp{decir} $\rightarrow$ \esp{di}).
    \item \esp{No des} (Négatif \textit{tú} : \esp{no} + subj. \esp{dar} $\rightarrow$ \esp{des}. Formes du subj. présent : \esp{dé, des, dé, demos, deis, den}. La 2\textsuperscript{e} pers. \esp{des} ne porte pas d'accent, contrairement à \esp{dé} (1\textsuperscript{e}/3\textsuperscript{e} pers.)).
    \item \esp{Pon} ... \esp{cámbiala} (Affirmatif \textit{tú} : \esp{poner} $\rightarrow$ \esp{pon} irrégulier. \esp{Cambiarla} : affirmatif + pronom \esp{la} soudé $\rightarrow$ accent sur \emph{á} : \esp{cámbiala}).
    \item \esp{No te fíes} (Négatif \textit{tú} réfléchi : \esp{no} + pronom \esp{te} + subj. \esp{fiarse} $\rightarrow$ \esp{fíes} — accent sur \emph{í} pour marquer le hiatus \emph{í-e}).
    \item \esp{Habla} (Affirmatif \textit{tú} : 3\textsuperscript{e} pers. sing. \esp{hablar} $\rightarrow$ \esp{habla}).
\end{enumerate}
\end{corrige}

\emph{Fin de la Grammaire 1. La section suivante (Grammaire 2) traitera du subjonctif dans les propositions subordonnées après \esp{es importante que}, \esp{es necesario que}, structure clé pour argumenter sur les risques numériques.}

\coursec{Grammaire 2 : le subjonctif après « es importante que »}

Dans la section précédente, nous avons appris à donner des conseils directs avec l'\cle{impératif} : \esp{No compartas tu contraseña.}, \esp{Protege tu privacidad.} Mais il existe une autre manière, plus nuancée et très fréquente à l'écrit, de conseiller ou de formuler une recommandation : les \cle{expressions impersonnelles} comme \esp{es importante que}, \esp{es necesario que}, \esp{es mejor que}. Ces structures exigent un mode particulier : le \cle{subjonctif}. C'est la grammaire reine du commentaire de texte et de l'expression écrite au Bac : impossible de « proposer des comportements responsables » sans elle.

\courssub{Observation : à la découverte de la structure}

Lisons d'abord ce court extrait d'une campagne de sensibilisation numérique destinée aux lycéens.

\begin{textebox}[Campaña escolar : « Navega con cabeza »]
Queridos alumnos: internet es una herramienta maravillosa, pero exige responsabilidad. Es importante que hables con un adulto si recibes mensajes extraños. Es necesario que protejas tu privacidad y que no compartas tu contraseña con nadie. Es mejor que limites el tiempo de pantalla y que verifiques las fuentes antes de creer una noticia. Conviene que aprendas a descargar solo aplicaciones seguras. Recuerda: es bueno que descanses sin el móvil por la noche. ¡Tu seguridad digital empieza hoy!
\end{textebox}

\textbf{Glossaire :} \esp{herramienta} = outil ; \esp{extraño} = étrange, bizarre ; \esp{fuentes} = sources ; \esp{creer una noticia} = croire une information ; \esp{descansar} = se reposer.

\textbf{Traduction :} « Chers élèves : internet est un outil merveilleux, mais il exige de la responsabilité. Il est important que tu parles à un adulte si tu reçois des messages étranges. Il est nécessaire que tu protèges ta vie privée et que tu ne partages ton mot de passe avec personne. Il vaut mieux que tu limites le temps d'écran et que tu vérifies les sources avant de croire une information. Il convient que tu apprennes à ne télécharger que des applications sûres. Souviens-toi : il est bon que tu te reposes sans le portable le soir. Ta sécurité numérique commence aujourd'hui ! »

\textbf{Questions de compréhension :}
\begin{enumerate}
\item Quels sont les quatre conseils donnés aux élèves ?
\item Relevez toutes les expressions qui commencent par \esp{es...} ou \esp{conviene...}. Qu'ont-elles en commun ?
\item Observez les verbes \esp{hables}, \esp{protejas}, \esp{compartas}, \esp{limites}, \esp{verifiques}, \esp{aprendas}, \esp{descanses} : à quel mode appartiennent-ils ?
\end{enumerate}

\textbf{Ce que nous observons :} chaque recommandation suit le même schéma :

\begin{center}
\textbf{expression impersonnelle} + \esp{que} + \textbf{verbe au subjonctif présent}
\end{center}

\esp{Es importante que hables con un adulto.} = Il est important que tu parles à un adulte.\\
\esp{Es necesario que protejas tu privacidad.} = Il est nécessaire que tu protèges ta vie privée.

\courssub{La règle : expressions impersonnelles + subjonctif}

\begin{propriete}[Expressions impersonnelles de conseil + subjonctif]
Après les tournures impersonnelles qui expriment la \cle{nécessité}, l'\cle{importance}, le \cle{conseil} ou le \cle{jugement de valeur}, le verbe de la subordonnée introduite par \esp{que} se met au \cle{subjonctif présent} :
\begin{itemize}
\item \esp{es importante que} = il est important que
\item \esp{es necesario que} = il est nécessaire que
\item \esp{es bueno que} = il est bon que
\item \esp{es mejor que} = il vaut mieux que
\item \esp{conviene que} = il convient que
\end{itemize}
Le mot \esp{que} est \textbf{obligatoire} et ne se supprime jamais.
\end{propriete}

\begin{exemplebox}[Exemples traduits]
\esp{Es necesario que cambies tu contraseña.} = Il est nécessaire que tu changes ton mot de passe.\\
\esp{Es mejor que no uses el móvil por la noche.} = Il vaut mieux que tu n'utilises pas le portable le soir.\\
\esp{Es importante que verifiques las fuentes.} = Il est important que tu vérifies les sources.\\
\esp{Es bueno que descanses sin pantallas.} = Il est bon que tu te reposes sans écrans.\\
\esp{Conviene que los jóvenes aprendan a usar internet con responsabilidad.} = Il convient que les jeunes apprennent à utiliser internet de manière responsable.
\end{exemplebox}

\begin{attention}[Deux erreurs à bannir]
\begin{enumerate}
\item \textbf{Ne jamais oublier} \esp{que} : *\esp{es importante protejas tu privacidad} est \textbf{faux}. Il faut : \esp{es importante que protejas tu privacidad}.
\item \textbf{Ne jamais mettre l'indicatif} après ces expressions : *\esp{es importante que proteges} est fautif ; la bonne forme est \esp{protejas} (subjonctif). L'indicatif exprime un fait ; ici, nous exprimons une nécessité, pas une réalité constatée.
\end{enumerate}
\end{attention}

\courssub{Formation du subjonctif présent}

\begin{propriete}[Formation régulière]
On part du \cle{radical de la 1\iere{} personne du singulier du présent de l'indicatif} (la forme en \esp{yo}, sans le \esp{-o} final), puis on ajoute les \cle{terminaisons inversées} :
\begin{itemize}
\item verbes en \esp{-ar} : \textbf{-e, -es, -e, -emos, -éis, -en}
\item verbes en \esp{-er} et \esp{-ir} : \textbf{-a, -as, -a, -amos, -áis, -an}
\end{itemize}
Astuce : le \esp{a} de l'indicatif devient \esp{e}, et le \esp{e/i} devient \esp{a}. On « croise » les voyelles.
\end{propriete}

\begin{center}
\begin{tabular}{|l|l|l|l|}
\hline
\rowcolor{popblueL} Personne & \esp{hablar} (-ar) & \esp{comer} (-er) & \esp{escribir} (-ir) \\
\hline
yo & hable & coma & escriba \\
\hline
tú & hables & comas & escribas \\
\hline
él / ella / usted & hable & coma & escriba \\
\hline
nosotros & hablemos & comamos & escribamos \\
\hline
vosotros & habléis & comáis & escribáis \\
\hline
ellos / ustedes & hablen & coman & escriban \\
\hline
\end{tabular}
\end{center}

\textbf{Attention aux radicaux irréguliers :} comme on part de la forme \esp{yo}, les verbes à radical changeant ou à 1\iere{} personne irrégulière gardent cette irrégularité : \esp{tener} → \esp{tengo} → \esp{tenga} ; \esp{hacer} → \esp{hago} → \esp{haga} ; \esp{poder} → \esp{puedo} → \esp{pueda} ; \esp{pensar} → \esp{pienso} → \esp{piense}.

\begin{popfigure}
\begin{center}
\begin{tikzpicture}[
  every node/.style={font=\small},
  head/.style={fill=popblue, text=white, rounded corners=2pt, minimum width=1.7cm, align=center},
  arv/.style={fill=poporangeL, rounded corners=2pt, minimum width=1.7cm, align=center},
  erv/.style={fill=popgreenL, rounded corners=2pt, minimum width=1.7cm, align=center},
  irv/.style={fill=popgoldL, rounded corners=2pt, minimum width=1.7cm, align=center},
  row/.style={font=\small\bfseries, text=popdark}
]
\node[head] (h1) at (0,0) {yo};
\node[head, right=0.15cm of h1] (h2) {tú};
\node[head, right=0.15cm of h2] (h3) {él};
\node[head, right=0.15cm of h3] (h4) {nosotros};
\node[head, right=0.15cm of h4] (h5) {vosotros};
\node[head, right=0.15cm of h5] (h6) {ellos};
\node[arv, below=0.12cm of h1] (a1) {hable};
\node[arv, below=0.12cm of h2] (a2) {hables};
\node[arv, below=0.12cm of h3] (a3) {hable};
\node[arv, below=0.12cm of h4] (a4) {hablemos};
\node[arv, below=0.12cm of h5] (a5) {habléis};
\node[arv, below=0.12cm of h6] (a6) {hablen};
\node[erv, below=0.12cm of a1] (e1) {coma};
\node[erv, below=0.12cm of a2] (e2) {comas};
\node[erv, below=0.12cm of a3] (e3) {coma};
\node[erv, below=0.12cm of a4] (e4) {comamos};
\node[erv, below=0.12cm of a5] (e5) {comáis};
\node[erv, below=0.12cm of a6] (e6) {coman};
\node[irv, below=0.12cm of e1] (i1) {escriba};
\node[irv, below=0.12cm of e2] (i2) {escribas};
\node[irv, below=0.12cm of e3] (i3) {escriba};
\node[irv, below=0.12cm of e4] (i4) {escribamos};
\node[irv, below=0.12cm of e5] (i5) {escribáis};
\node[irv, below=0.12cm of e6] (i6) {escriban};
\node[row, left=0.25cm of a1] {-ar};
\node[row, left=0.25cm of e1] {-er};
\node[row, left=0.25cm of i1] {-ir};
\end{tikzpicture}
\end{center}
\legende{Terminaisons du subjonctif présent : les voyelles se croisent (a $\rightarrow$ e ; e/i $\rightarrow$ a).}
\end{popfigure}

\begin{propriete}[Les neuf irréguliers à connaître par cœur]
\begin{center}
\begin{tabular}{|l|l|l|}
\hline
\rowcolor{popblueL} Infinitif & Subjonctif (yo / tú / él...) & Français \\
\hline
ser & sea, seas, sea, seamos, seáis, sean & être \\
\hline
ir & vaya, vayas, vaya, vayamos, vayáis, vayan & aller \\
\hline
haber & haya, hayas, haya, hayamos, hayáis, hayan & avoir (auxiliaire) \\
\hline
saber & sepa, sepas, sepa, sepamos, sepáis, sepan & savoir \\
\hline
estar & esté, estés, esté, estemos, estéis, estén & être (état) \\
\hline
dar & dé, des, dé, demos, deis, den & donner \\
\hline
tener & tenga, tengas, tenga... & avoir \\
\hline
hacer & haga, hagas, haga... & faire \\
\hline
poder & pueda, puedas, pueda... & pouvoir \\
\hline
\end{tabular}
\end{center}
Notez l'accent sur \esp{dé} (pour le distinguer de la préposition \esp{de}) et sur \esp{esté}. En revanche, \esp{deis} (vosotros) ne prend pas d'accent.
\end{propriete}

\courssub{Comparaison avec le français : un point d'appui précieux}

Bonne nouvelle pour les francophones : le français fonctionne \textbf{exactement de la même manière}. Après « il est important que », « il faut que », « il vaut mieux que », le français exige lui aussi le subjonctif.

\begin{center}
\begin{tabularx}{\linewidth}{|X|X|X|}
\hline
\rowcolor{popblueL} Français & Español & Remarque \\
\hline
Il est important que tu \emph{parles}. & Es importante que \emph{hables}. & Structure parfaitement parallèle. \\
\hline
Il faut que tu \emph{protèges} ta vie privée. & Es necesario que \emph{protejas} tu privacidad. & « Il faut que » = \esp{es necesario que} ; « il faut » + infinitif = \esp{hay que} + infinitif. \\
\hline
Il vaut mieux que tu ne \emph{partages} pas. & Es mejor que no \emph{compartas}. & En espagnol, \esp{no} suffit (pas de « ne » explétif). \\
\hline
Il convient que les jeunes \emph{apprennent}. & Conviene que los jóvenes \emph{aprendan}. & Tournure soutenue, fréquente au Bac. \\
\hline
\end{tabularx}
\end{center}

\begin{attention}[Faux amis et calques]
\begin{itemize}
\item \esp{es importante que} ne se traduit jamais par « c'est important de » quand il y a un sujet précis : « Il est important que tu changes » = \esp{es importante que cambies}, et non *\esp{es importante de cambiar tú}.
\item Ne traduisez pas « il faut que je fasse » par *\esp{es necesario que yo hago} : le subjonctif est obligatoire → \esp{es necesario que (yo) haga}.
\item Le pronom sujet (\esp{tú}, \esp{yo}) est le plus souvent omis en espagnol : \esp{es importante que protejas tu privacidad}.
\end{itemize}
\end{attention}

\courssub{Subjonctif ou infinitif ? Le choix décisif}

\begin{propriete}[Sujet précis → subjonctif ; sujet général → infinitif]
\begin{itemize}
\item Si la recommandation vise une \cle{personne précise} (sujet différent de celui de l'expression impersonnelle), on utilise \esp{que} + \textbf{subjonctif} :\\
\esp{Es importante que tú protejas tu privacidad.} = Il est important que \emph{toi}, tu protèges ta vie privée.
\item Si le conseil est \cle{général} (il vaut pour tout le monde, sans sujet déterminé), l'espagnol préfère l'\textbf{infinitif}, sans \esp{que} :\\
\esp{Es importante proteger la privacidad.} = Il est important de protéger la vie privée.
\end{itemize}
\end{propriete}

\begin{exemplebox}[Les deux structures en miroir]
\esp{Es importante que descargues solo aplicaciones seguras.} = Il est important que tu ne télécharges que des applications sûres. (conseil adressé à toi)\\
\esp{Es importante descargar aplicaciones seguras.} = Il est important de télécharger des applications sûres. (principe général)\\
\esp{Es necesario que limites el tiempo de pantalla.} = Il est nécessaire que tu limites ton temps d'écran.\\
\esp{Es necesario limitar el tiempo de pantalla.} = Il est nécessaire de limiter le temps d'écran.
\end{exemplebox}

\courssub{Méthode : traduire la structure en version}

\begin{methode}[Traduire « es importante que + subjonctif » en français]
\begin{enumerate}
\item \textbf{Repérer} la structure impersonnelle : \esp{es importante / necesario / mejor / bueno que}, \esp{conviene que}.
\item \textbf{Identifier le mode} : le verbe qui suit \esp{que} est au subjonctif → il exprime un conseil, une nécessité, pas un fait.
\item \textbf{Identifier la personne} du verbe subjonctif (souvent \esp{tú} sous-entendu : \esp{hables} = « que tu parles »).
\item \textbf{Traduire} par « il est important que » + subjonctif français : \esp{es mejor que no compartas tu contraseña} = « il vaut mieux que tu ne partages pas ton mot de passe ».
\item \textbf{Vérifier} : si la phrase espagnole a un infinitif (\esp{es importante proteger...}), traduire par « il est important de + infinitif ».
\end{enumerate}
\end{methode}

\begin{exoresolu}[Version guidée]
Énoncé : traduisez en français : \esp{Conviene que los alumnos sepan reconocer las noticias falsas en internet.}
\tcblower
Solution détaillée :
\begin{enumerate}
\item Structure repérée : \esp{conviene que} = « il convient que » (expression impersonnelle de conseil).
\item Verbe \esp{sepan} : subjonctif présent de \esp{saber} (irrégulier), 3\iere{} personne du pluriel → sujet : \esp{los alumnos}.
\item Traduction : « Il convient que les élèves sachent reconnaître les fausses informations sur internet. »
\item Piège évité : on n'a pas traduit \esp{sepan} par un indicatif (« savent »), car le français exige le subjonctif après « il convient que ».
\end{enumerate}
\end{exoresolu}

\begin{savaistu}[« Conviene que » : la tournure du Bac]
\esp{Conviene que} (« il convient que ») est une expression soutenue, très fréquente à l'écrit — articles de presse, discours officiels, textes argumentatifs — et très appréciée dans les copies du Bac. Exemple modèle pour vos productions : \esp{Conviene que los jóvenes aprendan a usar internet con responsabilidad.} = « Il convient que les jeunes apprennent à utiliser internet de manière responsable. » Le verbe \esp{convenir} est irrégulier (il se conjugue comme \esp{venir}), mais à la forme impersonnelle on n'utilise que \esp{conviene}. À noter : \esp{convenir} est un faux ami partiel : il signifie « convenir, être approprié », et non « convenir d'un rendez-vous » au sens français de « se mettre d'accord ».
\end{savaistu}

\courssub{Entraînement : thème et version}

\begin{exercice}[Thème : du français vers l'espagnol]
Traduisez en espagnol en utilisant une expression impersonnelle + subjonctif :
\begin{enumerate}
\item Il est important que tu parles à un adulte.
\item Il est nécessaire que tu changes ton mot de passe.
\item Il vaut mieux que tu n'utilises pas le portable le soir.
\item Il est bon que les jeunes limitent leur temps d'écran.
\item Il convient que nous vérifiions les sources avant de partager une information.
\end{enumerate}
\end{exercice}

\begin{corrige}[Thème]
\begin{enumerate}
\item \esp{Es importante que hables con un adulto.}
\item \esp{Es necesario que cambies tu contraseña.}
\item \esp{Es mejor que no uses el móvil por la noche.}
\item \esp{Es bueno que los jóvenes limiten su tiempo de pantalla.}
\item \esp{Conviene que verifiquemos las fuentes antes de compartir una noticia.}
\end{enumerate}
Points de vigilance : subjonctif à chaque verbe (\esp{hables}, \esp{cambies}, \esp{uses}, \esp{limiten}, \esp{verifiquemos}) ; \esp{que} obligatoire ; négation simple avec \esp{no} à la phrase 3 ; à la phrase 5, \esp{verifiquemos} prend un \esp{qu} pour conserver le son « k » devant \esp{e} (comme \esp{buscar} → \esp{busque}).
\end{corrige}

\begin{exercice}[Version : de l'espagnol vers le français]
Traduisez en français :
\begin{enumerate}
\item Es importante que protejas tu privacidad en las redes sociales.
\item Es necesario que descargues solo aplicaciones seguras.
\item Es mejor que no compartas tus fotos personales con desconocidos.
\item Es bueno que haya normas contra el ciberacoso en los colegios.
\item Conviene que estés atento cuando navegas por internet.
\end{enumerate}
\end{exercice}

\begin{corrige}[Version]
\begin{enumerate}
\item Il est important que tu protèges ta vie privée sur les réseaux sociaux.
\item Il est nécessaire que tu ne télécharges que des applications sûres.
\item Il vaut mieux que tu ne partages pas tes photos personnelles avec des inconnus.
\item Il est bon qu'il y ait des règles contre le cyberharcèlement dans les collèges. (\esp{haya} = subjonctif de \esp{haber})
\item Il convient que tu sois attentif quand tu navigues sur internet. (\esp{estés} = subjonctif de \esp{estar})
\end{enumerate}
\end{corrige}

\begin{aretenir}[L'essentiel de la section]
\begin{itemize}
\item Après \esp{es importante / necesario / bueno / mejor que} et \esp{conviene que}, le verbe est au \cle{subjonctif présent}.
\item Formation : radical de \esp{yo} au présent + terminaisons inversées (-ar → \textbf{e} ; -er/-ir → \textbf{a}).
\item Irréguliers : \esp{sea, vaya, haya, sepa, esté, dé, tenga, haga, pueda}.
\item Sujet précis → \esp{que} + subjonctif ; conseil général → infinitif sans \esp{que}.
\item Le français fonctionne pareil : « il est important que tu \emph{parles} » = \esp{es importante que \emph{hables}}. C'est votre meilleur allié !
\end{itemize}
\end{aretenir}

\coursec{Méthode du commentaire et commentaire modèle}

Dans la section précédente, nous avons appris à employer le subjonctif après \esp{es importante que} pour formuler des conseils : \esp{Es importante que protejas tu privacidad.} Ces structures, ainsi que les impératifs étudiés avant elles, ne sont pas seulement des outils de grammaire : ce sont aussi des \cle{procédés d'écriture} que l'on retrouve dans le texte support \esp{Conectados, ¿pero protegidos?}. Nous allons maintenant apprendre à les \emph{analyser} dans un commentaire de texte, exercice central de l'évaluation en Première A. Le commentaire consiste à lire un texte, à dégager sa problématique, puis à montrer \emph{comment} l'auteur construit son message, en s'appuyant sur des citations précises.

\courssub{Qu'est-ce qu'un commentaire de texte ?}

\begin{definition}[Le commentaire de texte]
Le \cle{commentaire de texte} (\esp{comentario de texto}) est un exercice d'expression écrite qui consiste à analyser un texte de manière organisée : on présente le texte et son contexte (introduction), on étudie ses idées et ses procédés d'expression en deux ou trois axes (développement), puis on fait le bilan et on ouvre sur une question plus large (conclusion). En espagnol LV2, le commentaire se rédige en espagnol, avec des citations courtes du texte entre guillemets.
\end{definition}

La différence essentielle avec un simple résumé : le résumé dit \emph{ce que} le texte raconte ; le commentaire montre \emph{comment} le texte le dit et \emph{pourquoi} c'est efficace. On ne se contente donc jamais de paraphraser : on cite, puis on analyse la citation (lexique, temps verbaux, impératifs, ton, exemples).

\begin{methode}[Rédiger l'introduction du commentaire]
L'introduction se construit en quatre étapes, en quatre ou cinq phrases :
\begin{enumerate}
\item \textbf{Phrase de mise en contexte} : situer le thème général. Exemple : \esp{Hoy en día, las tecnologías digitales forman parte de la vida cotidiana de los jóvenes.} (Aujourd'hui, les technologies numériques font partie de la vie quotidienne des jeunes.)
\item \textbf{Présentation du texte} : indiquer le titre, la nature du texte (article de presse, récit, discours), sa source et son thème. Exemple : \esp{El texto « Conectados, ¿pero protegidos? » es un artículo de prensa que trata de los riesgos de internet para los adolescentes.} (Le texte « Connectés, mais protégés ? » est un article de presse qui traite des risques d'internet pour les adolescents.)
\item \textbf{Annonce de la problématique} : formuler la question directrice, sous forme interrogative, avec les signes \esp{¿} et \esp{?}. Exemple : \esp{¿Cómo denuncia el texto los riesgos del mundo digital y qué soluciones propone?} (Comment le texte dénonce-t-il les risques du monde numérique et quelles solutions propose-t-il ?)
\item \textbf{Annonce du plan} : présenter brièvement les axes. Exemple : \esp{Estudiaremos primero la presentación de los riesgos y después los consejos y soluciones propuestos por el autor.} (Nous étudierons d'abord la présentation des risques, puis les conseils et solutions proposés par l'auteur.)
\end{enumerate}
\end{methode}

\begin{propriete}[La structure complète du commentaire]
Le commentaire suit toujours la même architecture :
\begin{itemize}
\item \textbf{Introduction} : contexte + présentation du texte + problématique + annonce du plan.
\item \textbf{Développement} : deux ou trois axes. Chaque axe contient une idée principale, des \cle{citations} du texte en espagnol et l'analyse de ces citations (procédés relevés et leur effet).
\item \textbf{Conclusion} : bilan de l'analyse (réponse à la problématique) + \cle{ouverture} (une question ou un élargissement vers un thème voisin, par exemple les réseaux sociaux au Cameroun).
\end{itemize}
Les axes sont reliés par des \cle{connecteurs logiques} qui donnent de la clarté à la copie.
\end{propriete}

\begin{popfigure}
\begin{center}
\begin{tikzpicture}[
node distance=0.55cm,
etape/.style={rectangle, rounded corners=4pt, draw=popblue, fill=popblueL, align=center, text width=6.4cm, minimum height=0.9cm, font=\small},
axe/.style={rectangle, rounded corners=4pt, draw=poporange, fill=poporangeL, align=center, text width=6.4cm, minimum height=0.9cm, font=\small},
fleche/.style={-{Stealth[length=2.5mm]}, thick, popdark}
]
\node[etape, align=center] (intro){\textbf{Introducción}\\ contexto + texto + problemática + plan};
\node[etape, below=of intro, draw=poppurple, fill=poppinkL, align=center] (prob){\textbf{Problemática}\\ ¿Cómo denuncia el texto los riesgos\\ y qué soluciones propone?};
\node[axe, below=of prob, align=center] (axe1){\textbf{Eje 1 : los riesgos}\\ ciberacoso, adicción, privacidad\\ (campos léxicos, ejemplo de Andrés)};
\node[axe, below=of axe1, align=center] (axe2){\textbf{Eje 2 : las soluciones}\\ imperativos y subjuntivo\\ (visada didáctica : aconsejar)};
\node[etape, below=of axe2, draw=popgreen, fill=popgreenL, align=center] (concl){\textbf{Conclusión}\\ balance + apertura};
\draw[fleche] (intro) -- (prob);
\draw[fleche] (prob) -- (axe1);
\draw[fleche] (axe1) -- (axe2);
\draw[fleche] (axe2) -- (concl);
\end{tikzpicture}
\end{center}
\legende{Schéma de la structure du commentaire de texte : de l'introduction à la conclusion, chaque partie enchaîne logiquement avec la suivante.}
\end{popfigure}

\courssub{Les connecteurs logiques : la charpente de la copie}

Un commentaire sans connecteurs ressemble à une liste d'idées désordonnées. Les connecteurs (\esp{conectores lógicos}) permettent d'annoncer les étapes, d'opposer, de conclure. Voici les plus utiles au niveau Première :

\begin{center}
\begin{tabularx}{\linewidth}{|X|X|X|}
\hline
\rowcolor{popblueL} \textbf{Español} & \textbf{Français} & \textbf{Emploi} \\
\hline
\esp{en primer lugar} & en premier lieu, d'abord & annoncer le premier axe \\
\hline
\esp{además} & de plus, en outre & ajouter une idée \\
\hline
\esp{sin embargo} & cependant, pourtant & opposer deux idées \\
\hline
\esp{por lo tanto} & par conséquent, donc & exprimer une conséquence \\
\hline
\esp{por último} & enfin, en dernier lieu & annoncer le dernier axe \\
\hline
\esp{en conclusión} & en conclusion & ouvrir la conclusion \\
\hline
\end{tabularx}
\end{center}

\begin{exemplebox}[Connecteurs en contexte]
\esp{En primer lugar, el autor presenta los riesgos de internet.} « En premier lieu, l'auteur présente les risques d'internet. »

\esp{Además, utiliza el ejemplo de Andrés para ilustrar el ciberacoso.} « De plus, il utilise l'exemple d'Andrés pour illustrer le cyberharcèlement. »

\esp{El texto describe los peligros; sin embargo, no se limita a criticar.} « Le texte décrit les dangers ; cependant, il ne se limite pas à critiquer. »

\esp{Los jóvenes están muy expuestos; por lo tanto, el autor propone consejos concretos.} « Les jeunes sont très exposés ; par conséquent, l'auteur propose des conseils concrets. »

\esp{En conclusión, el texto alerta y educa al mismo tiempo.} « En conclusion, le texte alerte et éduque en même temps. »
\end{exemplebox}

\begin{attention}[L'erreur classique : paraphraser sans analyser]
Au commentaire, il ne suffit pas de répéter le texte avec d'autres mots. Il faut \textbf{citer en espagnol} (entre guillemets, citation courte) puis \textbf{analyser le procédé} : champ lexical, temps verbal, impératif, subjonctif, exemple vécu, chiffre, opposition.

Mauvais : \esp{El texto dice que internet es peligroso para los jóvenes.} (simple paraphrase, aucune analyse).

Bon : \esp{El campo léxico del peligro (« ciberacoso », « adicción », « robo de datos ») crea un tono de alerta.} (citation + procédé + effet).

Autre piège : ne pas traduire les citations en français dans la copie rédigée en espagnol ; la citation reste en espagnol, l'analyse aussi. Attention enfin à l'orthographe des mots techniques : \esp{ciberacoso} (avec \emph{c} et un seul \emph{r}), \esp{privacidad} (sans accent), \esp{adicción} (avec \emph{cc} et accent sur le \emph{ó}).
\end{attention}

\courssub{Construire les axes d'analyse}

Reprenons la problématique retenue pour notre texte support : \esp{¿Cómo denuncia el texto los riesgos del mundo digital y qué soluciones propone?} Elle appelle naturellement deux axes.

\textbf{Axe 1 : la présentation des risques.} On relève d'abord les \cle{champs lexicaux} : le champ lexical du danger (\esp{ciberacoso, adicción, robo de datos, privacidad, amenaza}) et celui de la technologie (\esp{móvil, aplicación, contraseña, descargar, pantalla}). On relève ensuite les \cle{procédés} : l'exemple vécu du jeune Andrés, qui rend le danger concret et proche du lecteur ; les formulations qui généralisent (\esp{muchos jóvenes, cada día}) ; le ton d'alerte créé par l'accumulation des termes négatifs.

\textbf{Axe 2 : les conseils et les solutions.} On relève les \cle{impératifs} (\esp{no compartas, protege, limita, habla}) et les structures au \cle{subjonctif} (\esp{es importante que controles, es necesario que hables}). On montre alors la \cle{visée didactique} du texte (\esp{visada didáctica}) : il ne veut pas seulement informer, il veut \emph{former} le lecteur et transformer son comportement. C'est la fonction argumentative du texte.

\begin{propriete}[Citer et analyser : la formule en trois temps]
Pour chaque idée du développement, suis la séquence :
\begin{enumerate}
\item \textbf{Idée} : \esp{El autor denuncia el ciberacoso.} (L'auteur dénonce le cyberharcèlement.)
\item \textbf{Citation} : \esp{... a través del caso de Andrés, que recibió « mensajes insultantes » durante semanas.} (... à travers le cas d'Andrés, qui a reçu des « messages insultants » pendant des semaines.)
\item \textbf{Analyse} : \esp{El ejemplo concreto y el adjetivo « insultantes » producen un efecto de realidad y de alarma en el lector.} (L'exemple concret et l'adjectif « insultants » produisent un effet de réalité et d'alarme chez le lecteur.)
\end{enumerate}
Cette séquence idée -- citation -- analyse est l'unité de base de tout bon paragraphe de commentaire.
\end{propriete}

\courssub{Commentaire modèle rédigé}

Lisons maintenant un commentaire modèle complet, tel qu'on pourrait le rédiger en copie. Pendant la lecture collective, souligne les connecteurs et repère la séquence idée -- citation -- analyse dans chaque paragraphe.

\begin{textebox}[Comentario modelo : « Conectados, ¿pero protegidos? »]
Hoy en día, las tecnologías digitales forman parte de la vida cotidiana de los jóvenes, en España, en América Latina y también en Camerún. El texto « Conectados, ¿pero protegidos? » es un artículo de prensa que trata de los riesgos de internet para los adolescentes. Podemos preguntarnos : ¿cómo denuncia el texto los riesgos del mundo digital y qué soluciones propone? Estudiaremos primero la presentación de los riesgos y después los consejos del autor.

En primer lugar, el texto presenta los riesgos del mundo digital a través del caso de Andrés, un adolescente que sufrió ciberacoso. El campo léxico del peligro (« ciberacoso », « adicción », « robo de datos ») crea un tono de alerta, mientras que el ejemplo vivido produce un efecto de realidad : el lector comprende que el peligro no es abstracto. Además, las expresiones que generalizan, como « muchos jóvenes » y « cada día », muestran que el problema concierne a toda una generación.

Sin embargo, el texto no se limita a criticar : también propone soluciones. Los imperativos (« no compartas », « protege », « limita ») muestran la voluntad de aconsejar directamente al lector, y las estructuras con subjuntivo, como « es importante que controles tu tiempo », refuerzan esta visada didáctica. Por lo tanto, el autor no quiere asustar, sino educar : transformar los comportamientos de los jóvenes internautas.

En conclusión, el texto denuncia los riesgos del mundo digital con un tono de alerta y propone comportamientos responsables gracias a los imperativos y al subjuntivo. Es un texto útil que combina información y consejo. Podemos abrir la reflexión : ¿cómo proteger también a los jóvenes cameruneses de las noticias falsas en las redes sociales?
\end{textebox}

\textbf{Glossaire} : \esp{la prensa} = la presse ; \esp{el caso} = le cas ; \esp{vivido} = vécu ; \esp{el tono de alerta} = le ton d'alerte ; \esp{concierne a} = concerne ; \esp{la visada didáctica} = la visée didactique ; \esp{asustar} = effrayer ; \esp{sino} = mais (après une négation) ; \esp{las noticias falsas} = les fausses informations (fake news) ; \esp{cameruneses} = camerounais.

\begin{aretenir}[Les qualités du commentaire modèle]
\begin{itemize}
\item Une introduction complète en quatre étapes (contexte, texte, problématique, plan).
\item Deux axes annoncés par \esp{en primer lugar} et \esp{sin embargo}, reliés logiquement.
\item Des citations courtes entre guillemets, toujours suivies d'une analyse de procédé.
\item Une conclusion avec bilan (\esp{en conclusión}) et ouverture vers un thème voisin.
\item Une langue simple mais correcte : phrases claires, connecteurs variés, aucune paraphrase vide.
\end{itemize}
\end{aretenir}

\begin{exoresolu}[Analyser une phrase du commentaire modèle]
Énoncé. Dans la phrase \esp{« El campo léxico del peligro (ciberacoso, adicción, robo de datos) crea un tono de alerta, mientras que los imperativos (no compartas, protege) muestran la voluntad de aconsejar al lector »}, identifie : a) les citations ; b) les procédés analysés ; c) les effets annoncés.
\tcblower
Solution détaillée. a) Les citations sont les mots entre parenthèses, repris du texte support : \esp{ciberacoso, adicción, robo de datos} et les impératifs \esp{no compartas, protege}. b) Les procédés analysés sont le \cle{champ lexical} du danger (procédé lexical) et l'emploi de l'\cle{impératif} (procédé grammatical). c) Les effets annoncés sont le \esp{tono de alerta} (ton d'alerte) pour le champ lexical et la \esp{voluntad de aconsejar} (volonté de conseiller) pour les impératifs. La phrase respecte donc parfaitement la séquence idée -- citation -- analyse : c'est un modèle de paragraphe d'analyse.
\end{exoresolu}

\courssub{Entraînement : construire problématique et plan sur un nouveau texte}

Pour vérifier que la méthode est acquise, appliquons-la à un texte nouveau. Lis le texte suivant, puis construis seul la problématique et le plan en deux axes.

\begin{textebox}[Texto de entrenamiento : « Noticias falsas, ¿verdades peligrosas? »]
En los últimos años, las noticias falsas se han convertido en un problema mundial. En las redes sociales, una mentira puede viajar más rápido que la verdad : en pocos minutos, miles de personas comparten una información sin verificarla.

El caso de María, una estudiante de Malabo, ilustra este fenómeno. Un día, la joven leyó en su móvil que su escuela cerraría durante un mes. Sin comprobar la fuente, compartió la noticia con sus compañeros. Al día siguiente, toda la ciudad hablaba del cierre, pero era falso. « Me sentí avergonzada », confiesa María.

Las consecuencias de las noticias falsas son graves : crean miedo, destruyen reputaciones y confunden a la opinión pública. Además, los jóvenes son especialmente vulnerables porque pasan muchas horas conectados.

¿Qué podemos hacer? Verifica siempre la fuente antes de compartir. No creas todo lo que lees en internet. Consulta varios medios de comunicación serios. Es importante que los jóvenes aprendan a pensar con espíritu crítico, y es necesario que las familias y los profesores hablen de este tema. La lucha contra las noticias falsas empieza con cada uno de nosotros.
\end{textebox}

\textbf{Glossaire} : \esp{la mentira} = le mensonge ; \esp{verificar / comprobar} = vérifier ; \esp{la fuente} = la source ; \esp{avergonzada} = honteuse ; \esp{el cierre} = la fermeture ; \esp{el miedo} = la peur ; \esp{confunden} = ils trompent, ils induisent en erreur ; \esp{el espíritu crítico} = l'esprit critique ; \esp{la lucha} = la lutte.

\begin{exercice}[Problématique et plan sur « Noticias falsas »]
\begin{enumerate}
\item Présente le texte en deux phrases (nature, titre, thème).
\item Formule une problématique interrogative adaptée à ce texte.
\item Construis un plan en deux axes : pour chaque axe, donne le titre et deux procédés à relever (champ lexical, exemple vécu, impératifs, subjonctif).
\item Rédige la première phrase de chaque axe avec un connecteur approprié.
\end{enumerate}
\end{exercice}

\begin{corrige}[Exercice : problématique et plan]
1. Présentation : \esp{« Noticias falsas, ¿verdades peligrosas? » es un artículo de prensa que trata del problema de la desinformación en las redes sociales.} (C'est un article de presse qui traite du problème de la désinformation sur les réseaux sociaux.)

2. Problématique possible : \esp{¿Cómo denuncia el texto el peligro de las noticias falsas y qué comportamientos propone para combatirlas?} (Comment le texte dénonce-t-il le danger des fake news et quels comportements propose-t-il pour les combattre ?) Remarque : \esp{combatirlas} = les combattre ; le pronom \esp{las} reprend \esp{las noticias falsas} et se colle après l'infinitif.

3. Plan en deux axes. Axe 1 : \esp{la denuncia del fenómeno} — procédés : l'exemple vécu de María (\esp{el caso de María}) et le champ lexical du mensonge et de la peur (\esp{mentira, falso, miedo, avergonzada}). Axe 2 : \esp{las soluciones propuestas} — procédés : les impératifs (\esp{verifica, no creas, consulta}) et les structures au subjonctif (\esp{es importante que los jóvenes aprendan, es necesario que las familias hablen}).

4. Premières phrases : \esp{En primer lugar, el autor denuncia el fenómeno de las noticias falsas a través del caso de María.} ; \esp{Sin embargo, el texto no se queda en la crítica : propone soluciones concretas para el lector.} (Cependant, le texte ne s'arrête pas à la critique : il propose des solutions concrètes pour le lecteur.)
\end{corrige}

\begin{experience}[Atelier d'écriture en binômes]
Par groupes de deux, rédigez l'introduction complète du commentaire sur \esp{Noticias falsas} (quatre phrases : contexte, présentation, problématique, plan). Puis échangez votre production avec un autre binôme : le lecteur doit vérifier la présence des quatre étapes et souligner les connecteurs. Enfin, deux volontaires lisent leur introduction à voix haute et la classe vote pour la plus claire.
\end{experience}

\begin{savaistu}[El comentario de texto, un exercice emblématique du monde hispanique]
Le \esp{comentario de texto} est un exercice emblématique des lycées d'Espagne : il figure dans l'épreuve de \esp{Lengua castellana} de la \esp{Selectividad}, l'équivalent espagnol du baccalauréat. Les élèves espagnols y commentent des articles de presse ou des textes littéraires, exactement comme vous le faites en espagnol LV2 au Cameroun. Maîtriser cette méthode, c'est donc acquérir une compétence universelle du monde hispanophone, utile aussi pour comprendre la presse de Guinée équatoriale, seul pays hispanophone d'Afrique subsaharienne — c'est là, à Malabo, que se déroule l'histoire de María dans notre texte d'entraînement.
\end{savaistu}

\begin{aretenir}[L'essentiel de la section]
\begin{itemize}
\item Le commentaire suit une architecture fixe : introduction (contexte + texte + problématique + plan), développement en deux axes, conclusion (bilan + ouverture).
\item Chaque paragraphe du développement suit la séquence \cle{idée -- citation -- analyse} ; on ne paraphrase jamais sans relever un procédé précis.
\item Les connecteurs (\esp{en primer lugar, además, sin embargo, por lo tanto, en conclusión}) structurent la copie et guident le correcteur.
\item Pour le texte sur le numérique : axe 1 = les risques (champs lexicaux, exemple vécu) ; axe 2 = les solutions (impératifs et subjonctif, visée didactique).
\end{itemize}
\end{aretenir}

\coursec{Production guidée : proposer des comportements responsables}

Dans la section précédente, nous avons appris à \cle{commenter un texte} : observer, relever les idées, analyser les procédés et formuler un jugement personnel. Nous passons maintenant de la lecture à l'\cle{écriture} : c'est à ton tour de produire un texte en espagnol. La tâche est concrète et proche de ta vie quotidienne : rédiger, pour le journal du lycée, un court article de conseils sur l'usage responsable du numérique. C'est exactement le type de production demandée à l'épreuve d'expression écrite du baccalauréat.

\courssub{La tâche finale}

\begin{definition}[La consigne]
Rédige un article de 12 à 15 lignes pour le journal de ton lycée, intitulé \esp{« Jóvenes conectados: consejos para un uso responsable »}. Ton article doit informer les élèves sur les risques du numérique et leur proposer des comportements responsables.
\end{definition}

Avant d'écrire, observons ce qu'on attend de toi. Un bon article de lycéen en espagnol possède quatre qualités : un \cle{plan clair} (introduction, développement, conclusion), le \cle{lexique du thème} (ici, le monde numérique), des \cle{structures grammaticales variées} (impératif pour conseiller, subjonctif après \esp{es importante que}) et des \cle{connecteurs} qui relient les idées. Nous allons construire ton article étape par étape.

\begin{popfigure}
\begin{tikzpicture}[
  etape/.style={rectangle, rounded corners=4pt, draw=popblue, fill=popblueL, align=center, minimum width=2.6cm, minimum height=1cm, font=\small\bfseries},
  fleche/.style={-{Stealth[length=3mm]}, thick, poporange}
]
\node[etape, align=center] (e1) at (0,0){1. Brainstorm\\ \\ormalfont\itshape ideas};
\node[etape, right=0.7cm of e1, align=center] (e2){2. Plan\\ \\ormalfont\itshape estructura};
\node[etape, right=0.7cm of e2, align=center] (e3){3. Rédaction\\ \\ormalfont\itshape redacción};
\node[etape, right=0.7cm of e3, align=center] (e4){4. Relecture\\ \\ormalfont\itshape revisión};
\node[etape, right=0.7cm of e4, align=center] (e5){5. Publication\\ \\ormalfont\itshape publicación};
\draw[fleche] (e1) -- (e2);
\draw[fleche] (e2) -- (e3);
\draw[fleche] (e3) -- (e4);
\draw[fleche] (e4) -- (e5);
\end{tikzpicture}
\legende{Les cinq étapes de la production écrite, à afficher en classe.}
\end{popfigure}

\courssub{Étape 1 — Le brainstorm : rassembler les idées}

On n'écrit jamais directement : on commence par noter des idées, sans ordre, en puisant dans le lexique de la leçon. Consigne : liste \textbf{trois risques} du numérique et \textbf{trois conseils} pour un usage responsable. Voici un exemple de tableau d'idées déjà rempli.

\begin{center}
\begin{tabularx}{\linewidth}{|X|X|X|}\hline
\rowcolor{popblueL} \textbf{Riesgos} (risques) & \textbf{Consejos} (conseils) & \textbf{Traduction du conseil} \\ \hline
el ciberacoso & no compartas tu contraseña & ne partage pas ton mot de passe \\ \hline
la adicción al móvil & limita el tiempo de pantalla & limite le temps d'écran \\ \hline
la pérdida de privacidad & configura la privacidad de tu perfil & règle la confidentialité de ton profil \\ \hline
descargar aplicaciones peligrosas & descarga solo aplicaciones seguras & ne télécharge que des applications sûres \\ \hline
hablar con desconocidos & no aceptes a desconocidos & n'accepte pas d'inconnus \\ \hline
\end{tabularx}
\end{center}

\begin{attention}[Faux amis et accords à surveiller dès le brainstorm]
\begin{itemize}
\item \esp{una aplicación segura} : l'adjectif \esp{segura} s'accorde en genre et en nombre avec \esp{aplicación} (féminin singulier). De même : \esp{datos personales} (masculin pluriel), \esp{una red social peligrosa}.
\item \esp{descargar} signifie « télécharger » (depuis internet vers ton appareil) ; \esp{cargar} signifie « charger » (la batterie). Ne confonds pas.
\item \esp{la contraseña} = le mot de passe ; n'écris jamais \esp{la contraseñe} et ne cherche pas à calquer « mot de passe » mot à mot.
\item \esp{no aceptes a desconocidos} : devant une personne, complément d'objet direct, l'espagnol exige la \textbf{préposition} \esp{a} (la « a personnelle ») : \esp{no aceptes \textbf{a} desconocidos}, \esp{ayudo \textbf{a} mi hermano}.
\end{itemize}
\end{attention}

\courssub{Étape 2 — Le plan imposé}

Un article de 12 à 15 lignes suit une architecture simple en quatre blocs. Respecte-le scrupuleusement : au baccalauréat, le respect du plan est un critère de notation.

\begin{enumerate}
\item \textbf{Introducción} (2-3 lignes) : un constat. \esp{Hoy en día, los jóvenes pasan muchas horas en internet y en las redes sociales.} (« Aujourd'hui, les jeunes passent beaucoup d'heures sur internet et sur les réseaux sociaux. »)
\item \textbf{Párrafo 1 : los riesgos} (3-4 lignes) : présente deux ou trois dangers (ciberacoso, adicción, privacidad).
\item \textbf{Párrafo 2 : los consejos} (4-5 lignes) : propose des solutions avec l'impératif et le subjonctif.
\item \textbf{Conclusión} (2-3 lignes) : un appel à l'action. \esp{¡Protege tu vida digital!} (« Protège ta vie numérique ! »)
\end{enumerate}

\begin{exemplebox}[Phrases-modèles à réutiliser (traduites)]
\begin{itemize}
\item \esp{En mi opinión, es importante que los jóvenes limiten el tiempo de pantalla.} = « À mon avis, il est important que les jeunes limitent le temps d'écran. »
\item \esp{No compartas tu contraseña con nadie.} = « Ne partage ton mot de passe avec personne. »
\item \esp{Es necesario que configuremos la privacidad de nuestras redes sociales.} = « Il est nécessaire que nous réglions la confidentialité de nos réseaux sociaux. »
\item \esp{Sin embargo, internet también tiene peligros.} = « Cependant, internet a aussi des dangers. »
\item \esp{Por eso, debemos usar el móvil de forma responsable.} = « C'est pourquoi nous devons utiliser le portable de façon responsable. »
\end{itemize}
\end{exemplebox}

\courssub{Étape 3 — Les contraintes linguistiques}

Ton article n'est pas libre à 100 \% : il doit obligatoirement contenir les éléments suivants. Coche-les pendant la relecture.

\begin{center}
\begin{tabularx}{\linewidth}{|X|X|}\hline
\rowcolor{popgreenL} \textbf{Contrainte} & \textbf{Exemple attendu} \\ \hline
Au moins 2 impératifs négatifs & \esp{no compartas}, \esp{no aceptes}, \esp{no descargues} \\ \hline
Au moins 2 structures « es importante que » + subjonctif & \esp{es importante que limites}, \esp{es importante que protejas} \\ \hline
Au moins 5 mots du lexique thématique & \esp{móvil, aplicación, contraseña, ciberacoso, privacidad, descargar, adicción} \\ \hline
Au moins 3 connecteurs & \esp{sin embargo, además, por eso, en primer lugar, en conclusión} \\ \hline
\end{tabularx}
\end{center}

\begin{attention}[Le piège du subjonctif]
Après \esp{es importante que}, le verbe est au \textbf{subjonctif présent}, jamais à l'indicatif ni à l'infinitif. On écrit \esp{es importante que \textbf{protejas} tus datos} et non \esp{es importante que proteges tus datos}. Rappel des formes : \esp{proteger} → que yo \esp{proteja}, tú \esp{protejas}, él \esp{proteja}, nosotros \esp{protejamos}, vosotros \esp{protejáis}, ellos \esp{protejan}. Attention aux verbes irréguliers : \esp{tener} → \esp{tengas} ; \esp{hacer} → \esp{hagas} ; \esp{ser} → \esp{seas} ; \esp{ir} → \esp{vayas}.
\end{attention}

\begin{aretenir}[Les deux outils pour conseiller]
\begin{itemize}
\item \textbf{Impératif affirmatif} (forme \esp{tú}) = 3\textsuperscript{e} personne du singulier de l'indicatif : \esp{limita}, \esp{configura}, \esp{protege}, \esp{descarga}.
\item \textbf{Impératif négatif} = \esp{no} + subjonctif présent : \esp{no compartas}, \esp{no aceptes}, \esp{no descargues}.
\item \textbf{Conseil impersonnel} = \esp{es importante que} / \esp{es necesario que} + subjonctif : \esp{es importante que protejas tus datos}.
\end{itemize}
\end{aretenir}

\courssub{Étape 4 — Rédaction, relecture en binôme et modèle d'article}

Tu rédiges d'abord \textbf{individuellement} (15 minutes), puis tu échanges ta copie avec un camarade qui la corrige à l'aide de la grille ci-dessous. Voici d'abord un modèle d'article réussi, rédigé par un élève de Première de Yaoundé.

\begin{textebox}[Jóvenes conectados: consejos para un uso responsable (article modèle)]
Hoy en día, los jóvenes pasan muchas horas en internet. El móvil es una herramienta fantástica para estudiar y comunicarse; sin embargo, también tiene peligros.

En primer lugar, existe el ciberacoso: algunos alumnos reciben mensajes crueles en las redes sociales. Además, la adicción a las pantallas roba el tiempo del sueño y de los estudios. Finalmente, si no configuramos la privacidad, nuestros datos personales están en peligro.

¿Qué podemos hacer? En mi opinión, es importante que los jóvenes limiten el tiempo de pantalla. También es importante que protejas tu información personal. No compartas tu contraseña con nadie y no aceptes a desconocidos en tus redes. Descarga solo aplicaciones seguras.

En conclusión, internet es un buen amigo si lo usamos bien. Por eso, ¡protégete y sé un internauta responsable!
\end{textebox}

\textbf{Glossaire} : \esp{herramienta} = outil ; \esp{roba} = vole (verbe \esp{robar}) ; \esp{el sueño} = le sommeil ; \esp{datos personales} = données personnelles ; \esp{desconocidos} = inconnus ; \esp{internauta} = internaute.

\textbf{Vérification du modèle} : 2 impératifs négatifs (\esp{no compartas}, \esp{no aceptes}) ; 2 subjonctifs (\esp{limiten}, \esp{protejas}) ; 6 mots du lexique ; 5 connecteurs (\esp{sin embargo, en primer lugar, además, finalmente, en conclusión}). Le plan est respecté. Remarque aussi les deux impératifs affirmatifs irréguliers de la conclusion : \esp{protégete} (impératif pronominal de \esp{protegerse}, avec accent sur le \esp{e} tonique) et \esp{sé} (impératif irrégulier de \esp{ser}).

\begin{methode}[Grille de relecture en binôme]
Corrige la copie de ton camarade en quatre passages, dans cet ordre :
\begin{enumerate}
\item \textbf{Plan} : y a-t-il une introduction, un paragraphe « riesgos », un paragraphe « consejos », une conclusion ? Souligne chaque partie.
\item \textbf{Lexique} : compte les mots du thème numérique (minimum 5). Entoure-les.
\item \textbf{Grammaire} : vérifie les 2 impératifs négatifs (\esp{no} + subjonctif) et les 2 subjonctifs après \esp{es importante que}. Vérifie aussi l'accord des adjectifs (\esp{una aplicación segura}).
\item \textbf{Connecteurs} : surligne au moins 3 connecteurs différents. S'il en manque, propose : \esp{además}, \esp{sin embargo}, \esp{por eso}.
\end{enumerate}
Rends la copie avec une remarque positive et un conseil d'amélioration.
\end{methode}

\begin{exoresolu}[Complète ce début d'article]
Énoncé : complète avec la forme correcte du verbe entre parenthèses. \esp{Es importante que tú (proteger) \ldots{} tus datos personales. No (compartir) \ldots{} tu contraseña con nadie y (limitar) \ldots{} el tiempo de pantalla. Es necesario que nosotros (configurar) \ldots{} la privacidad de nuestras redes sociales.}
\tcblower
Solution détaillée :
\begin{itemize}
\item \esp{Es importante que tú \textbf{protejas} tus datos personales.} → subjonctif présent, 2\textsuperscript{e} personne du singulier de \esp{proteger}.
\item \esp{No \textbf{compartas} tu contraseña con nadie} → impératif négatif = \esp{no} + subjonctif (\esp{compartas}).
\item \esp{y \textbf{limita} el tiempo de pantalla} → impératif affirmatif, forme \esp{tú} : \esp{limita} (comme la 3\textsuperscript{e} personne du singulier de l'indicatif).
\item \esp{Es necesario que nosotros \textbf{configuremos} la privacidad} → subjonctif présent, 1\textsuperscript{re} personne du pluriel.
\end{itemize}
Piège évité : après \esp{es importante que}, l'infinitif français (« il est important de protéger ») ne se calque pas en espagnol quand le sujet est précisé : il faut \esp{que} + subjonctif.
\end{exoresolu}

\courssub{Variante orale : el debate}

\begin{experience}[¿El móvil: amigo o enemigo?]
La classe se divise en deux groupes. Le groupe A défend la thèse \esp{« El móvil es un amigo »} ; le groupe B défend \esp{« El móvil es un enemigo »}.
\begin{enumerate}
\item Préparation (5 min) : chaque groupe note trois arguments avec le lexique de la leçon (\esp{comunicarse, aprender, adicción, ciberacoso, privacidad}).
\item Débat (10 min) : chaque élève prend la parole avec une structure argumentée : \esp{En mi opinión\ldots{} / Estoy de acuerdo porque\ldots{} / No estoy de acuerdo porque\ldots{} / Por ejemplo\ldots{}}
\item Conclusion : la classe vote et rédige ensemble une phrase de synthèse au tableau : \esp{El móvil es útil, pero es importante que lo usemos con responsabilidad.}
\end{enumerate}
\end{experience}

\courssub{Complément Bac (hors strict programme)}

\begin{methode}[Réussir l'expression écrite au baccalauréat]
Sujet type : \esp{« Escribe un artículo sobre las ventajas y los peligros de las redes sociales »} (80-100 mots). Quatre réflexes gagnants :
\begin{enumerate}
\item \textbf{Respecte le nombre de mots} : compte-les ; en dessous de 80 ou au-dessus de 120, tu perds des points.
\item \textbf{Utilise le lexique du thème} : le correcteur cherche les mots du champ lexical numérique (\esp{redes sociales, contraseña, privacidad, ciberacoso, descargar}).
\item \textbf{Varie les structures} : alterne impératif (\esp{no compartas}), subjonctif (\esp{es importante que protejas}) et conditionnel si possible (\esp{sería bueno que usaras el móvil con moderación}).
\item \textbf{Soigne les connecteurs} : \esp{en primer lugar, por un lado / por otro lado, sin embargo, además, en conclusión}. Ils donnent de la cohérence et font gagner des points.
\end{enumerate}
\end{methode}

\begin{exercice}[À toi d'écrire]
Rédige ton article \esp{« Jóvenes conectados: consejos para un uso responsable »} en 12 à 15 lignes en respectant le plan imposé et toutes les contraintes linguistiques de l'étape 3. Fais ensuite relire ta copie par ton binôme avec la grille, corrige-la, puis recopie la version finale pour l'affichage du journal du lycée.
\end{exercice}

\begin{savaistu}[El Día de Internet Segura]
Le 6 février est célébrée la Journée mondiale pour un internet plus sûr, en espagnol \esp{el Día de Internet Segura} (aussi appelé \esp{Día de un Internet más Seguro}). Dans de nombreuses écoles d'Espagne et d'Amérique latine, les élèves participent ce jour-là à des ateliers sur le \esp{ciberacoso}, la protection de la \esp{contraseña} et le respect de la \esp{privacidad}. Et si ton lycée organisait, lui aussi, sa journée de l'internet responsable ? Ce serait une belle occasion de publier ton article !
\end{savaistu}

\newpage

\coursec{Activité d'intégration}

\coursec{Leçon EA18 : Las tecnologías digitales y la comunicación}

\courssub{Objectifs de l'activité}

Cette activité d'intégration clôt la leçon \esp{Las tecnologías digitales y la comunicación}. Elle fait travailler les élèves en situation réelle de communication : concevoir et présenter une campagne de sensibilisation aux risques du numérique dans leur propre lycée. À l'issue de l'activité, l'élève doit être capable de :

\begin{itemize}
\item mobiliser le \cle{lexique} du monde numérique (\esp{internet, móvil, aplicación, contraseña, ciberacoso, adicción, privacidad, descargar}) dans un texte de production personnelle ;
\item employer correctement l'\cle{impératif} affirmatif et négatif pour formuler des slogans et des conseils (\esp{¡Protege tus datos!}, \esp{¡No compartas tu contraseña!}) ;
\item utiliser le \cle{subjonctif} présent après \esp{es importante que}, \esp{es necesario que}, \esp{conviene que} pour exprimer un conseil général ;
\item évaluer les risques du numérique (cyberharcèlement, addiction, vie privée) et proposer des comportements responsables, à l'écrit comme à l'oral ;
\item travailler en groupe, présenter un document en deux minutes et argumenter un choix collectif (vote de la classe).
\end{itemize}

\courssub{Situation}

Le club bilingue du lycée de Nkol-Eton (Yaoundé) organise la \emph{Semana de la Ciudadanía Digital}. Le proviseur a constaté que de nombreux élèves partagent leurs mots de passe, publient des photos sans réfléchir et passent plus de cinq heures par jour sur leur téléphone. Il demande donc aux classes de Première A de créer, en espagnol, des affiches de sensibilisation qui seront exposées dans le hall et dans la salle multimédia. Le thème officiel de la semaine est : \esp{« Piensa antes de publicar »}.

Pour lancer la réflexion, le professeur projette ce message, inspiré de faits réels observés dans plusieurs établissements :

\begin{textebox}[Mensaje del club bilingüe — Semana de la Ciudadanía Digital]
¡Atención, alumnos de Primera A!\\
La semana pasada, un alumno de nuestra escuela compartió su contraseña con un amigo. Ese amigo publicó fotos falsas en su nombre y ahora el alumno sufre ciberacoso. Otros estudiantes descargan aplicaciones sin leer las condiciones y pierden su privacidad. Algunos pasan seis horas al día con el móvil: es una verdadera adicción.\\
El club bilingüe os invita a crear una campaña de sensibilización: « Piensa antes de publicar ». Diseñad un cartel con un eslogan, tres riesgos y tres consejos. El mejor cartel decorará la sala de informática durante todo el año.\\
¡Participad y proteged vuestra vida digital!
\end{textebox}

\textbf{Glossaire} : \esp{compartir} = partager ; \esp{en su nombre} = en son nom ; \esp{las condiciones} = les conditions (d'utilisation) ; \esp{un cartel} = une affiche ; \esp{un eslogan} = un slogan ; \esp{diseñad} = concevez (impératif, 2\up{e} pers. du pluriel) ; \esp{decorará} = décorera (futur).

\courssub{Consignes}

Par groupes de quatre, vous créez une affiche de sensibilisation en espagnol destinée aux élèves du lycée. Suivez les étapes dans l'ordre.

\begin{enumerate}
\item \textbf{Comprendre la situation.} Lisez le message du club bilingüe. Relevez les trois risques du numérique qui y sont mentionnés et soulignez les mots du lexique de la leçon. Résumez oralement le problème en deux phrases en espagnol.

\item \textbf{Choisir un slogan.} Rédigez un slogan court et accrocheur avec un impératif (affirmatif ou négatif), par exemple : \esp{¡No compartas tu contraseña!} ou \esp{¡Protege tu privacidad!}. Le slogan doit être visible et mémorable.

\item \textbf{Présenter trois risques.} Illustrez sur l'affiche les trois risques étudiés : \esp{ciberacoso}, \esp{adicción}, \esp{privacidad}. Pour chaque risque, écrivez une phrase simple au présent de l'indicatif qui explique le danger (ex. : \esp{El ciberacoso destruye la autoestima de los jóvenes.}).

\item \textbf{Formuler trois conseils.} Rédigez trois conseils avec la structure \esp{es importante que} + subjonctif (ou \esp{es necesario que}, \esp{conviene que}). Attention à la conjugaison du subjonctif présent.

\item \textbf{Mettre en page.} Disposez le slogan en haut, les trois risques au centre (avec un petit dessin ou un symbole pour chacun) et les trois conseils en bas. Vérifiez l'orthographe, les accents et la ponctuation espagnole (¡ ! ¿ ?).

\item \textbf{Présenter à l'oral.} Chaque groupe présente son affiche en deux minutes : un membre lit le slogan, un autre explique les risques, un troisième lit les conseils, le quatrième conclut. Parlez lentement et articulez.

\item \textbf{Voter et justifier.} La classe vote pour la meilleure campagne. Chaque électeur justifie son vote en une phrase : \esp{Voto por el grupo tres porque...}. L'affiche gagnante sera exposée dans la salle multimédia.
\end{enumerate}

\begin{attention}[Critères d'évaluation]
L'affiche et la présentation sont notées sur trois critères : la richesse et la justesse du \cle{lexique} du numérique ; la correction \cle{grammaticale} (impératif et subjonctif sans faute, accents, ponctuation) ; la qualité \cle{argumentative} (risques clairement expliqués, conseils pertinents et adaptés aux lycéens camerounais).
\end{attention}

\courssub{Ressources à mobiliser}

\begin{aretenir}[Lexique utile : el mundo digital]
\esp{internet} (m., souvent sans article : \esp{acceso a internet} ; \esp{el internet} possible à l'oral) ; \esp{el móvil} (le portable) ; \esp{la aplicación / la app} ; \esp{la contraseña} (le mot de passe) ; \esp{el ciberacoso} (le cyberharcèlement) ; \esp{la adicción} ; \esp{la privacidad} ; \esp{descargar} (télécharger) ; \esp{compartir} (partager) ; \esp{publicar} (publier) ; \esp{proteger} (protéger) ; \esp{las redes sociales} (les réseaux sociaux) ; \esp{los datos personales} (les données personnelles) ; \esp{el acoso en línea} (synonyme de \esp{ciberacoso}) ; \esp{la huella digital} (l'empreinte numérique).
\end{aretenir}

\begin{propriete}[Rappel : l'impératif pour conseiller]
\begin{itemize}
\item \textbf{Impératif affirmatif (tú)} : 3\up{e} personne du singulier du présent de l'indicatif. \esp{protege}, \esp{comparte}, \esp{piensa} (diphtongaison \emph{e} $\to$ \emph{ie}).
\item \textbf{Irréguliers (tú) :} \esp{ten} (tener), \esp{haz} (hacer), \esp{pon} (poner), \esp{sal} (salir), \esp{ve} (ir), \esp{sé} (ser), \esp{di} (decir), \esp{ven} (venir).
\item \textbf{Impératif négatif (tú) :} \esp{no} + subjonctif présent. \esp{no compartas}, \esp{no publiques}, \esp{no descargues}, \esp{no protejas}.
\item \textbf{Changements orthographiques (impératif négatif / subjonctif) :}
\begin{center}
\begin{tabular}{|l|l|l|}\hline
\rowcolor{popblueL} Infinitif & Impératif affirmatif (tú) & Impératif négatif (tú) \\ \hline
publicar & publica & no publiques (c $\to$ qu) \\ \hline
descargar & descarga & no descargues (g $\to$ gu) \\ \hline
proteger & protege & no protejas (g $\to$ j) \\ \hline
jugar & juega & no juegues (u $\to$ ue, g $\to$ gu) \\ \hline
\end{tabular}
\end{center}
\item \textbf{Impératif (vosotros) :} Affirmatif : infinitif + \emph{-d} (\esp{participad}, \esp{proteged}). Négatif : \esp{no} + subjonctif (\esp{no participéis}, \esp{no protejáis}).
\end{itemize}
\end{propriete}

\begin{propriete}[Rappel : es importante que + subjonctif]
Après les expressions impersonnelles de conseil, nécessité ou convenance (\esp{es importante que}, \esp{es necesario que}, \esp{conviene que}, \esp{es mejor que}, \esp{es fundamental que}), le verbe qui suit est au \cle{subjonctif} présent.
\begin{itemize}
\item \textbf{Formation (réguliers) :} Radical de la 1\up{re} pers. du sg. de l'indicatif (\esp{yo}) + terminaisons \emph{inversées} : \emph{-ar} $\to$ \emph{-e, -es, -e, -emos, -éis, -en} ; \emph{-er/-ir} $\to$ \emph{-a, -as, -a, -amos, -áis, -an}.
\item \textbf{Verbes à diphtongaison (pensar, cerrar, empezar...) :} \esp{pienso} $\to$ \esp{pienses, piensa, pensemos, penséis, piensen}.
\item \textbf{Verbes en -ger/-gir (proteger, elegir...) :} \esp{protejo} $\to$ \esp{protejas} (g $\to$ j) ; \esp{elijo} $\to$ \esp{elijas} (g $\to$ j).
\item \textbf{Verbes en -car/-gar/-zar (buscar, jugar, empezar...) :} \esp{busco} $\to$ \esp{busque} (c $\to$ qu) ; \esp{juego} $\to$ \esp{juegue} (g $\to$ gu) ; \esp{empiezo} $\to$ \esp{empiece} (z $\to$ c).
\end{itemize}
\textbf{Exemples :} \esp{Es importante que protejas tus datos.} ; \esp{Es necesario que uses contraseñas seguras.} ; \esp{Conviene que limites el tiempo en el móvil.}
\end{propriete}

\courssub{Entraînement guidé (grammaire)}

\begin{exercice}[Mise en forme : impératif et subjonctif]
Transformez les phrases suivantes selon les consignes.
\begin{enumerate}
\item \esp{Compartes tu contraseña con tu amigo.} $\to$ Interdiction (tú) : \esp{¡No \underline{\hspace{3cm}}!}
\item \esp{Descargas aplicaciones sin leer las condiciones.} $\to$ Conseil négatif (tú) : \esp{¡No \underline{\hspace{3cm}}!}
\item \esp{Proteges tus datos personales.} $\to$ Conseil affirmatif (tú) : \esp{\underline{\hspace{3cm}} tus datos personales.}
\item \esp{Es importante que (tú) piensas antes de publicar.} $\to$ Corrigez le mode verbal : \esp{Es importante que \underline{\hspace{3cm}} antes de publicar.}
\item \esp{Es necesario que (ustedes) usan contraseñas difíciles.} $\to$ Corrigez : \esp{Es necesario que \underline{\hspace{3cm}} contraseñas difíciles.}
\item \esp{Conviene que (nosotros) hablamos cara a cara.} $\to$ Corrigez : \esp{Conviene que \underline{\hspace{3cm}} cara a cara.}
\end{enumerate}
\end{exercice}

\begin{corrige}[Exercice : Mise en forme]
\begin{enumerate}
\item \esp{¡No compartas!} (subjonctif \emph{compartir}, négatif).
\item \esp{¡No descargues!} (subjonctif \emph{descargar}, g $\to$ gu).
\item \esp{Protege} (impératif affirmatif régulier).
\item \esp{pienses} (subjonctif \emph{pensar}, diphtongaison).
\item \esp{usen} (subjonctif \emph{usar}, 3\up{e} pers. pl. pour \emph{ustedes}).
\item \esp{hablemos} (subjonctif \emph{hablar}, 1\up{re} pers. pl. pour \emph{nosotros}).
\end{enumerate}
\end{corrige}

\courssub{Production guidée : Campagne « Piensa antes de publicar »}

\begin{exoresolu}[Affiche modèle du groupe 1 — « Piensa antes de publicar »]
Énoncé : rédiger le contenu complet d'une affiche (slogan, trois risques, trois conseils) et la conclusion orale de la présentation.
\tcblower
\textbf{Production modèle (contenu de l'affiche) :}

\esp{¡PIENSA ANTES DE PUBLICAR! ¡Tu vida digital es tu vida!}

\esp{\textbf{Riesgo 1 — Ciberacoso:} El ciberacoso en las redes sociales destruye la autoestima y la reputación de los jóvenes.}

\esp{\textbf{Riesgo 2 — Adicción:} Pasar muchas horas con el móvil crea adicción y perjudica los estudios y el sueño.}

\esp{\textbf{Riesgo 3 — Privacidad:} Si compartes tu contraseña o tus fotos, pierdes tu privacidad para siempre.}

\esp{\textbf{Consejo 1:} Es importante que protejas tus datos personales con contraseñas difíciles.}

\esp{\textbf{Consejo 2:} Es necesario que pienses antes de publicar una foto o un comentario.}

\esp{\textbf{Consejo 3:} Conviene que limites el tiempo con el móvil y que hables con tus amigos cara a cara.}

\esp{¡No compartas tu contraseña! ¡No descargues aplicaciones sin leer las condiciones! ¡Denuncia el ciberacoso!}

\textbf{Conclusion orale :} \esp{En conclusión, el móvil es una herramienta útil, pero puede ser peligrosa. Pensad antes de publicar y proteged vuestra vida digital. ¡Gracias por vuestra atención!}

\textbf{Commentaire des choix de langue :}
\begin{itemize}
\item Le slogan principal utilise l'impératif affirmatif de \esp{pensar} : \esp{piensa} (diphtongaison \emph{e} $\to$ \emph{ie}), forme courte et percutante.
\item Les trois risques sont au présent de l'indicatif (vérités générales) : \esp{destruye}, \esp{crea}, \esp{perjudica}. Le troisième risque emploie \esp{si + présent, présent} (\esp{si compartes... pierdes}), structure correcte pour une conséquence réelle et immédiate.
\item Les trois conseils varient les expressions impersonnelles (\esp{es importante que}, \esp{es necesario que}, \esp{conviene que}) pour éviter les répétitions, chacune suivie du subjonctif : \esp{protejas} (\emph{proteger}, g $\to$ j), \esp{pienses} (\emph{pensar}, diphtongaison), \esp{limites} (régulier), \esp{hables} (\emph{hablar}, régulier).
\item La série finale d'impératifs négatifs (\esp{no compartas}, \esp{no descargues}) respecte les changements orthographiques : \emph{c $\to$ qu} pour \esp{publicar} (\esp{publiques}) et \emph{g $\to$ gu} pour \esp{descargar} (\esp{descargues}).
\item La conclusion orale passe au pluriel (\esp{pensad}, \esp{proteged}) car le groupe s'adresse à l'auditoire : impératif de la 2\up{e} personne du pluriel (vosotros), formé sur l'infinitif + \emph{-d}.
\item Le lexique de la leçon est entièrement mobilisé : \esp{ciberacoso}, \esp{adicción}, \esp{privacidad}, \esp{contraseña}, \esp{descargar}, \esp{aplicaciones}, \esp{móvil}, \esp{redes sociales}, \esp{datos personales}.
\end{itemize}
\end{exoresolu}

\courssub{Bilan de l'activité}

Cette campagne \esp{« Piensa antes de publicar »} a permis de réinvestir l'ensemble des acquis de la leçon dans une tâche finale concrète et motivante. Sur le plan de la \cle{compréhension}, les élèves ont lu et exploité un message authentique du club bilingüe. Sur le plan du \cle{lexique}, ils ont manipulé le vocabulaire du monde numérique dans un contexte qui leur parle : le téléphone portable, les réseaux sociaux, la vie du lycée à Yaoundé. Sur le plan \cle{grammatical}, l'impératif (affirmatif et négatif) et le subjonctif après \esp{es importante que} ont été employés non comme des exercices mécaniques, mais comme de véritables outils pour convaincre. Enfin, la présentation orale et le vote final ont développé l'\cle{expression orale}, l'argumentation et l'esprit de collaboration.

Le point de vigilance reste la conjugaison du subjonctif (changements orthographiques \emph{c/qu, g/gu, g/j, z/c} et diphtongaison) : une reprise collective des affiches après le vote permet de corriger les dernières erreurs avant l'exposition de la campagne gagnante. Cette activité prépare naturellement à l'évaluation de fin de module, où l'élève devra, seul, rédiger un court texte de conseils sur l'usage responsable du numérique.

\begin{savaistu}[Cameroun, Guinée équatoriale et monde hispanophone : la fracture numérique]
Au Cameroun, le taux de pénétration d'internet dépasse 40\% (données UIT récentes), porté par la 3G/4G et le \emph{Mobile Money} (Orange Money, MTN MoMo), véritable levier d'inclusion financière. Cependant, le coût des données reste élevé pour les lycéens. En Guinée équatoriale, seul pays hispanophone d'Afrique, l'espagnol est la langue de l'administration et de l'éducation ; l'accès au numérique y est un enjeu majeur de développement. En Espagne et en Amérique latine, des plans nationaux (\esp{Plan de Alfabetización Digital}, \esp{Conectar Igualdad} en Argentine) visent à réduire la fracture. Le \cle{ciberacoso} et la \cle{adicción a las pantallas} sont des préoccupations partagées : l'UNICEF alerte sur l'exposition précoce aux réseaux sociaux. Le slogan \esp{« Piensa antes de publicar »} s'applique donc de Yaoundé à Madrid en passant par Malabo ou Buenos Aires.
\end{savaistu}

\newpage

\coursec{Méthodes et exercices résolus}

\courssub{Méthodes et exercices résolus}

\begin{methode}[Lire et comprendre un texte sur le numérique]
\textbf{Objectif :} répondre à des questions de compréhension (examen : 6 à 8 points) en s'appuyant strictement sur le texte.
\begin{enumerate}
    \item \textbf{Lecture globale :} lire une première fois sans s'arrêter au vocabulaire inconnu. Identifier le \cle{thème} (les technologies numériques), la \cle{thèse} (avantages et risques) et la \cle{structure} (introduction, développement sur les risques, conclusion et conseils).
    \item \textbf{Lecture détaillée :} souligner les \cle{mots-clés} du titre et des questions. Repérer les \cle{connecteurs logiques} (\esp{sin embargo}, \esp{por tanto}, \esp{además}, \esp{en conclusión}) qui organisent l'argumentation.
    \item \textbf{Analyse des questions :} distinguer les questions \textbf{droit au texte} (réponse explicite dans le texte) des questions \textbf{d'inférence} (déduire le sens d'un mot, l'intention de l'auteur). Pour les questions « Justifiez par une citation », citer le texte entre guillemets est obligatoire.
    \item \textbf{Rédaction des réponses :} répondre en \textbf{espagnol}, en phrases complètes. Utiliser ses propres mots (paraphrase) sauf pour les citations demandées. Vérifier l'accord sujet-verbe et l'orthographe des mots recopiés.
\end{enumerate}
\cle{Réflexe examen} : pour « Explica con tus propias palabras », ne recopiez pas la phrase du texte ; reformulez avec du vocabulaire connu (synonymes, structures simples).
\end{methode}

\begin{textebox}[Conectados, ¿pero protegidos? (texte original)]
Vivimos en la era de la hiperconexión. El smartphone se ha convertido en una extensión de nuestra mano: descargamos aplicaciones para todo, compartimos nuestra ubicación, nuestras fotos y nuestros pensamientos sin leer las condiciones de privacidad. Sin embargo, esta comodidad tiene un precio. El ciberacoso afecta a miles de jóvenes que no saben cómo bloquear a un agresor ni a quién denunciar. La adicción a las pantallas provoca insomnio, ansiedad y fracaso escolar. Además, nuestras contraseñas débiles abren la puerta al robo de identidad. Es fundamental educar en un uso crítico: no compartas datos sensibles, configura la privacidad de tus perfiles y desconecta una hora antes de dormir. La tecnología es una herramienta, no un fin en sí misma.
\end{textebox}

\emph{Traduction :} « Nous vivons à l'ère de l'hyperconnexion. Le smartphone est devenu une extension de notre main : nous téléchargeons des applications pour tout, nous partageons notre localisation, nos photos et nos pensées sans lire les conditions de confidentialité. Cependant, ce confort a un prix. Le cyberharcèlement touche des milliers de jeunes qui ne savent ni comment bloquer un agresseur ni à qui le signaler. L'addiction aux écrans provoque insomnie, anxiété et échec scolaire. De plus, nos mots de passe faibles ouvrent la porte au vol d'identité. Il est fondamental d'éduquer à un usage critique : ne partage pas de données sensibles, configure la confidentialité de tes profils et déconnecte une heure avant de dormir. La technologie est un outil, pas une fin en soi. »

\begin{exoresolu}[Compréhension écrite : « Conectados, ¿pero protegidos? »]
\textbf{Questions (sur le texte ci-dessus) :}
\begin{enumerate}
    \item ¿Cuál es el tema principal del texto? (1 punto)
    \item Según el texto, ¿cuáles son tres consecuencias negativas de la adicción a las pantallas? (3 puntos)
    \item « El smartphone se ha convertido en una extensión de nuestra mano ». Explica esta imagen con tus propias palabras. (2 puntos)
    \item ¿Qué tres consejos da el autor para un uso responsable? Cita el texto. (3 puntos)
    \item Justifica con una cita que el autor no está en contra de la tecnología en sí. (1 punto)
\end{enumerate}

\tcblower

\textbf{Solution détaillée :}

\textbf{1. Tema principal :} \esp{El texto trata de los riesgos de la hiperconexión digital (ciberacoso, adicción, privacidad) y de la necesidad de educar en un uso responsable y crítico de la tecnología.}
\emph{Justification :} la lecture globale montre l'opposition (connecteur \esp{sin embargo}) entre l'omniprésence du smartphone et ses dangers, suivie de solutions (\esp{es fundamental}).

\textbf{2. Conséquences de l'addiction :}
\begin{itemize}
    \item \esp{insomnio} (troubles du sommeil) ;
    \item \esp{ansiedad} (anxiété) ;
    \item \esp{fracaso escolar} (échec scolaire).
\end{itemize}
\emph{Note :} recopier exactement les mots du texte assure les points pour une question « droit au texte ».

\textbf{3. Explication de l'image :}
\esp{Significa que el teléfono móvil se ha vuelto indispensable e inseparable de la vida diaria: lo llevamos siempre con nosotros y lo usamos para casi todas las actividades, como si fuera una parte más de nuestro cuerpo.}
\emph{Analyse :} la métaphore de l'extension du corps exprime la dépendance et l'omniprésence. Réponse en espagnol, structures simples (\esp{se ha vuelto}, \esp{lo usamos}, \esp{como si fuera} + imparfait du subjonctif, correct après \esp{como si}).

\textbf{4. Trois conseils (citation obligatoire) :}
\begin{enumerate}
    \item \esp{« no compartas datos sensibles »} (impératif négatif) ;
    \item \esp{« configura la privacidad de tus perfiles »} (impératif affirmatif) ;
    \item \esp{« desconecta una hora antes de dormir »} (impératif affirmatif).
\end{enumerate}

\textbf{5. Citation pro-technologie (dernière phrase) :}
\esp{« La tecnología es una herramienta, no un fin en sí misma. »}
\emph{Analyse :} l'auteur valorise l'outil (\esp{herramienta}) mais refuse l'aliénation (\esp{fin en sí misma}).

\textbf{Bilan :} 10/10. Respect des consignes : réponses en espagnol, citations exactes entre guillemets, phrases complètes.
\end{exoresolu}

\begin{methode}[Évaluer les risques du numérique : vocabulaire et structures d'analyse]
\textbf{Objectif :} identifier, nommer et expliquer les risques (cyberharcèlement, addiction, vie privée) avec le lexique précis du module et des structures de cause et de conséquence.
\begin{enumerate}
    \item \textbf{Identifier le risque :} utiliser le nom exact : \cle{el ciberacoso}, \cle{la adicción a las pantallas}, \cle{la pérdida de privacidad}, \cle{el robo de identidad}, \cle{las contraseñas débiles}.
    \item \textbf{Exprimer la cause :} \esp{debido a / a causa de + nom} (\esp{Debido a la adicción...}) ; \esp{como / porque + indicatif} (\esp{Como no configuras la privacidad...}) ; \esp{el uso excesivo de + nom + provoca + nom}.
    \item \textbf{Exprimer la conséquence :} \esp{provoca / causa / genera + nom} ; \esp{por eso / por tanto / como resultado + phrase}.
    \item \textbf{Porter un jugement critique :} \esp{es grave que / es preocupante que / es inaceptable que} + \textbf{subjonctif} (évaluation subjective). Ex. : \esp{Es grave que los jóvenes sufran ciberacoso.} « Il est grave que les jeunes subissent du cyberharcèlement. »
\end{enumerate}
\cle{Tableau de correspondance FR/ES (à mémoriser) :}
\begin{center}
\begin{tabularx}{\linewidth}{|X|X|X|}\hline
\rowcolor{popblueL} \textbf{Français} & \textbf{Español} & \textbf{Remarque} \\ \hline
Cyberharcèlement & \esp{el ciberacoso} & \esp{Acoso} = harcèlement ; verbe : \esp{acosar} ; en ligne : \esp{ciberacosar}. \\ \hline
Addiction (aux écrans) & \esp{la adicción a las pantallas} & Adjectif : \esp{adicto/a a} ; on dit \esp{ser adicto a las redes}. \\ \hline
Vie privée / confidentialité & \esp{la privacidad} & \esp{Privado/a} = privé (faux ami partiel : \esp{educación privada} = enseignement privé, sens correct). \\ \hline
Mot de passe & \esp{la contraseña} & Féminin ; synonyme : \esp{la clave}. \\ \hline
Télécharger & \esp{descargar} & Verbe régulier en \esp{-ar} ; antonyme : \esp{subir}. \\ \hline
Application & \esp{la aplicación} (la \esp{app}) & Féminin ; accent sur la \esp{ó}. \\ \hline
Signaler / bloquer & \esp{denunciar} / \esp{bloquear} & Verbes réguliers en \esp{-ar}. \\ \hline
\end{tabularx}
\end{center}
\end{methode}

\begin{textebox}[Los peligros de la red (texte à traduire)]
Muchos adolescentes ignoran los riesgos de publicar su rutina diaria en las redes sociales. Comparten su ubicación en tiempo real, lo que facilita el trabajo de los ciberdelincuentes. Además, el uso de contraseñas débiles, como «123456» o su fecha de nacimiento, permite el acceso no autorizado a sus cuentas. Las consecuencias pueden ser devastadoras: desde el robo de identidad hasta el acoso psicológico. Por eso, los expertos recomiendan activar la verificación en dos pasos y revisar la configuración de privacidad cada mes.
\end{textebox}

\begin{exoresolu}[Version et analyse des risques : « Los peligros de la red »]
\textbf{Énoncé :}
\begin{enumerate}
    \item \textbf{Version :} traduisez en français le texte « Los peligros de la red » ci-dessus.
    \item \textbf{Analyse lexicale :} classez les mots suivants dans trois colonnes : \textbf{Nom du risque} / \textbf{Verbe d'action (conseil ou protection)} / \textbf{Conséquence}.
    \esp{ignorar, riesgos, publicar, ubicación, facilitar, ciberdelincuentes, contraseñas, acceso, consecuencias, devastadoras, robo, acoso, recomiendan, activar, verificación, revisar, configuración.}
\end{enumerate}

\tcblower

\textbf{Solution détaillée :}

\textbf{1. Version (français) :}
« De nombreux adolescents ignorent les risques de publier leur routine quotidienne sur les réseaux sociaux. Ils partagent leur localisation en temps réel, ce qui facilite le travail des cybercriminels. De plus, l'utilisation de mots de passe faibles, comme "123456" ou leur date de naissance, permet l'accès non autorisé à leurs comptes. Les conséquences peuvent être dévastatrices : du vol d'identité au harcèlement psychologique. C'est pourquoi les experts recommandent d'activer la vérification en deux étapes et de vérifier les paramètres de confidentialité chaque mois. »

\emph{Points de vigilance (francophones) :}
\begin{itemize}
    \item \esp{Ignoran los riesgos} $\rightarrow$ « ignorent les risques » est ici la bonne traduction : \esp{ignorar} signifie « ne pas connaître » ou « ne pas tenir compte de ». Attention au contexte : \esp{ignoro la respuesta} = « je ne connais pas la réponse ».
    \item \esp{Ubicación en tiempo real} $\rightarrow$ « localisation en temps réel » (ne pas garder « ubicación »).
    \item \esp{Ciberdelincuentes} $\rightarrow$ « cybercriminels ».
    \item \esp{Acceso no autorizado} $\rightarrow$ « accès non autorisé » (terme technique).
    \item \esp{Verificación en dos pasos} $\rightarrow$ « vérification en deux étapes » (double authentification).
    \item \esp{Revisar} $\rightarrow$ « vérifier, contrôler » (et non « réviser » au sens scolaire ici).
\end{itemize}

\textbf{2. Classement lexical :}

\begin{center}
\begin{tabularx}{\linewidth}{|X|X|X|}\hline
\rowcolor{popblueL} \textbf{Nom du risque / problème} & \textbf{Verbe d'action} & \textbf{Conséquence / résultat} \\ \hline
\esp{los riesgos} & \esp{ignorar} (comportement à risque) & \esp{las consecuencias} \\ \hline
\esp{la ubicación} (donnée sensible) & \esp{publicar} (comportement à risque) & \esp{devastadoras} (adjectif) \\ \hline
\esp{los ciberdelincuentes} (acteurs) & \esp{facilitar} (effet négatif) & \esp{el robo (de identidad)} \\ \hline
\esp{las contraseñas (débiles)} & \esp{recomiendan} (conseil) & \esp{el acoso (psicológico)} \\ \hline
\esp{el acceso (no autorizado)} & \esp{activar} (protection) &  \\ \hline
\esp{la verificación} & \esp{revisar} (protection) &  \\ \hline
\esp{la configuración} &  &  \\ \hline
\end{tabularx}
\end{center}

\emph{Note pédagogique :} \esp{publicar}, \esp{compartir}, \esp{ignorar} désignent des \textbf{comportements à risque}. \esp{recomendar}, \esp{activar}, \esp{revisar}, \esp{configurar}, \esp{bloquear}, \esp{denunciar} sont des \textbf{verbes de protection et de conseil}, souvent employés à l'impératif ou au subjonctif.
\end{exoresolu}

\begin{methode}[L'impératif pour conseiller et interdire (Grammaire 1)]
\textbf{Objectif :} maîtriser la formation et l'emploi de l'impératif (affirmatif et négatif) pour formuler des conseils de sécurité numérique (\esp{No compartas}, \esp{Configura}, \esp{Desconecta}).
\begin{enumerate}
    \item \textbf{Impératif AFFIRMATIF :}
    \begin{itemize}
        \item \textbf{tú} : 3\up{e} personne du singulier du présent de l'indicatif (\esp{habla, come, vive, configura, desconecta, bloquea}). \cle{Exceptions} : 8 verbes irréguliers : \esp{di} (decir), \esp{haz} (hacer), \esp{ve} (ir), \esp{pon} (poner), \esp{sal} (salir), \esp{sé} (ser), \esp{ten} (tener), \esp{ven} (venir).
        \item \textbf{usted / ustedes} : subjonctif présent, 3\up{e} personne : \esp{configure, configuren}.
        \item \textbf{vosotros} : remplacer le \esp{-r} de l'infinitif par \esp{-d} : \esp{hablad, comed, vivid, configurad}. \cle{Exception} : \esp{id} (ir).
    \end{itemize}
    \item \textbf{Impératif NÉGATIF (toutes les personnes) :} \esp{no} + \textbf{subjonctif présent}.
    \begin{itemize}
        \item \esp{tú : no hables, no comas, no vivas, no compartas, no descargues}.
        \item \textbf{Irrégularités du subjonctif (donc de l'impératif négatif) :} verbes en \esp{-car / -gar / -zar} (\esp{no saques, no pagues, no empieces}), verbes à diphtongaison (\esp{no cuentes, no pienses, no duermas}), et \esp{ir} (\esp{no vayas}), \esp{ser} (\esp{no seas}), \esp{estar} (\esp{no estés}), \esp{saber} (\esp{no sepas}), \esp{dar} (\esp{no des}), \esp{haber} (\esp{no hayas}).
    \end{itemize}
    \item \textbf{Place des pronoms :}
    \begin{itemize}
        \item Affirmatif : pronoms \textbf{enclitiques} (collés après le verbe, avec accent écrit pour conserver l'accentuation) : \esp{Configúrala, Bloquéalo, Dímelo, Desconéctate}.
        \item Négatif : pronoms \textbf{proclitiques} (entre \esp{no} et le verbe) : \esp{No la configures, No lo bloquees, No me lo digas, No te desconectes}.
    \end{itemize}
    \item \textbf{Emploi :} ordres, conseils, instructions, interdictions. Dans un conseil amical ou éducatif : \esp{tú} ; dans un registre formel : \esp{usted}.
\end{enumerate}
\cle{Tableau récapitulatif des verbes du module :}
\begin{center}
\begin{tabular}{|l|l|l|l|}\hline
\rowcolor{popblueL} \textbf{Infinitif} & \textbf{Imp. aff. (tú)} & \textbf{Imp. nég. (tú)} & \textbf{Avec pronom (aff. / nég.)} \\ \hline
\esp{configurar} & \esp{configura} & \esp{no configures} & \esp{configúralo / no lo configures} \\ \hline
\esp{compartir} & \esp{comparte} & \esp{no compartas} & \esp{compártelo / no lo compartas} \\ \hline
\esp{descargar} & \esp{descarga} & \esp{no descargues} & \esp{descárgala / no la descargues} \\ \hline
\esp{bloquear} & \esp{bloquea} & \esp{no bloquees} & \esp{bloquéalo / no lo bloquees} \\ \hline
\esp{denunciar} & \esp{denuncia} & \esp{no denuncies} & \esp{denúncialo / no lo denuncies} \\ \hline
\esp{desconectar(se)} & \esp{desconecta} & \esp{no te desconectes} & \esp{desconéctate / no te desconectes} \\ \hline
\esp{poner} & \esp{pon} (irrég.) & \esp{no pongas} & \esp{pónla / no la pongas} \\ \hline
\end{tabular}
\end{center}
\end{methode}

\begin{exoresolu}[Transformation et thème : l'impératif en situation]
\textbf{Énoncé :}
\begin{enumerate}
    \item \textbf{Transformation :} mettez les verbes entre parenthèses à l'impératif (conseil à un camarade = \esp{tú}).
    \begin{enumerate}
        \item \esp{\_\_\_\_\_\_\_\_\_\_ (configurar) la privacidad de tu perfil ahora.}
        \item \esp{\_\_\_\_\_\_\_\_\_\_ (no / compartir) tu contraseña con nadie.}
        \item \esp{\_\_\_\_\_\_\_\_\_\_ (bloquear) al usuario que te insulta.}
        \item \esp{\_\_\_\_\_\_\_\_\_\_ (no / descargar) aplicaciones de fuentes desconocidas.}
        \item \esp{\_\_\_\_\_\_\_\_\_\_ (poner) una contraseña segura (mayúsculas, números, símbolos).}
        \item \esp{\_\_\_\_\_\_\_\_\_\_ (no / conectarse) a redes wifi públicas sin VPN.}
    \end{enumerate}
    \item \textbf{Thème (français $\rightarrow$ espagnol) :} traduisez ces conseils avec l'impératif \esp{tú} et les pronoms nécessaires.
    \begin{enumerate}
        \item Protège ta vie privée. (verbe : \esp{proteger})
        \item Ne l'accepte pas (la demande d'ami). (verbe : \esp{aceptar})
        \item Signale-le (le harceleur) immédiatement. (verbe : \esp{denunciar})
        \item Ne t'inquiète pas. (verbe : \esp{preocuparse})
        \item Déconnecte-toi une heure avant de dormir. (verbe : \esp{desconectarse})
    \end{enumerate}
\end{enumerate}

\tcblower

\textbf{Solution détaillée :}

\textbf{1. Transformation (impératif \esp{tú}) :}
\begin{enumerate}
    \item \esp{Configura} (régulier en \esp{-ar} : 3\up{e} pers. sg. du présent de l'indicatif).
    \item \esp{No compartas} (négatif $\rightarrow$ subjonctif : \esp{compartir} $\rightarrow$ \esp{no compartas}). \cle{Attention :} jamais \esp{no compartes} (l'indicatif est interdit à l'impératif négatif).
    \item \esp{Bloquea} (régulier en \esp{-ar}).
    \item \esp{No descargues} (négatif $\rightarrow$ subjonctif). \cle{Orthographe :} \esp{-gar} $\rightarrow$ \esp{-gue} devant \esp{e} pour garder le son [g] : \esp{descargue, descargues}.
    \item \esp{Pon} (impératif affirmatif irrégulier : \esp{di, haz, ve, pon, sal, sé, ten, ven}). \cle{Piège :} ne pas écrire \esp{pone}.
    \item \esp{No te conectes} (verbe pronominal \esp{conectarse} ; négatif $\rightarrow$ \esp{no} + pronom \esp{te} + subjonctif \esp{conectes}).
\end{enumerate}

\textbf{2. Thème :}
\begin{enumerate}
    \item \esp{Protege tu privacidad.} (impératif affirmatif régulier ; COD « ta vie privée » = nom, pas de pronom nécessaire).
    \item \esp{No la aceptes.} (négatif : \esp{no} + pronom COD \esp{la} (la solicitud, féminin) + subjonctif \esp{aceptes}).
    \item \esp{Denúncialo inmediatamente.} (affirmatif : \esp{denuncia} + \esp{lo} $\rightarrow$ \esp{denúncialo}, accent écrit sur l'antépénultième syllabe).
    \item \esp{No te preocupes.} (pronominal \esp{preocuparse} ; négatif : \esp{no} + \esp{te} + subjonctif \esp{preocupes}). \cle{Remarque :} \esp{preocupar} n'est pas un verbe à diphtongaison ; le radical reste \esp{preocup-} partout.
    \item \esp{Desconéctate una hora antes de dormir.} (affirmatif pronominal : \esp{desconecta} + \esp{te} $\rightarrow$ \esp{desconéctate}, accent sur l'antépénultième syllabe).
\end{enumerate}

\textbf{Bilan grammatical :} formation correcte (indicatif à l'affirmatif, subjonctif au négatif), irrégulier \esp{pon}, orthographe \esp{-gue}, place des pronoms (enclitiques à l'affirmatif, proclitiques au négatif) et accents écrits.
\end{exoresolu}

\begin{methode}[Le subjonctif après les expressions de nécessité et de conseil (Grammaire 2)]
\textbf{Objectif :} utiliser le subjonctif présent dans les subordonnées introduites par \esp{es importante que}, \esp{es necesario que}, \esp{es fundamental que}, \esp{recomiendo que}, \esp{quiero que}, pour exprimer un conseil, une exigence ou une évaluation subjective.
\begin{enumerate}
    \item \textbf{La règle de base :} volonté, émotion, expression impersonnelle d'évaluation, recommandation, doute $\rightarrow$ \textbf{subjonctif} dans la subordonnée s'il y a \textbf{changement de sujet}.
    \begin{center}
    \esp{Sujet 1 + verbe de volonté/évaluation (indicatif) + QUE + sujet 2 + verbe (subjonctif)}
    \end{center}
    \item \textbf{Expressions impersonnelles du module :}
    \esp{es importante que}, \esp{es fundamental que}, \esp{es necesario que}, \esp{es vital que}, \esp{es urgente que}, \esp{está prohibido que}.
    \item \textbf{Verbes de volonté et de conseil :}
    \esp{recomiendo que}, \esp{sugiero que}, \esp{aconsejo que}, \esp{exijo que}, \esp{pido que}, \esp{quiero que}, \esp{prefiero que}.
    \item \textbf{Formation du subjonctif présent (rappel) :}
    \begin{enumerate}
        \item prendre la 1\up{re} personne (\esp{yo}) du présent de l'indicatif : \esp{configuro, tengo, pongo, sé, voy} ;
        \item enlever le \esp{-o} ;
        \item ajouter les \textbf{terminaisons inversées} : \esp{-ar} $\rightarrow$ \esp{-e, -es, -e, -emos, -éis, -en} ; \esp{-er / -ir} $\rightarrow$ \esp{-a, -as, -a, -amos, -áis, -an}.
    \end{enumerate}
    \item \textbf{6 verbes totalement irréguliers au subjonctif (à connaître par cœur) :}
    \begin{center}
    \begin{tabular}{|l|l|l|}\hline
    \rowcolor{popblueL} \textbf{Infinitif} & \textbf{Subjonctif présent} & \textbf{Exemple} \\ \hline
    \esp{ir} & \esp{vaya, vayas, vaya, vayamos, vayáis, vayan} & \esp{Es importante que vayas al ciberespacio seguro.} \\ \hline
    \esp{ser} & \esp{sea, seas, sea, seamos, seáis, sean} & \esp{Es fundamental que la contraseña sea segura.} \\ \hline
    \esp{estar} & \esp{esté, estés, esté, estemos, estéis, estén} & \esp{Quiero que estés protegido.} \\ \hline
    \esp{saber} & \esp{sepa, sepas, sepa, sepamos, sepáis, sepan} & \esp{Es necesario que sepas bloquear.} \\ \hline
    \esp{dar} & \esp{dé, des, dé, demos, deis, den} & \esp{Te pido que no des tus datos.} \\ \hline
    \esp{haber} & \esp{haya, hayas, haya, hayamos, hayáis, hayan} & \esp{Es importante que no haya ciberacoso.} \\ \hline
    \end{tabular}
    \end{center}
    \item \textbf{Cas sans changement de sujet :} si le sujet est général ou identique, on utilise l'\textbf{infinitif} : \esp{Es importante desconectarse.} « Il est important de se déconnecter. »
\end{enumerate}
\cle{Attention francophones :} ne dites pas \esp{Es importante que tú configuras} (indicatif). Le conseil projette une action non réalisée $\rightarrow$ \textbf{subjonctif} : \esp{Es importante que tú configures}.
\end{methode}

\begin{exoresolu}[Choix du mode : indicatif ou subjonctif ? + thème]
\textbf{Énoncé :}
\begin{enumerate}
    \item \textbf{Complétion :} conjuguez le verbe entre parenthèses au présent de l'indicatif ou du subjonctif selon le sens.
    \begin{enumerate}
        \item \esp{Es fundamental que los jóvenes \_\_\_\_\_\_\_\_\_\_ (proteger) su privacidad.}
        \item \esp{Los expertos confirman que el ciberacoso \_\_\_\_\_\_\_\_\_\_ (afectar) a muchos estudiantes.}
        \item \esp{No creo que esa aplicación \_\_\_\_\_\_\_\_\_\_ (ser) segura.}
        \item \esp{Mi padre quiere que yo \_\_\_\_\_\_\_\_\_\_ (apagar) el móvil a las diez.}
        \item \esp{Es verdad que internet \_\_\_\_\_\_\_\_\_\_ (tener) riesgos.}
        \item \esp{Te recomiendo que no \_\_\_\_\_\_\_\_\_\_ (aceptar) a desconocidos.}
    \end{enumerate}
    \item \textbf{Thème :} traduisez en espagnol avec une structure impersonnelle ou un verbe de volonté + subjonctif.
    \begin{enumerate}
        \item Il est nécessaire que les adolescents bloquent les harceleurs.
        \item Je recommande que tu ne partages pas ta localisation.
        \item Il est interdit que les mineurs accèdent à ce contenu.
        \item Les professeurs exigent que nous changions nos mots de passe.
        \item Il est important de se déconnecter le soir.
    \end{enumerate}
\end{enumerate}

\tcblower

\textbf{Solution détaillée :}

\textbf{1. Complétion (mode justifié) :}
\begin{enumerate}
    \item \esp{protejan} (\textbf{subjonctif}). \esp{Es fundamental que} (nécessité, expression impersonnelle) + changement de sujet. \esp{Proteger} : \esp{yo protejo} $\rightarrow$ \esp{protejan}.
    \item \esp{afecta} (\textbf{indicatif}). \esp{Confirman que} : verbe de certitude, fait constaté.
    \item \esp{sea} (\textbf{subjonctif}). \esp{No creo que} : doute, négation de la certitude. \esp{Ser} est irrégulier : \esp{sea}.
    \item \esp{apague} (\textbf{subjonctif}). \esp{Quiere que} : volonté + changement de sujet. \esp{Apagar} : \esp{yo apago} $\rightarrow$ \esp{apague} (orthographe \esp{-gue}).
    \item \esp{tiene} (\textbf{indicatif}). \esp{Es verdad que} : certitude, fait objectif.
    \item \esp{aceptes} (\textbf{subjonctif}). \esp{Te recomiendo que} : conseil + changement de sujet ; la négation \esp{no} précède le verbe.
\end{enumerate}

\textbf{2. Thème :}
\begin{enumerate}
    \item \esp{Es necesario que los adolescentes bloqueen a los acosadores.} (\esp{bloquear} $\rightarrow$ \esp{bloqueen}, 3\up{e} pers. pl.).
    \item \esp{Recomiendo que no compartas tu ubicación.} (\esp{compartir} $\rightarrow$ \esp{compartas} ; \esp{no} devant le verbe).
    \item \esp{Está prohibido que los menores accedan a ese contenido.} (\esp{acceder} $\rightarrow$ \esp{accedan}).
    \item \esp{Los profesores exigen que cambiemos nuestras contraseñas.} (\esp{cambiar} $\rightarrow$ \esp{cambiemos} ; attention, la forme est identique à l'indicatif \esp{cambiamos} sauf pour la voyelle thématique : \esp{-emos} au subjonctif).
    \item \esp{Es importante desconectarse por la noche.} \cle{Règle :} pas de changement de sujet (sens général « on ») $\rightarrow$ \textbf{infinitif}. Avec un sujet précis : \esp{Es importante que te desconectes por la noche.}
\end{enumerate}

\textbf{Points clés :} distinction certitude (indicatif) / volonté et subjectivité (subjonctif) ; irrégulier \esp{sea} ; terminaisons inversées ; place de la négation.
\end{exoresolu}

\begin{methode}[Production écrite : rédiger une liste de conseils responsables (Savoir-faire 3)]
\textbf{Objectif :} produire un texte cohérent (80 à 120 mots, niveau A2+/B1) proposant des comportements responsables face aux risques numériques. Formats possibles : article de blog, affiche de prévention, lettre ouverte.
\begin{enumerate}
    \item \textbf{Analyse du sujet :} identifier le \cle{format}, le \cle{destinataire} (camarades, parents, direction), le \cle{ton} (engagé, pédagogique) et les \cle{contraintes} (mots obligatoires : \esp{ciberacoso, contraseña, privacidad, adicción} ; structures imposées : impératif, subjonctif).
    \item \textbf{Plan type (3 parties) :}
    \begin{enumerate}
        \item \textbf{Accroche / constat :} \esp{Actualmente... / Hoy en día... / En la era digital...} + problème (\esp{El ciberacoso aumenta. La adicción al móvil es real.}).
        \item \textbf{Conseils (corps) :} 3 ou 4 mesures concrètes. \textbf{Structures obligatoires :}
        \begin{itemize}
            \item \esp{Es importante / fundamental / necesario que + subjonctif} ;
            \item \esp{Recomiendo / Aconsejo que + subjonctif} ;
            \item \textbf{impératif direct} : \esp{¡Bloquea al acosador! ¡No descargues aplicaciones desconocidas!} ;
            \item connecteurs : \esp{en primer lugar, además, por otro lado, finalmente}.
        \end{itemize}
        \item \textbf{Conclusion / appel à l'action :} \esp{En conclusión... / Para terminar...} + phrase choc : \esp{La tecnología nos sirve, no nos esclaviza. ¡Piensa antes de hacer clic!}
    \end{enumerate}
    \item \textbf{Relecture (check-list) :}
    \begin{itemize}
        \item accents : \esp{configuración, contraseña, móvil, público, jóvenes} ;
        \item accords sujet-verbe et adjectifs ;
        \item subjonctif après \esp{que} (volonté, nécessité) ;
        \item impératif : formation et place des pronoms ;
        \item vocabulaire précis du module (\esp{riesgos, ciberacoso, adicción} plutôt que le vague \esp{problemas}).
    \end{itemize}
\end{enumerate}
\end{methode}

\begin{textebox}[Article modèle : Conectados, sí; seguros, también]
Actualmente, en nuestro colegio de Yaundé, el smartphone es indispensable. Sin embargo, existen riesgos graves: el ciberacoso en WhatsApp y la adicción a TikTok causan ansiedad y bajan las notas. Es fundamental que todos protejamos nuestra privacidad.

En primer lugar, configura tu perfil en «privado» y no compartas tu ubicación. En segundo lugar, usa contraseñas fuertes y activa la verificación en dos pasos. Finalmente, desconéctate una hora antes de dormir para evitar el insomnio.

Recomiendo que denunciéis cualquier acoso a un adulto de confianza. La tecnología es una herramienta para aprender, no para sufrir. ¡Piénsalo antes de hacer clic!
\end{textebox}

\emph{Traduction :} « Aujourd'hui, dans notre lycée de Yaoundé, le smartphone est indispensable. Cependant, il existe des risques graves : le cyberharcèlement sur WhatsApp et l'addiction à TikTok provoquent de l'anxiété et font baisser les notes. Il est fondamental que nous protégions tous notre vie privée. D'abord, configure ton profil en "privé" et ne partage pas ta localisation. Ensuite, utilise des mots de passe forts et active la vérification en deux étapes. Enfin, déconnecte-toi une heure avant de dormir pour éviter l'insomnie. Je recommande que vous signaliez tout harcèlement à un adulte de confiance. La technologie est un outil pour apprendre, pas pour souffrir. Réfléchis avant de cliquer ! »

\begin{exoresolu}[Rédaction : « Consejos para una vida digital sana »]
\textbf{Énoncé :} vous êtes délégué(e) de votre classe à Yaoundé. Rédigez un court article (90 à 110 mots) pour le journal du lycée intitulé \esp{« Conectados, sí; seguros, también »}. Vous devez :
\begin{itemize}
    \item dénoncer deux risques (cyberharcèlement, addiction) ;
    \item donner trois conseils concrets en utilisant l'impératif et le subjonctif ;
    \item utiliser au moins 5 mots du lexique du module.
\end{itemize}
(Le texte modèle est donné dans l'encadré ci-dessus.)

\tcblower

\textbf{Analyse de la copie modèle (grille d'évaluation) :}
\begin{itemize}
    \item \textbf{Contenu :} 2 risques cités (\esp{ciberacoso, adicción}), 3 conseils structurés (\esp{en primer lugar, en segundo lugar, finalmente}), plus de 5 mots du lexique (\esp{ciberacoso, adicción, privacidad, ubicación, contraseñas, herramienta, denunciar}).
    \item \textbf{Grammaire :} subjonctifs corrects (\esp{protejamos}, \esp{denunciéis}), impératifs variés (affirmatifs : \esp{configura, usa, activa, desconéctate} ; négatif : \esp{no compartas}), expression impersonnelle \esp{es fundamental que}.
    \item \textbf{Vocabulaire :} précis, contexte camerounais (\esp{colegio de Yaundé}), registre adapté à un journal scolaire.
    \item \textbf{Orthographe et accents :} \esp{contraseña, móvil, ubicación, verificación, desconéctate, piénsalo, denunciéis} (accent sur le \esp{í} du hiatus \esp{-éis}). Zéro faute d'accent.
    \item \textbf{Cohérence :} connecteurs, progression logique du constat vers les solutions, ton engageant.
\end{itemize}
\textbf{Note estimée : 14/16 (Très bien).}
\end{exoresolu}

\begin{methode}[Méthode du commentaire de texte (épreuve écrite)]
\textbf{Objectif :} réussir l'exercice de commentaire (souvent 10 à 12 points) en suivant une structure rigoureuse : introduction $\rightarrow$ analyse détaillée (2 ou 3 axes) $\rightarrow$ conclusion ouverte.
\begin{enumerate}
    \item \textbf{Introduction (3 étapes) :} \textbf{1. Accroche} (contexte général : \esp{En la sociedad actual, la tecnología...}). \textbf{2. Présentation du document} (nature : \esp{artículo de opinión / blog} ; titre ; source ; thème). \textbf{3. Problématique + annonce du plan} (\esp{¿Hasta qué punto...?} $\rightarrow$ \esp{Analizaremos primero..., luego... y finalmente...}).
    \item \textbf{Développement (2 ou 3 axes) :} chaque axe = \textbf{idée directrice} + \textbf{citation courte} intégrée dans la phrase + \textbf{analyse} (lien idée-texte, vocabulaire, grammaire, intention de l'auteur). \cle{Utiliser :} \esp{El autor destaca que...}, \esp{Se observa el uso de...}, \esp{Esta frase revela...}, \esp{Por tanto...}.
    \item \textbf{Axes types pour ce module :}
    \begin{itemize}
        \item Axe 1 : \textbf{le constat alarmant} (hyperconnexion, ignorance des risques ; indicatifs \esp{vivimos, descargamos, compartimos}).
        \item Axe 2 : \textbf{la typologie des dangers} (champ lexical du danger : \esp{precio, ciberacoso, adicción, robo de identidad} ; verbes de conséquence : \esp{provoca, afecta}).
        \item Axe 3 : \textbf{la proposition éducative} (structures de conseil : \esp{es fundamental que + subjonctif}, impératifs ; le ton passe de l'alerte à la solution).
    \end{itemize}
    \item \textbf{Conclusion (2 étapes) :} \textbf{1. Synthèse} (réponse à la problématique). \textbf{2. Ouverture} (lien avec l'actualité camerounaise ou équato-guinéenne, expérience personnelle, autre document).
\end{enumerate}
\cle{Gestion du temps conseillée :} 15 min de lecture et d'annotation, 45 min de rédaction, 10 min de relecture.
\end{methode}

\begin{exoresolu}[Commentaire modèle : « Conectados, ¿pero protegidos? »]
\textbf{Énoncé :} commentez le texte « Conectados, ¿pero protegidos? » en répondant à la problématique : \esp{« ¿Cómo alerta el texto sobre los peligros digitales y propone una ciudadanía digital responsable? »} (150 à 200 mots).

\tcblower

\textbf{Commentaire rédigé (modèle) :}

\textbf{Introducción.} \esp{En la era de la hiperconexión, el smartphone se ha vuelto inseparable de la vida de los jóvenes, en Camerún como en el mundo hispano. El texto « Conectados, ¿pero protegidos? », un artículo de divulgación, denuncia la distancia entre la comodidad digital y la vulnerabilidad real de los usuarios. ¿Cómo alerta el autor sobre los peligros para construir después una ciudadanía digital responsable? Analizaremos primero el diagnóstico alarmante, luego la tipología de los riesgos y finalmente la pedagogía de la protección.}

\textbf{1. Un diagnóstico alarmante.} \esp{El autor describe una realidad cotidiana: « Vivimos en la era de la hiperconexión ». El presente de indicativo (vivimos, descargamos, compartimos) presenta los hechos como constatables y normales. La metáfora « extensión de nuestra mano » sugiere una dependencia casi física. La frase « sin leer las condiciones de privacidad » revela la negligencia de los usuarios.}

\textbf{2. La tipología de los riesgos.} \esp{El conector adversativo « sin embargo » introduce el campo léxico del peligro: precio, ciberacoso, adicción, insomnio, ansiedad, fracaso escolar, robo de identidad. La acumulación de sustantivos crea un efecto de amenaza creciente. El autor muestra que la tecnología daña tanto la seguridad (robo de identidad) como la salud mental (ansiedad, insomnio).}

\textbf{3. La pedagogía de la protección.} \esp{El giro final (« Es fundamental educar ») cambia el registro: de la descripción a la prescripción. Tres estructuras gramaticales operan aquí: la expresión impersonal « es fundamental » seguida de infinitivo, los imperativos directos (no compartas, configura, desconecta) que responsabilizan al lector, y la definición final: « La tecnología es una herramienta, no un fin en sí misma », que resume la filosofía del uso crítico.}

\textbf{Conclusión.} \esp{En síntesis, el texto pasa de la denuncia (indicativo, hechos) a la acción (imperativo, consejos). En Camerún y en Guinea Ecuatorial, donde el acceso a internet crece rápidamente, este mensaje es esencial: la mejor contraseña sigue siendo la educación.}

\textbf{Auto-évaluation :} structure introduction-développement-conclusion respectée ; citations intégrées ; analyse grammaticale (indicatif du constat contre impératif du conseil) reliée au sens ; vocabulaire précis ; ouverture contextuelle (Cameroun, Guinée équatoriale). \textbf{Note estimée : 11/12.}
\end{exoresolu}

\begin{attention}[Les 5 erreurs qui coûtent des points à l'examen (EA18)]
\begin{enumerate}
    \item \textbf{Accents oubliés ou mal placés :}
    \begin{itemize}
        \item \esp{contraseña} (tilde sur le \esp{ñ}), \esp{móvil, público, jóvenes, configuración, ubicación, adicción} : accents obligatoires.
        \item \textbf{Impératif + pronom :} \esp{configúralo, desconéctate, dímelo, denúncialo} (accent écrit pour conserver l'accentuation). Mais à la forme négative : \esp{no lo configures} (pas d'accent ajouté).
    \end{itemize}
    \item \textbf{Faux ami : \esp{ignorar}.} Il signifie « ne pas connaître » ou « ne pas tenir compte de » selon le contexte : \esp{Ignoran los riesgos} = « Ils ignorent les risques » (ils n'en tiennent pas compte). Mais \esp{Ignoro la respuesta} = « Je ne connais pas la réponse ». Ne pas traduire mécaniquement.
    \item \textbf{Calque du français « conseiller à quelqu'un de faire ».} \cle{Interdit :} \esp{Aconsejo a ti configurar...} \cle{Correct :} \esp{Te aconsejo que configures...} (verbe de conseil + \esp{que} + subjonctif).
    \item \textbf{Indicatif au lieu du subjonctif} après \esp{es importante que}, \esp{recomiendo que}, \esp{quiero que}. \cle{Faux :} \esp{Es importante que tú configuras}. \cle{Correct :} \esp{Es importante que tú configures}.
    \item \textbf{Accords de base :}
    \begin{itemize}
        \item féminins : \esp{la privacidad, la seguridad, la adicción, la contraseña} $\rightarrow$ \esp{una contraseña segura} ;
        \item masculins : \esp{el ciberacoso, el riesgo, el acceso, el robo, el móvil} ;
        \item pluriel : \esp{los jóvenes usan / sufren / deben} (jamais \esp{los jóvenes usa}).
    \end{itemize}
\end{enumerate}
\end{attention}

\newpage

\coursec{Exercices}

\courssub{Je vérifie mes connaissances}

\begin{exercice}[QCM : el mundo digital]
Pour chaque question, entoure la bonne réponse.

\begin{enumerate}
\item \esp{Una contraseña segura debe tener...}
\begin{enumerate}
\item \esp{solo tu nombre y tu fecha de nacimiento.}
\item \esp{letras, números y símbolos, y no debe compartirse.}
\item \esp{la misma palabra para todas las cuentas.}
\end{enumerate}
\item \esp{El ciberacoso es...}
\begin{enumerate}
\item \esp{una aplicación para editar fotos.}
\item \esp{un acoso o burla a través de internet o del móvil.}
\item \esp{un programa para proteger el ordenador.}
\end{enumerate}
\item \esp{Descargar una aplicación significa...}
\begin{enumerate}
\item \esp{instalarla en el móvil desde internet.}
\item \esp{borrarla del teléfono.}
\item \esp{comprar un teléfono nuevo.}
\end{enumerate}
\item \esp{La privacidad en internet consiste en...}
\begin{enumerate}
\item \esp{publicar todo lo que hacemos cada día.}
\item \esp{controlar quién puede ver nuestros datos personales.}
\item \esp{no usar nunca el móvil.}
\end{enumerate}
\end{enumerate}
\end{exercice}

\begin{exercice}[Vrai ou faux ? Justifie ta réponse]
Dis si les affirmations suivantes sont vraies (V) ou fausses (F) et justifie ta réponse en une phrase en français.

\begin{enumerate}
\item \esp{« No compartas tu contraseña »} est un impératif négatif à la deuxième personne du singulier.
\item Après \esp{es importante que}, on emploie l'indicatif présent.
\item \esp{No descargues} est la forme correcte de l'impératif négatif de \esp{descargar} avec \esp{tú}.
\item \esp{Es importante que protejas tu privacidad} signifie « Il est important que tu protèges ta vie privée ».
\item Le verbe \esp{hablar} à l'impératif affirmatif avec \esp{tú} donne \esp{hablas}.
\end{enumerate}
\end{exercice}

\begin{exercice}[Lexique : trouve le mot]
Complète chaque phrase avec le mot du champ lexical qui convient : \esp{aplicación -- contraseña -- ciberacoso -- adicción -- privacidad -- descargar -- móvil -- datos personales}.

\begin{enumerate}
\item Para entrar en mi cuenta, escribo mi nombre de usuario y mi \ldots\ldots\ldots
\item No debes \ldots\ldots\ldots{} programas de páginas desconocidas : pueden tener virus.
\item WhatsApp es una \ldots\ldots\ldots{} muy popular entre los jóvenes cameruneses.
\item Los insultos y las burlas por internet son una forma de \ldots\ldots\ldots
\item Pasar ocho horas al día delante de la pantalla puede ser una \ldots\ldots\ldots
\item En las redes sociales, hay que proteger la \ldots\ldots\ldots{} y no publicarlo todo.
\item Mi hermano usa el \ldots\ldots\ldots{} para llamar, chatear y hacer fotos.
\item Nunca des tu dirección ni tu número de teléfono : son \ldots\ldots\ldots
\end{enumerate}
\end{exercice}

\begin{exercice}[Conjugaison à compléter]
Complète le tableau avec la forme correcte de l'impératif affirmatif et négatif (\esp{tú}) des verbes suivants.

\begin{center}
\begin{tabular}{|l|l|l|}
\hline
\rowcolor{popblueL} Infinitivo & Imperativo afirmativo (tú) & Imperativo negativo (tú) \\
\hline
\esp{hablar} & \ldots\ldots\ldots & \ldots\ldots\ldots \\
\hline
\esp{compartir} & \ldots\ldots\ldots & \ldots\ldots\ldots \\
\hline
\esp{proteger} & \ldots\ldots\ldots & \ldots\ldots\ldots \\
\hline
\esp{verificar} & \ldots\ldots\ldots & \ldots\ldots\ldots \\
\hline
\esp{descargar} & \ldots\ldots\ldots & \ldots\ldots\ldots \\
\hline
\esp{leer} & \ldots\ldots\ldots & \ldots\ldots\ldots \\
\hline
\esp{escribir} & \ldots\ldots\ldots & \ldots\ldots\ldots \\
\hline
\esp{apagar} & \ldots\ldots\ldots & \ldots\ldots\ldots \\
\hline
\esp{hacer} & \ldots\ldots\ldots & \ldots\ldots\ldots \\
\hline
\esp{ser} & \ldots\ldots\ldots & \ldots\ldots\ldots \\
\hline
\end{tabular}
\end{center}
\end{exercice}

\courssub{Je m'entraîne}

\begin{exercice}[Thème : cinq conseils à traduire]
Traduis en espagnol les cinq conseils suivants. Attention à l'impératif négatif et au subjonctif.

\begin{enumerate}
\item Ne partage pas tes photos.
\item Il est important que tu protèges ta vie privée.
\item Parle à un adulte.
\item Il vaut mieux que tu vérifies l'information.
\item Ne télécharge pas cette application.
\end{enumerate}
\end{exercice}

\begin{exercice}[Transformations grammaticales]
\textbf{A.} Transforme les phrases suivantes en conseils à l'impératif (\esp{tú}), affirmatif ou négatif selon le sens.

\begin{enumerate}
\item \esp{(Tú) compartir tu contraseña con tus amigos.} $\rightarrow$ \ldots\ldots\ldots
\item \esp{(Tú) hablar con tus padres si tienes un problema.} $\rightarrow$ \ldots\ldots\ldots
\item \esp{(Tú) publicar fotos de otras personas sin permiso.} $\rightarrow$ \ldots\ldots\ldots
\item \esp{(Tú) verificar las noticias antes de compartirlas.} $\rightarrow$ \ldots\ldots\ldots
\item \esp{(Tú) usar el móvil durante toda la noche.} $\rightarrow$ \ldots\ldots\ldots
\end{enumerate}

\textbf{B.} Transforme les phrases à l'infinitif en phrases avec \esp{es importante que} + subjonctif.

\begin{enumerate}
\item \esp{Proteger tus datos personales.} $\rightarrow$ \esp{Es importante que} \ldots\ldots\ldots
\item \esp{Apagar el móvil antes de dormir.} $\rightarrow$ \esp{Es importante que} \ldots\ldots\ldots
\item \esp{Leer bien los mensajes antes de responder.} $\rightarrow$ \esp{Es importante que} \ldots\ldots\ldots
\item \esp{Ser prudente con los desconocidos en internet.} $\rightarrow$ \esp{Es importante que} \ldots\ldots\ldots
\item \esp{Hacer deporte y no solo mirar la pantalla.} $\rightarrow$ \esp{Es importante que} \ldots\ldots\ldots
\end{enumerate}
\end{exercice}

\begin{exercice}[Texte à trous : la privacidad]
Complète le texte suivant avec dix mots de la liste : \esp{contraseña -- descargar -- ciberacoso -- datos personales -- privacidad -- aplicaciones -- compartir -- verificar -- proteger -- redes sociales}.

\begin{textebox}[Mi vida digital]
Hoy casi todos los jóvenes usan las \ldots\ldots\ldots{} para comunicarse con sus amigos. Pero hay que tener cuidado: es fundamental \ldots\ldots\ldots{} nuestra \ldots\ldots\ldots{} en internet. Primero, no debemos publicar nuestros \ldots\ldots\ldots{}, como la dirección o el número de teléfono. Segundo, hay que crear una \ldots\ldots\ldots{} difícil y no dársela a nadie. Tercero, antes de \ldots\ldots\ldots{} una aplicación nueva, debemos leer bien la información. Algunas \ldots\ldots\ldots{} piden demasiados permisos. También es importante \ldots\ldots\ldots{} las noticias antes de difundirlas. Por último, si alguien sufre insultos en internet, es decir, \ldots\ldots\ldots{}, debe hablar con un adulto. No hay que \ldots\ldots\ldots{} fotos íntimas con desconocidos. ¡Naveguemos con responsabilidad!
\end{textebox}
\end{exercice}

\begin{exercice}[Appariements : conseils et situations]
Associe chaque situation (1--6) au conseil qui convient (a--f). Écris la réponse sous la forme 1--c.

\textbf{Situations :}
\begin{enumerate}
\item Un desconocido te pide fotos personales por WhatsApp.
\item Ves una noticia sorprendente en Facebook.
\item Un compañero recibe insultos en un grupo de chat.
\item Pasas cinco horas seguidas jugando con el móvil.
\item Quieres instalar un juego nuevo de una página desconocida.
\item Tu amiga quiere usar tu cuenta de Instagram.
\end{enumerate}

\textbf{Conseils :}
\begin{enumerate}[label=\alph*.]
\item \esp{No compartas tu contraseña, ni siquiera con tus amigos.}
\item \esp{No respondas y habla inmediatamente con un adulto.}
\item \esp{Es importante que verifiques la información antes de compartirla.}
\item \esp{No descargues aplicaciones de páginas no oficiales.}
\item \esp{Es importante que descanses y apagues la pantalla con frecuencia.}
\item \esp{Es importante que le aconsejes hablar con un profesor o con sus padres.}
\end{enumerate}
\end{exercice}

\courssub{Je me prépare au Bac}

\begin{exercice}[Compréhension écrite : la adicción a las pantallas]
Lis le texte suivant, puis réponds aux questions.

\begin{textebox}[Jóvenes y pantallas: ¿una nueva adicción?]
En España, como en muchos países del mundo, los adolescentes pasan cada vez más tiempo delante de las pantallas. Según los profesores y los padres, muchos jóvenes miran su móvil más de cinco horas al día: redes sociales, videojuegos, vídeos y mensajes ocupan casi todo su tiempo libre.

Esta situación preocupa a los especialistas. La adicción a las pantallas puede tener consecuencias graves: los jóvenes duermen peor, tienen dificultades para concentrarse en clase y a veces abandonan el deporte y las actividades con los amigos. Además, el uso excesivo de las redes sociales puede provocar ansiedad, porque los adolescentes se comparan constantemente con las imágenes perfectas que ven en internet.

Sin embargo, no todo es negativo. Las tecnologías digitales también son herramientas útiles para aprender, comunicarse con la familia y descubrir el mundo. El problema no es la tecnología, sino el uso que hacemos de ella.

Los expertos dan varios consejos. Primero, es importante que los jóvenes apaguen el móvil una hora antes de dormir. Segundo, es necesario que limiten el tiempo de las pantallas durante la semana. Tercero, es fundamental que hablen con sus padres o profesores si se sienten mal. Finalmente, hay que practicar deporte y mantener relaciones reales con los amigos.

Como dice un profesor de Madrid: «La tecnología debe ser una herramienta en nuestras manos, no una cadena en nuestros pies.»
\end{textebox}

\textbf{Questions de compréhension} (réponds en espagnol, avec des phrases complètes) :
\begin{enumerate}
\item ¿Cuántas horas al día miran el móvil muchos jóvenes, según el texto?
\item Menciona dos consecuencias negativas de la adicción a las pantallas.
\item ¿Por qué las redes sociales pueden provocar ansiedad en los adolescentes?
\item Según el texto, ¿las tecnologías digitales son siempre negativas? Justifica tu respuesta.
\item Escribe dos consejos que dan los expertos para usar bien las pantallas.
\end{enumerate}

\textbf{Question d'expression personnelle} :
\begin{enumerate}[start=6]
\item ¿Y tú? ¿Cuánto tiempo pasas delante de las pantallas cada día? ¿Qué haces para limitar tu uso del móvil? (4--5 frases)
\end{enumerate}

\textbf{Questions de langue} :
\begin{enumerate}[start=7]
\item Trouve dans le texte deux subjonctifs après \esp{es importante que} ou une expression équivalente, et recopie les phrases complètes.
\item Mets à l'impératif négatif (\esp{tú}) : \esp{mirar el móvil toda la noche} ; \esp{compararte con los demás}.
\end{enumerate}
\end{exercice}

\begin{exercice}[Thème et version]
\textbf{A. Thème.} Traduis en espagnol les phrases suivantes (attention aux impératifs et aux subjonctifs) :
\begin{enumerate}
\item Ne partage pas tes photos personnelles sur internet.
\item Il est important que tu protèges ta vie privée.
\item Parle à un adulte si tu as un problème.
\item Il vaut mieux que tu vérifies l'information avant de la publier.
\item Ne télécharge pas cette application : elle est dangereuse.
\end{enumerate}

\textbf{B. Version.} Traduis en français le paragraphe suivant :

\begin{textebox}[Contra el ciberacoso]
El ciberacoso es un problema serio entre los jóvenes. Si recibes mensajes insultantes, no respondas con agresividad y no guardes silencio: habla con tus padres o con un profesor. Es importante que conserves los mensajes como prueba y que bloquees a la persona que te ataca. Es importante también que los amigos de la víctima la ayuden y no compartan los insultos. Recuerda: detrás de cada pantalla hay una persona real, con sentimientos reales.
\end{textebox}
\end{exercice}

\begin{exercice}[Expression écrite : artículo para el periódico del instituto]
\textbf{Sujet :} \esp{Escribe un artículo (80--100 palabras) para el periódico del instituto sobre los riesgos de las redes sociales y los consejos para usarlas con responsabilidad.}

\textbf{Consignes :}
\begin{itemize}
\item Ton article doit comporter : un titre, une introduction (présentation du sujet), un paragraphe sur les risques (\esp{ciberacoso, adicción, privacidad}), un paragraphe avec au moins trois conseils (utilise l'impératif et \esp{es importante que} + subjonctif), et une courte conclusion.
\item Utilise au moins six mots du lexique étudié : \esp{móvil, aplicación, contraseña, ciberacoso, adicción, privacidad, descargar, datos personales}.
\item Soigne l'orthographe et les accents.
\end{itemize}

\textbf{Grille d'évaluation indicative :} respect de la consigne et du nombre de mots (4 pts) ; richesse du lexique (4 pts) ; correction grammaticale : impératif et subjonctif (6 pts) ; organisation et cohérence du texte (4 pts) ; orthographe et accents (2 pts).
\end{exercice}

\newpage

\coursec{Corrigés des exercices}

\coursec{Corrigés des exercices — Leçon EA18 : Las tecnologías digitales y la comunicación}

\begin{corrige}[Exercice 1 — QCM : El mundo digital]
\begin{enumerate}
\item \textbf{Bonnes réponses : 1-b ; 2-b ; 3-a ; 4-b.}
\end{enumerate}

\textbf{Justifications :}
\begin{enumerate}
\item \textbf{b.} \esp{Una contraseña segura debe tener letras, números y símbolos, y no debe compartirse.} « Un mot de passe sûr doit contenir des lettres, des chiffres et des symboles, et ne doit pas être partagé. » Les réponses a et c décrivent des mots de passe faibles et dangereux.
\item \textbf{b.} \esp{El ciberacoso es un acoso o burla a través de internet o del móvil.} « Le cyberharcèlement est un harcèlement ou une moquerie à travers internet ou le téléphone portable. »
\item \textbf{a.} \esp{Descargar una aplicación significa instalarla en el móvil desde internet.} « Télécharger une application signifie l'installer sur le portable depuis internet. » Attention au faux ami : \esp{descargar} = « télécharger », et non « décharger ».
\item \textbf{b.} \esp{La privacidad en internet consiste en controlar quién puede ver nuestros datos personales.} « La vie privée sur internet consiste à contrôler qui peut voir nos données personnelles. »
\end{enumerate}
\end{corrige}

\begin{corrige}[Exercice 2 — Vrai ou faux ?]
\begin{enumerate}
\item \textbf{Vrai.} \esp{No compartas} est bien un impératif négatif à la deuxième personne du singulier (\esp{tú}) : l'impératif négatif se construit avec \esp{no} + subjonctif présent (\esp{compartir} $\rightarrow$ \esp{no compartas}).
\item \textbf{Faux.} Après \esp{es importante que}, on emploie le \textbf{subjonctif} présent, car cette expression impersonnelle exprime un jugement, un souhait ou une nécessité : \esp{Es importante que protejas tu privacidad} (et non \emph{proteges}).
\item \textbf{Vrai.} \esp{Descargar} : subjonctif présent \esp{descargues} $\rightarrow$ impératif négatif \esp{no descargues}. « Ne télécharge pas. »
\item \textbf{Vrai.} \esp{Es importante que protejas tu privacidad} signifie exactement « Il est important que tu protèges ta vie privée » ; \esp{protejas} est le subjonctif présent de \esp{proteger}, deuxième personne du singulier.
\item \textbf{Faux.} L'impératif affirmatif de \esp{hablar} avec \esp{tú} est \esp{habla} (sans \emph{-s}) : \esp{¡Habla con un adulto!} « Parle à un adulte ! » La forme \esp{hablas} est le présent de l'indicatif. En revanche, l'impératif \emph{négatif} reprend le \emph{-s} : \esp{no hables}.
\end{enumerate}
\end{corrige}

\begin{corrige}[Exercice 3 — Lexique : trouve le mot]
\begin{enumerate}
\item \esp{Para entrar en mi cuenta, escribo mi nombre de usuario y mi \textbf{contraseña}.} « Pour entrer dans mon compte, j'écris mon nom d'utilisateur et mon mot de passe. »
\item \esp{No debes \textbf{descargar} programas de páginas desconocidas : pueden tener virus.} « Tu ne dois pas télécharger de programmes de pages inconnues : ils peuvent contenir des virus. »
\item \esp{WhatsApp es una \textbf{aplicación} muy popular entre los jóvenes cameruneses.} « WhatsApp est une application très populaire parmi les jeunes Camerounais. »
\item \esp{Los insultos y las burlas por internet son una forma de \textbf{ciberacoso}.} « Les insultes et les moqueries sur internet sont une forme de cyberharcèlement. »
\item \esp{Pasar ocho horas al día delante de la pantalla puede ser una \textbf{adicción}.} « Passer huit heures par jour devant l'écran peut être une addiction. »
\item \esp{En las redes sociales, hay que proteger la \textbf{privacidad} y no publicarlo todo.} « Sur les réseaux sociaux, il faut protéger sa vie privée et ne pas tout publier. »
\item \esp{Mi hermano usa el \textbf{móvil} para llamar, chatear y hacer fotos.} « Mon frère utilise le portable pour appeler, discuter et faire des photos. »
\item \esp{Nunca des tu dirección ni tu número de teléfono : son \textbf{datos personales}.} « Ne donne jamais ton adresse ni ton numéro de téléphone : ce sont des données personnelles. »
\end{enumerate}
\end{corrige}

\begin{corrige}[Exercice 4 — Conjugaison à compléter]
\begin{center}
\begin{tabular}{|l|l|l|}
\hline
\rowcolor{popblueL} Infinitivo & Imperativo afirmativo (tú) & Imperativo negativo (tú) \\
\hline
\esp{hablar} & \esp{habla} & \esp{no hables} \\
\hline
\esp{compartir} & \esp{comparte} & \esp{no compartas} \\
\hline
\esp{proteger} & \esp{protege} & \esp{no protejas} \\
\hline
\esp{verificar} & \esp{verifica} & \esp{no verifiques} \\
\hline
\esp{descargar} & \esp{descarga} & \esp{no descargues} \\
\hline
\esp{leer} & \esp{lee} & \esp{no leas} \\
\hline
\esp{escribir} & \esp{escribe} & \esp{no escribas} \\
\hline
\esp{apagar} & \esp{apaga} & \esp{no apagues} \\
\hline
\esp{hacer} & \esp{haz} & \esp{no hagas} \\
\hline
\esp{ser} & \esp{sé} & \esp{no seas} \\
\hline
\end{tabular}
\end{center}

\textbf{Rappels grammaticaux :}
\begin{itemize}
\item Impératif affirmatif (\esp{tú}) : il reprend la 3e personne du singulier du présent de l'indicatif : \esp{habla}, \esp{comparte}, \esp{escribe}.
\item Impératif négatif : \esp{no} + subjonctif présent : \esp{no hables}, \esp{no compartas}, \esp{no escribas}.
\item Verbes irréguliers : \esp{hacer} $\rightarrow$ \esp{haz} / \esp{no hagas} ; \esp{ser} $\rightarrow$ \esp{sé} / \esp{no seas}.
\item Changements orthographiques au subjonctif : \esp{proteger} $\rightarrow$ \esp{protejas} (\emph{g} $\rightarrow$ \emph{j} devant \emph{a}) ; \esp{verificar} $\rightarrow$ \esp{verifiques} (\emph{c} $\rightarrow$ \emph{qu}) ; \esp{apagar} $\rightarrow$ \esp{apagues} (\emph{g} $\rightarrow$ \emph{gu}).
\end{itemize}

\textbf{Points de vigilance :} Ne pas confondre l'impératif affirmatif \esp{habla} (sans \emph{s}) et le présent de l'indicatif \esp{hablas} (avec \emph{s}). À l'impératif négatif, le \emph{s} réapparaît : \esp{no hables}.
\end{corrige}

\begin{corrige}[Exercice 5 — Thème : cinq conseils à traduire]
\begin{enumerate}
\item « Ne partage pas tes photos. » $\rightarrow$ \esp{No compartas tus fotos.} (impératif négatif : \esp{no} + subjonctif \esp{compartas})
\item « Il est important que tu protèges ta vie privée. » $\rightarrow$ \esp{Es importante que protejas tu privacidad.} (\esp{protejas} : subjonctif de \esp{proteger}, avec \emph{g} $\rightarrow$ \emph{j})
\item « Parle à un adulte. » $\rightarrow$ \esp{Habla con un adulto.} (impératif affirmatif, sans \emph{-s})
\item « Il vaut mieux que tu vérifies l'information. » $\rightarrow$ \esp{Es mejor que verifiques la información.} (\esp{es mejor que} + subjonctif ; \emph{c} $\rightarrow$ \emph{qu} devant \emph{e})
\item « Ne télécharge pas cette application. » $\rightarrow$ \esp{No descargues esta aplicación.} (impératif négatif ; attention au faux ami : \esp{descargar} = « télécharger »)
\end{enumerate}
\end{corrige}

\begin{corrige}[Exercice 6 — Transformations grammaticales]
\textbf{A. Conseils à l'impératif :}
\begin{enumerate}
\item \esp{No compartas tu contraseña con tus amigos.} « Ne partage pas ton mot de passe avec tes amis. » (conseil négatif : danger)
\item \esp{Habla con tus padres si tienes un problema.} « Parle avec tes parents si tu as un problème. » (conseil positif)
\item \esp{No publiques fotos de otras personas sin permiso.} « Ne publie pas de photos d'autres personnes sans permission. » (\esp{publicar} : \emph{c} $\rightarrow$ \emph{qu} au subjonctif)
\item \esp{Verifica las noticias antes de compartirlas.} « Vérifie les nouvelles avant de les partager. » (impératif affirmatif)
\item \esp{No uses el móvil durante toda la noche.} « N'utilise pas le portable pendant toute la nuit. » (\esp{usar} : forme régulière \esp{no uses})
\end{enumerate}

\textbf{B. Phrases avec \esp{es importante que} + subjonctif :}
\begin{enumerate}
\item \esp{Es importante que protejas tus datos personales.} « Il est important que tu protèges tes données personnelles. »
\item \esp{Es importante que apagues el móvil antes de dormir.} « Il est important que tu éteignes le portable avant de dormir. » (\emph{g} $\rightarrow$ \emph{gu})
\item \esp{Es importante que leas bien los mensajes antes de responder.} « Il est important que tu lises bien les messages avant de répondre. »
\item \esp{Es importante que seas prudente con los desconocidos en internet.} « Il est important que tu sois prudent avec les inconnus sur internet. » (subjonctif irrégulier de \esp{ser})
\item \esp{Es importante que hagas deporte y no solo mires la pantalla.} « Il est important que tu fasses du sport et que tu ne regardes pas seulement l'écran. » (subjonctif irrégulier de \esp{hacer})
\end{enumerate}
\end{corrige}

\begin{corrige}[Exercice 7 — Texte à trous : La privacidad]
\textbf{Texte complété :}

\esp{Hoy casi todos los jóvenes usan las \textbf{redes sociales} para comunicarse con sus amigos. Pero hay que tener cuidado: es fundamental \textbf{proteger} nuestra \textbf{privacidad} en internet. Primero, no debemos publicar nuestros \textbf{datos personales}, como la dirección o el número de teléfono. Segundo, hay que crear una \textbf{contraseña} difícil y no dársela a nadie. Tercero, antes de \textbf{descargar} una aplicación nueva, debemos leer bien la información. Algunas \textbf{aplicaciones} piden demasiados permisos. También es importante \textbf{verificar} las noticias antes de difundirlas. Por último, si alguien sufre insultos en internet, es decir, \textbf{ciberacoso}, debe hablar con un adulto. No hay que \textbf{compartir} fotos íntimas con desconocidos. ¡Naveguemos con responsabilidad!}

\textbf{Traduction :} « Aujourd'hui, presque tous les jeunes utilisent les réseaux sociaux pour communiquer avec leurs amis. Mais il faut faire attention : il est fondamental de protéger notre vie privée sur internet. D'abord, nous ne devons pas publier nos données personnelles, comme l'adresse ou le numéro de téléphone. Ensuite, il faut créer un mot de passe difficile et ne le donner à personne. Troisièmement, avant de télécharger une nouvelle application, nous devons bien lire les informations. Certaines applications demandent trop de permissions. Il est important aussi de vérifier les nouvelles avant de les diffuser. Enfin, si quelqu'un subit des insultes sur internet, c'est-à-dire du cyberharcèlement, il doit parler avec un adulte. Il ne faut pas partager de photos intimes avec des inconnus. Naviguons avec responsabilité ! »

\textbf{Points de vigilance :} \esp{es fundamental} + infinitif (sujet général) ; \esp{antes de} + infinitif ; \esp{no hay que} + infinitif pour exprimer l'interdiction générale.
\end{corrige}

\begin{corrige}[Exercice 8 — Appariements : conseils et situations]
\begin{itemize}
\item \textbf{1--b.} Un inconnu demande des photos personnelles $\rightarrow$ \esp{No respondas y habla inmediatamente con un adulto.} « Ne réponds pas et parle immédiatement avec un adulte. »
\item \textbf{2--c.} Une nouvelle surprenante sur Facebook $\rightarrow$ \esp{Es importante que verifiques la información antes de compartirla.} « Il est important que tu vérifies l'information avant de la partager. »
\item \textbf{3--f.} Un camarade reçoit des insultes dans un groupe de chat $\rightarrow$ \esp{Es importante que le aconsejes hablar con un profesor o con sus padres.} « Il est important que tu lui conseilles de parler avec un professeur ou avec ses parents. »
\item \textbf{4--e.} Cinq heures de jeu d'affilée $\rightarrow$ \esp{Es importante que descanses y apagues la pantalla con frecuencia.} « Il est important que tu te reposes et que tu éteignes l'écran fréquemment. »
\item \textbf{5--d.} Un jeu d'une page inconnue $\rightarrow$ \esp{No descargues aplicaciones de páginas no oficiales.} « Ne télécharge pas d'applications de pages non officielles. »
\item \textbf{6--a.} Ton amie veut utiliser ton compte $\rightarrow$ \esp{No compartas tu contraseña, ni siquiera con tus amigos.} « Ne partage pas ton mot de passe, même pas avec tes amis. »
\end{itemize}
\end{corrige}

\begin{corrige}[Exercice 9 — Compréhension écrite : La adicción a las pantallas]
\textbf{Barème indicatif (20 points) :} questions de compréhension 1--5 : 5 x 2 pts = 10 pts ; question d'expression personnelle : 4 pts ; questions de langue 7--8 : 2 x 3 pts = 6 pts.

\textbf{Réponses modèles (en espagnol, phrases complètes) :}
\begin{enumerate}
\item \esp{Según el texto, muchos jóvenes miran su móvil más de cinco horas al día.} « Selon le texte, beaucoup de jeunes regardent leur portable plus de cinq heures par jour. »
\item \esp{Dos consecuencias negativas son que los jóvenes duermen peor y tienen dificultades para concentrarse en clase.} (Autres réponses possibles : \esp{abandonan el deporte y las actividades con los amigos} ; \esp{puede provocar ansiedad}.) « Deux conséquences négatives sont que les jeunes dorment moins bien et ont des difficultés à se concentrer en classe. »
\item \esp{Las redes sociales pueden provocar ansiedad porque los adolescentes se comparan constantemente con las imágenes perfectas que ven en internet.} « Les réseaux sociaux peuvent provoquer de l'anxiété parce que les adolescents se comparent constamment aux images parfaites qu'ils voient sur internet. »
\item \esp{No, las tecnologías no son siempre negativas: también son herramientas útiles para aprender, comunicarse con la familia y descubrir el mundo. El problema es el uso que hacemos de ellas.} « Non, les technologies ne sont pas toujours négatives : ce sont aussi des outils utiles pour apprendre, communiquer avec la famille et découvrir le monde. Le problème, c'est l'usage que nous en faisons. »
\item \esp{Los expertos aconsejan apagar el móvil una hora antes de dormir y limitar el tiempo de las pantallas durante la semana.} (Autres réponses possibles : \esp{hablar con los padres o profesores si se sienten mal} ; \esp{practicar deporte y mantener relaciones reales con los amigos}.) « Les experts conseillent d'éteindre le portable une heure avant de dormir et de limiter le temps des écrans pendant la semaine. »
\item \textbf{Expression personnelle — rédaction modèle :} \esp{Yo paso unas tres horas al día delante de las pantallas, sobre todo con el móvil. Uso mucho WhatsApp para hablar con mis amigos y también veo vídeos. Para limitar mi uso, apago el móvil cuando hago los deberes y no lo uso durante las comidas. Además, los domingos juego al fútbol con mis amigos del barrio. Creo que es importante que los jóvenes tengan actividades reales y no solo virtuales.} « Je passe environ trois heures par jour devant les écrans, surtout avec le portable. J'utilise beaucoup WhatsApp pour parler avec mes amis et je regarde aussi des vidéos. Pour limiter mon usage, j'éteins le portable quand je fais mes devoirs et je ne l'utilise pas pendant les repas. De plus, le dimanche, je joue au football avec mes amis du quartier. Je pense qu'il est important que les jeunes aient des activités réelles et pas seulement virtuelles. »
\item \textbf{Subjonctifs du texte :} \esp{«Es importante que los jóvenes apaguen el móvil una hora antes de dormir.»} et \esp{«Es necesario que limiten el tiempo de las pantallas durante la semana.»} (On accepte aussi : \esp{«Es fundamental que hablen con sus padres o profesores si se sienten mal.»}) Les subjonctifs sont \esp{apaguen}, \esp{limiten}, \esp{hablen}, \esp{sientan} : après une expression impersonnelle de nécessité ou d'importance, l'espagnol exige le subjonctif.
\item \textbf{Impératifs négatifs :} \esp{No mires el móvil toda la noche.} « Ne regarde pas le portable toute la nuit. » ; \esp{No te compares con los demás.} « Ne te compare pas aux autres. » (verbe pronominal : le pronom \esp{te} se place \emph{avant} le verbe à l'impératif négatif).
\end{enumerate}

\textbf{Points de vigilance :} Répondre par des phrases complètes en reprenant les mots de la question ; ne pas recopier tout le texte ; à la question 8, bien placer le pronom réfléchi devant le verbe (\esp{no te compares}, jamais \emph{no compares te}).
\end{corrige}

\begin{corrige}[Exercice 10 — Thème et version]
\textbf{A. Thème — traductions :}
\begin{enumerate}
\item « Ne partage pas tes photos personnelles sur internet. » $\rightarrow$ \esp{No compartas tus fotos personales en internet.}
\item « Il est important que tu protèges ta vie privée. » $\rightarrow$ \esp{Es importante que protejas tu privacidad.}
\item « Parle à un adulte si tu as un problème. » $\rightarrow$ \esp{Habla con un adulto si tienes un problema.}
\item « Il vaut mieux que tu vérifies l'information avant de la publier. » $\rightarrow$ \esp{Es mejor que verifiques la información antes de publicarla.} (infinitif + pronom \esp{la} soudé : \esp{publicarla})
\item « Ne télécharge pas cette application : elle est dangereuse. » $\rightarrow$ \esp{No descargues esta aplicación: es peligrosa.}
\end{enumerate}

\textbf{Barème indicatif du thème (10 pts) :} 2 pts par phrase (1 pt pour la forme verbale correcte, 1 pt pour le reste de la phrase).

\textbf{Points de vigilance :} Impératif négatif = \esp{no} + subjonctif ; \esp{protejas} avec \emph{j} ; \esp{verifiques} avec \emph{qu} ; « dangereuse » se dit \esp{peligrosa} (faux ami : \esp{dangerosa} n'existe pas).

\textbf{B. Version — traduction modèle :}

« Le cyberharcèlement est un problème sérieux parmi les jeunes. Si tu reçois des messages insultants, ne réponds pas avec agressivité et ne garde pas le silence : parle avec tes parents ou avec un professeur. Il est important que tu conserves les messages comme preuve et que tu bloques la personne qui t'attaque. Il est important aussi que les amis de la victime l'aident et ne partagent pas les insultes. Souviens-toi : derrière chaque écran, il y a une personne réelle, avec des sentiments réels. »

\textbf{Barème indicatif de la version (10 pts) :} compréhension globale du sens (4 pts) ; exactitude des tournures (4 pts) ; qualité du français (2 pts).

\textbf{Points de vigilance :} \esp{no guardes silencio} = « ne garde pas le silence » (et non « ne gardes pas le silence ») ; \esp{como prueba} = « comme preuve » ; \esp{bloquees} et \esp{conserves} sont des subjonctifs après \esp{es importante que}, rendus en français par le subjonctif ou l'indicatif selon la tournure ; \esp{Recuerda} est un impératif : « Souviens-toi ».
\end{corrige}

\begin{corrige}[Exercice 11 — Expression écrite : Artículo para el periódico del instituto]
\textbf{Barème indicatif (20 pts) :} respect de la consigne et du nombre de mots (4 pts) ; richesse du lexique (4 pts) ; correction grammaticale, impératif et subjonctif (6 pts) ; organisation et cohérence (4 pts) ; orthographe et accents (2 pts).

\textbf{Proposition de rédaction modèle (environ 95 mots) :}

\esp{\textbf{Las redes sociales: ¡úsalas con cabeza!}

\esp{Hoy casi todos los estudiantes usan las redes sociales en su móvil. Son útiles para comunicarse, pero también tienen riesgos.}

\esp{Los principales peligros son el ciberacoso, la adicción a las pantallas y la pérdida de la privacidad. Muchos jóvenes publican sus datos personales sin pensar en las consecuencias.}

\esp{Por eso, sigue estos consejos: no compartas tu contraseña con nadie, no descargues aplicaciones de páginas desconocidas y no publiques fotos íntimas. Además, es importante que verifiques las noticias antes de compartirlas y que hables con un adulto si sufres ciberacoso.}

\esp{Las redes sociales son una herramienta: ¡depende de ti usarlas bien!}}

\textbf{Traduction :} « Les réseaux sociaux : utilise-les avec intelligence ! Aujourd'hui, presque tous les élèves utilisent les réseaux sociaux sur leur portable. Ils sont utiles pour communiquer, mais ils ont aussi des risques. Les principaux dangers sont le cyberharcèlement, l'addiction aux écrans et la perte de la vie privée. Beaucoup de jeunes publient leurs données personnelles sans penser aux conséquences. C'est pourquoi, suis ces conseils : ne partage ton mot de passe avec personne, ne télécharge pas d'applications de pages inconnues et ne publie pas de photos intimes. De plus, il est important que tu vérifies les nouvelles avant de les partager et que tu parles avec un adulte si tu subis du cyberharcèlement. Les réseaux sociaux sont un outil : c'est à toi de bien les utiliser ! »

\textbf{Points de vigilance :}
\begin{itemize}
\item Vérifier la présence du titre, de l'introduction, des deux paragraphes (risques / conseils) et de la conclusion.
\item Employer au moins six mots du lexique : ici \esp{redes sociales, móvil, ciberacoso, adicción, privacidad, datos personales, contraseña, descargar, aplicaciones} (neuf mots).
\item Impératifs négatifs corrects : \esp{no compartas}, \esp{no descargues}, \esp{no publiques} ; subjonctifs après \esp{es importante que} : \esp{verifiques}, \esp{hables}.
\item Soigner les accents : \esp{aplicación, adicción, información, íntimas, úsalas}.
\end{itemize}
\end{corrige}

\newpage

\coursec{Fiche bilan}

\begin{popfigure}
\begin{center}\tcbox[colback=popblue,colframe=popblue,coltext=white,arc=9pt,boxsep=4pt,left=6pt,right=6pt]{\bfseries\small Las tecnologías digitales}\end{center}
\vspace{-2pt}
\begin{tcbraster}[raster columns=3,raster equal height=rows,raster column skip=6pt,raster row skip=6pt,enhanced,blankest]
\begin{tcolorbox}[enhanced,colback=poporange,colframe=poporange,coltext=white,boxrule=0.9pt,arc=3pt,left=4pt,right=4pt,top=3pt,bottom=3pt,boxsep=1pt,halign=left]\bfseries\scriptsize Lexique : móvil, aplicación, contraseña, descargar\end{tcolorbox}
\begin{tcolorbox}[enhanced,colback=poppurple,colframe=poppurple,coltext=white,boxrule=0.9pt,arc=3pt,left=4pt,right=4pt,top=3pt,bottom=3pt,boxsep=1pt,halign=left]\bfseries\scriptsize Riesgos : ciberacoso, adicción, privacidad\end{tcolorbox}
\begin{tcolorbox}[enhanced,colback=popgreen,colframe=popgreen,coltext=white,boxrule=0.9pt,arc=3pt,left=4pt,right=4pt,top=3pt,bottom=3pt,boxsep=1pt,halign=left]\bfseries\scriptsize Impératif : no compartas, protégete\end{tcolorbox}
\begin{tcolorbox}[enhanced,colback=poppink,colframe=poppink,coltext=white,boxrule=0.9pt,arc=3pt,left=4pt,right=4pt,top=3pt,bottom=3pt,boxsep=1pt,halign=left]\bfseries\scriptsize Subjonctif : es importante que...\end{tcolorbox}
\begin{tcolorbox}[enhanced,colback=popgold,colframe=popgold,coltext=white,boxrule=0.9pt,arc=3pt,left=4pt,right=4pt,top=3pt,bottom=3pt,boxsep=1pt,halign=left]\bfseries\scriptsize Conseils : comportements responsables\end{tcolorbox}
\begin{tcolorbox}[enhanced,colback=popblue!70,colframe=popblue!70,coltext=white,boxrule=0.9pt,arc=3pt,left=4pt,right=4pt,top=3pt,bottom=3pt,boxsep=1pt,halign=left]\bfseries\scriptsize Méthode : commentaire de texte\end{tcolorbox}
\begin{tcolorbox}[enhanced,colback=popdark,colframe=popdark,coltext=white,boxrule=0.9pt,arc=3pt,left=4pt,right=4pt,top=3pt,bottom=3pt,boxsep=1pt,halign=left]\bfseries\scriptsize Expression : proposer, évaluer\end{tcolorbox}
\end{tcbraster}

\legende{Carte mentale de la leçon : le numérique, ses risques et les outils pour conseiller.}
\end{popfigure}

\begin{bilan}
\textbf{1. Lexique clé (à connaître par cœur)}
\begin{itemize}
  \item \esp{internet} : internet ; \esp{el móvil} : le portable ; \esp{la aplicación} : l'application
  \item \esp{la contraseña} : le mot de passe ; \esp{descargar} : télécharger
  \item \esp{el ciberacoso} : le cyberharcèlement ; \esp{la adicción} : l'addiction
  \item \esp{la privacidad} : la vie privée ; \esp{compartir} : partager ; \esp{proteger} : protéger
\end{itemize}

\textbf{2. Grammaire à connaître par cœur}
\begin{center}
\begin{tabularx}{\linewidth}{|X|X|X|}
\hline
\rowcolor{popblueL} Règle & Formation & Exemple \\
\hline
Impératif affirmatif (tú) & 3\textsuperscript{e} pers. sg. du présent & \esp{Protege tus datos.} \\
\hline
Impératif négatif (tú) & \esp{no} + subjonctif présent & \esp{No compartas tu contraseña.} \\
\hline
Subjonctif de conseil & \esp{es importante que} + subjonctif & \esp{Es importante que uses contraseñas seguras.} \\
\hline
Subjonctif présent & radical de la 1\textsuperscript{re} pers. + terminaisons inversées (-a/-e) & \esp{hable, coma, escriba} \\
\hline
\end{tabularx}
\end{center}

\textbf{3. Méthodes express}
\begin{itemize}
  \item \cle{Commentaire de texte} : présenter (auteur, thème, but) $\rightarrow$ résumer les idées $\rightarrow$ analyser les risques cités $\rightarrow$ donner son avis argumenté.
  \item \cle{Conseiller} : impératif pour l'ordre direct (\esp{No descargues...}), subjonctif pour le conseil nuancé (\esp{Te recomiendo que...}).
  \item \cle{Évaluer un risque} : nommer le danger, expliquer ses conséquences, proposer une solution concrète.
\end{itemize}
\end{bilan}

\begin{aretenir}[Checklist avant le Bac]
\begin{itemize}
  \item $\square$ Je connais le lexique du numérique : \esp{móvil, aplicación, contraseña, ciberacoso, privacidad, descargar}.
  \item $\square$ Je sais former l'impératif affirmatif de \esp{tú} : \esp{comparte, protege, descarga}.
  \item $\square$ Je sais former l'impératif négatif : \esp{no} + subjonctif (\esp{no compartas, no abras}).
  \item $\square$ Je connais les impératifs irréguliers : \esp{di, haz, ve, pon, sal, ten, ven, sé}.
  \item $\square$ Je sais conjuguer le subjonctif présent des verbes réguliers : \esp{hable, coma, escriba}.
  \item $\square$ J'emploie le subjonctif après \esp{es importante que}, \esp{es necesario que}, \esp{te recomiendo que}.
  \item $\square$ Je sais énumérer trois risques du numérique : cyberharcèlement, addiction, atteinte à la vie privée.
  \item $\square$ Je sais proposer des comportements responsables en deux ou trois phrases.
  \item $\square$ Je n'oublie pas les accents : \esp{aplicación, adicción, contraseña}.
  \item $\square$ Je n'écris pas \esp{privacité} : le mot correct est \esp{privacidad}.
  \item $\square$ Je sais structurer un commentaire : présentation, résumé, analyse, avis personnel.
  \item $\square$ Je relie mes idées avec \esp{además, sin embargo, por eso, en conclusión}.
\end{itemize}
\end{aretenir}

\findecours

\end{document}
